\documentclass[10pt,a4paper]{article}

\usepackage{amsmath,amssymb,amsthm,mathtools}
\usepackage[hmargin=2.5cm,vmargin=2.4cm]{geometry}
\usepackage{booktabs}
\usepackage{longtable}
\usepackage{microtype}
\usepackage{enumitem}
\usepackage{array}
\usepackage[hidelinks]{hyperref}

\theoremstyle{plain}
\newtheorem{theorem}{Theorem}[section]
\newtheorem{proposition}[theorem]{Proposition}
\newtheorem{lemma}[theorem]{Lemma}
\newtheorem{corollary}[theorem]{Corollary}
\theoremstyle{definition}
\newtheorem{definition}[theorem]{Definition}
\newtheorem{remark}[theorem]{Remark}
\newtheorem{convention}[theorem]{Convention}

\newcommand{\R}{\mathbb{R}}
\newcommand{\N}{\mathbb{N}}
\newcommand{\Q}{\mathbb{Q}}
\newcommand{\C}{\mathbb{C}}
\newcommand{\E}{\mathbb{E}}
\newcommand{\PR}{\mathbb{P}}
\newcommand{\Cii}{C_{\mathrm{iid}}}
\newcommand{\bigO}{O}
\newcommand{\lean}[1]{\texttt{#1}}
\newcommand{\file}[1]{\texttt{#1}}
\DeclareMathOperator{\sinc}{sinc}
\DeclareMathOperator{\Rea}{Re}
\DeclareMathOperator{\Ima}{Im}
\DeclareMathOperator*{\esssup}{ess\,sup}

\newenvironment{leanref}%
  {\par\addvspace{3pt}\begingroup\sloppy\emergencystretch=4em
   \noindent\footnotesize\textsc{Lean.}\ \ignorespaces}%
  {\par\endgroup\addvspace{3pt}}

\tolerance=2000
\setlength{\emergencystretch}{2.5em}

\setlength{\abovedisplayskip}{5pt plus 2pt minus 2pt}
\setlength{\belowdisplayskip}{5pt plus 2pt minus 2pt}
\setlength{\abovedisplayshortskip}{3pt plus 2pt}
\setlength{\belowdisplayshortskip}{4pt plus 2pt minus 2pt}

\title{\bf Two Certified Bounds on the Universal\\ i.i.d.\ Berry--Esseen Constant:\\[2mm]
\large $0.40 \le \Cii \le 0.4688$}

\author{A formalization report on the \file{BerryEsseen} Lean development}
\date{13 September 2026}

\begin{document}
\maketitle

\begin{abstract}
\noindent
Let $X_1,X_2,\dots$ be independent, identically distributed real random
variables with $\E X_1=0$, $\E X_1^2=1$ and $\beta:=\E|X_1|^3<\infty$, let
$F_n$ be the distribution function of $(X_1+\dots+X_n)/\sqrt n$, and let
$\Phi$ be the standard normal distribution function. The universal i.i.d.\
Berry--Esseen constant $\Cii$ is the least constant $C$ for which
$\sup_x|F_n(x)-\Phi(x)|\le C\beta/\sqrt n$ holds for every such law and every
$n\ge1$. This document presents, in ordinary mathematical language, the
complete proofs of the two bounds
\[
  0.40 \;\le\; \Cii \;\le\; 0.4688 \;=\; \tfrac{293}{625},
\]
as they are formalized in the Lean~4 development \file{BerryEsseen} on top of
mathlib. The two statements are exactly the two declarations that the
project's pinned, sandboxed comparator is configured to check.

The lower bound comes from an explicit six-summand standardized Bernoulli
witness with success probability $2/5$, evaluated by exact rational
arithmetic, together with one rational estimate of a normal probability. For
completeness we also give the (not yet formalized) asymptotic proof of
Esseen's stronger lower bound $(\sqrt{10}+3)/(6\sqrt{2\pi})$.

The upper bound is a Fourier-analytic argument that removes the unknown law
analytically and then verifies interval inequalities. Its ingredients are:
cosine-anchor estimates that retain the information that $\beta=1$ forces the
symmetric two-point law; a global exponential modulus bound obtained from
symmetrization and an explicitly constructed convex minorant; Prawitz's
smoothing inequality, proved from scratch through a shift recurrence and the
Riemann--Lebesgue lemma; a power-comparison step that keeps the Gaussian
damping inside a geometric average; and four exhaustive parameter regions ---
an analytic small-fraction region with the stronger regional constant
$183/400$, an elementary Cantelli region, $467$ certified interval cells
covering all sample sizes $n\ge20$ or third moments $\beta\ge2$, and $174$
certified finite-sample cells covering $1\le n\le19$, $\beta<2$.

All finite certificates are checked in Lean's kernel by \lean{decide +kernel};
the development contains no \lean{sorry} and no \lean{native\_decide}, and its
axiom closure is $\{\lean{propext},\lean{Classical.choice},\lean{Quot.sound}\}$.
The final pinned sandboxed comparator run is \emph{pending} at the time of
writing; nothing in this document should be read as reporting a completed
comparator verification.
\end{abstract}

\tableofcontents

\newpage

%=====================================================================
\section{Introduction}
%=====================================================================

\subsection{The constant}

Let $X$ be a real random variable with
\begin{equation}\label{eq:admissible}
  \E X = 0, \qquad \E X^2 = 1, \qquad \beta := \E|X|^3 < \infty ,
\end{equation}
let $X_1,X_2,\dots$ be independent copies of $X$, and put
\[
  S_n := \frac{X_1+\dots+X_n}{\sqrt n},
  \qquad F_n(x) := \PR(S_n \le x),
  \qquad \Delta_n := \sup_{x\in\R}\bigl|F_n(x)-\Phi(x)\bigr| .
\]
The Berry--Esseen theorem \cite{Berry1941,Esseen1942} asserts the existence of
an absolute constant $C$ with $\Delta_n \le C\,\beta/\sqrt n$ for all $n\ge1$
and all laws satisfying \eqref{eq:admissible}. We write $\Cii$ for the least
such constant and
\[
  \ell := \frac{\beta}{\sqrt n}
\]
for the \emph{Lyapunov fraction}, so that the assertion to be proved is
$\Delta_n \le \Cii\,\ell$ and $\Cii = \sup \Delta_n/\ell$, the supremum being
over all admissible laws and all $n\ge1$.

Table~\ref{tab:history} records the published history of two-sided bounds for
$\Cii$, following the survey in \cite{Tao2026}. Zolotarev's conjecture
\cite{Zolotarev1967} is that $\Cii$ equals Esseen's lower bound
\[
  c_E := \frac{\sqrt{10}+3}{6\sqrt{2\pi}} = 0.4097321837\ldots,
\]
and this is known in the i.i.d.\ Bernoulli (binomial) subcase
\cite{Schulz2016,Zolotukhin2018}.

\begin{table}[ht]
\centering
\begin{tabular}{@{}llp{7.4cm}@{}}
\toprule
Bound & Reference & Comment \\
\midrule
$\Cii \le 0.82$   & \cite{Zolotarev1967} & Zolotarev-type smoothing inequalities;
  \cite{Zolotarev1967} also gives $0.9051$ for the non-i.i.d.\ case. \\
$\Cii \le 0.7975$ & \cite{vanBeek1972}   & Fourier-analytic refinement. \\
$\Cii \le 0.7655$ & \cite{Shiganov1986}  & \\
$\Cii \le 0.7056$ & \cite{Shevtsova2006} & \\
$\Cii \le 0.5129$ & \cite{KorolevShevtsova2009}
  & From the structural bound $\Delta_n\le 0.34445(\beta+0.489)/\sqrt n$. \\
$\Cii \le 0.4785$ & \cite{Tyurin2009}    & \\
$\Cii \le 0.4748$ & \cite{Shevtsova2011} & \\
$\Cii \le 0.4690$ & \cite{Shevtsova2013}
  & From $\Delta_n\le n^{-1/2}\min\{0.4690\beta,\,0.3322(\beta+0.429),\,
    0.3031(\beta+0.646)\}$. \\
\midrule
$\Cii \ge c_E$    & \cite{Esseen1956}
  & $c_E=(\sqrt{10}+3)/(6\sqrt{2\pi})\approx0.4097321837$, attained
    asymptotically by normalized Bernoulli sums. \\
\bottomrule
\end{tabular}
\caption{Published bounds for the universal i.i.d.\ Berry--Esseen constant,
as collected in \cite{Tao2026}.}
\label{tab:history}
\end{table}

\subsection{The two theorems of this document}

The Lean development \file{BerryEsseen} proves two separate statements about
the same defined object. In Lean they read as follows (the extended nonnegative
reals are written \lean{ENNReal} here; the source file uses the Unicode
notation for that type):
\begin{verbatim}
theorem berry_esseen_constant_lower_bound :
    (0.40 : ENNReal) <= BerryEsseen.berryEsseenConstant

theorem berry_esseen_constant_upper_bound :
    BerryEsseen.berryEsseenConstant <= 0.4688
\end{verbatim}
and the object \lean{BerryEsseen.berryEsseenConstant} is defined as an
extended nonnegative real supremum, so that finiteness of $\Cii$ is a
consequence of the second theorem rather than a hypothesis. The precise
definitions are reproduced in Section~\ref{sec:statement}.

\begin{theorem}[Lower bound; \lean{BerryEsseen.berryEsseenConstant\_lower\_bound}]
\label{thm:main-lower}
$\;\dfrac{2}{5} \le \Cii$.
\end{theorem}

\begin{theorem}[Upper bound; \lean{BerryEsseen.berryEsseenConstant\_upper\_bound}]
\label{thm:main-upper}
$\;\Cii \le \dfrac{293}{625} = 0.4688$.
\end{theorem}

These two declarations are exactly the \emph{comparator targets}: the file
\file{comparator/config.json} lists \lean{berry\_esseen\_constant\_lower\_bound}
and \lean{berry\_esseen\_constant\_upper\_bound} as the two theorem names to be
compared against the trusted challenge file, with permitted axiom list
$\{\lean{propext},\lean{Quot.sound},\lean{Classical.choice}\}$.

\subsection{Shape of the upper-bound proof}

The obstruction to a short proof of a numerically good constant is that the
unknown law enters through its characteristic function $f(t)=\E e^{itX}$, and
that the Kolmogorov distance is sensitive to the whole frequency range,
including lattice laws whose characteristic functions return to modulus one.
The proof organizes this as follows.

\begin{enumerate}[leftmargin=2.2em]
\item \textbf{Anchors} (Section~\ref{sec:anchors}). The third absolute moment
  satisfies $\beta\ge1$ with equality only for the symmetric two-point law,
  whose characteristic function is $\cos t$. Quantitatively
  $\E[(|X|-1)^2(2|X|+1)]=2(\beta-1)$, and this identity yields
  \[
    |\Rea f(t)-\cos t| \le \tfrac{\beta-1}{6}|t|^3,
    \qquad
    |\Ima f(t)| \le \min\Bigl\{
      \tfrac{|t|^3}{6}\sqrt{(\beta-1)(\beta+\tfrac53)},\;
      \tfrac{t^2}{2}\sqrt{\beta-1},\;
      \tfrac{\beta|t|^3}{6}\Bigr\}.
  \]
  Both bounds degenerate as $\beta\downarrow1$, which the classical
  third-order Taylor estimate does not.
\item \textbf{Global modulus} (Section~\ref{sec:modulus}). Symmetrization and
  an explicitly constructed convex, nonincreasing minorant $h$ of
  $s\mapsto(1-\cos s)/s^2$ give
  $|f(t)|\le\exp\{-t^2h((\beta+1)|t|)\}$ at \emph{every} frequency.
\item \textbf{Prawitz smoothing} (Section~\ref{sec:prawitz}). A pointwise
  majorant of the indicator of $[0,\infty)$, built from a truncated cotangent
  Fourier kernel, converts the two characteristic-function estimates into four
  positive real integrals over a bounded frequency band plus one explicit
  Gaussian tail.
\item \textbf{Power comparison} (Section~\ref{sec:compression}). Because
  $f_n(t)=f(t/\sqrt n)^n$, a difference of $n$-th powers must be estimated.
  Keeping the differing decay rates of the two bases inside a geometric
  average, rather than bounding the geometric sum by $n$ times its largest
  term, is what makes a single inequality valid for infinitely many sample
  sizes.
\item \textbf{Four exhaustive regions} (Sections~\ref{sec:small}--\ref{sec:finite}).
  See Table~\ref{tab:regions}.
\end{enumerate}

\begin{table}[ht]
\centering
\begin{tabular}{@{}lll r@{}}
\toprule
Region & Method & Section & Regional bound \\
\midrule
$\ell\le\frac1{20}$
  & analytic, no computation & \ref{sec:small} & $\frac{183}{400}=0.4575$ \\
$\ell\ge\frac65$
  & Cantelli, no computation & \ref{sec:large} & $\frac{11}{24}=0.4583\overline{3}$ \\
$\ell\ge\frac1{20}$ and ($\beta\ge2$ or $n\ge20$)
  & $467$ certified interval cells & \ref{sec:scalar} & $\frac{293}{625}=0.4688$ \\
$1\le n\le19$, $1\le\beta<2$, $\ell<\frac65$
  & $174$ certified finite cells & \ref{sec:finite} & $\frac{293}{625}=0.4688$ \\
\bottomrule
\end{tabular}
\caption{The four parameter regions of the upper-bound proof. The union of the
regions is the whole admissible parameter set; see
Section~\ref{sec:assembly}.}
\label{tab:regions}
\end{table}

\subsection{What is proved, and what is verified}
\label{sec:intro-status}

It is important to keep three different claims apart.

\begin{enumerate}[leftmargin=2.2em]
\item \emph{Mathematics.} Theorems~\ref{thm:main-lower} and
  \ref{thm:main-upper} are formalized in Lean~4 with mathlib. There is no
  \lean{sorry} in the development, and no published Berry--Esseen bound is
  imported as an assumption. The axiom closure of the two theorems is
  $\{\lean{propext},\lean{Classical.choice},\lean{Quot.sound}\}$ --- the
  standard mathlib axioms --- and in particular \lean{native\_decide} is
  \emph{not} used anywhere in this development.
\item \emph{Finite checks.} All finite certificates are discharged by
  \lean{decide +kernel}, i.e.\ they are evaluated by Lean's kernel on exact
  integer and rational arithmetic. A direct count over the certificate
  directories gives $961{,}710$ occurrences of \lean{decide +kernel}
  (Table~\ref{tab:kernel}).
\item \emph{Comparator.} The project's intended final gate is a fresh,
  sandboxed Linux run of the pinned \file{leanprover/comparator}
  \cite{Comparator2026}, which re-elaborates the statement against the
  self-contained challenge file, checks the axiom whitelist, and replays the
  dependency closure in Lean's kernel. \textbf{That run is pending.} Neither
  the repository nor this document claims comparator completion. A passing
  ordinary \lean{lake build} is a different and weaker event, and is reported
  separately; see Section~\ref{sec:verification}.
\end{enumerate}

\subsection{A second, sharper development in the same repository}
\label{sec:0423-note}

The same repository contains a separate development that proves the sharper
bound $\Cii\le0.423$ (\lean{BerryEsseen.Upper0423.berryEsseenConstant\_le\_0423}
in \file{BerryEsseen/Upper0423/Main.lean}, documented in
\file{docs/UPPER\_0423.md}). That development uses \lean{native\_decide} for its
finite certificates, so its axiom report contains, besides the three standard
axioms, one auxiliary axiom per finite certificate ($579$ in total); it
therefore trusts Lean's compiler and evaluator in addition to the kernel, a
strictly weaker trust boundary than the one described in
Section~\ref{sec:intro-status}. It does not modify, and is not imported by,
the $0.4688$ development: the default library target does not import
\file{Upper0423}, so the comparator configuration is unaffected. The two
theorems presented in this document, and not the $0.423$ theorem, are the
comparator targets.

\subsection{Provenance and attribution}
\label{sec:provenance}

Because the same repository hosts two developments with different provenance,
we state explicitly what the repository's own documentation records.

The $0.4688$ development discussed in this document contains \emph{no}
statement attributing any of its ingredients to an external source beyond the
classical literature. A search of \file{docs/} and of the Lean sources outside
\file{BerryEsseen/Upper0423/} for attribution language (``due to'', ``adapted
from'', ``based on'', ``attribution'', author names) returns only the names of
classical results that the proofs invoke and prove for themselves: Prawitz's
smoothing inequality and his kernel \cite{Prawitz1972}; the one-sided
Chebyshev (Cantelli) inequality; Jordan's sine inequality; the
Riemann--Lebesgue lemma; Taylor's theorem; Cauchy--Schwarz; and Esseen's
lower-bound construction \cite{Esseen1956}, whose Bernoulli coefficient the
notes take from the historical record in \cite{Zolotukhin2018}. The broader
strategy --- Fourier smoothing with a finite frequency band, a
moment-dependent family of characteristic-function estimates, and a split of
the parameter space into regions treated by different arguments --- is the
strategy of the literature summarized in Table~\ref{tab:history}, in
particular of \cite{Zolotarev1967,vanBeek1972,Shiganov1986,Shevtsova2006,%
KorolevShevtsova2009,Tyurin2009,Shevtsova2011,Shevtsova2013}.

The one attribution section in the repository is in
\file{docs/UPPER\_0423.md}, and it concerns the \emph{separate} $0.423$
development of Section~\ref{sec:0423-note}. It reads as follows (its bullet list
is run together here, and its per-file cross-references are omitted; the
ellipsis marks the omission of a numerical bracketing of $\kappa$ and of the
coarser rational majorant used there):
\begin{quote}\small\sloppy
``The analytic framework of this proof is not new here. The research notes
\file{research/imaginary-envelope.md} (``Source alignment and attribution'')
and \file{research/recovery/RECOVERY.md} attribute the following ingredients to
[haonan-xiao/iid-berry-esseen, commit \texttt{f00d781c}], which proves a
$0.4395$ bound for the same constant: the shared moment parameters, the radial
rectangle and the disk/strip discussion, the shared convex profile, the moment
region, and the Prawitz normalization convention; the two-point sine-circle
mechanism; the stop-loss identity for $\E|X-X'|^3$ and the resulting bounds on
the symmetrization excess $\eta$; the separate radial estimates for the real
and imaginary parts; the cosine constant
$\kappa=\max_{v>0}(\cos v-1+v^2/2)/v^3$ \dots{} and the same convex profile
$1/2-\kappa v$ continued by $(1-\cos v)/v^2$.''
\end{quote}
The cited research notes are not part of the repository; on the machine used to
prepare this document they sit in a sibling \file{research/} directory next to
the checkout. They record, conversely, that two of the ingredients used by the
$0.423$ argument are taken \emph{from} the present $0.4688$ development: they
list ``our restored repository's \file{BerryEsseen/WeightedMoments.lean},
\file{ImaginaryPart.lean}, and \file{TrigonometricAnchors.lean}'' as the source
of ``the earlier beta-minus-one anchor estimates'', and state that ``the
restored project's \file{docs/ANCHORS.md} and \file{docs/PROOF\_OVERVIEW.md}
retain the original anchor and geometric-sum arguments''.

Two of the ingredients that \file{docs/UPPER\_0423.md} attributes to
\cite{XiaoLi2026} do have counterparts in the $0.4688$ development described
here, and we record this as a factual overlap without any claim about
priority:
\begin{enumerate}[leftmargin=2.2em]
\item the \emph{convex cosine profile}: the minorant $h$ of
  Section~\ref{sec:modulus} has exactly the shape ``a linear branch
  $\tfrac12-cv$ continued by a scaled $(1-\cos v)/v^2$''. The present
  development uses the coarser rational slope $c=\tfrac1{10}$, justified by
  its own lemma $\cos v-1+v^2/2\le|v|^3/10$
  (\lean{BerryEsseen.cosine\_tenth} in \file{BerryEsseen/CosineMinorant.lean}),
  and splices at $v=9/2$;
\item the \emph{Prawitz normalization convention}: the present development
  uses the doubled positive-frequency kernel
  $K_\sigma(s)=\sigma(1-s)+iB(s)$ with no compensating factor $\tfrac12$
  outside the integral, and says so explicitly in
  \file{docs/PRAWITZ\_MAJORANT.md} and \file{docs/SMALL\_FRACTION.md}
  (Convention~\ref{conv:kernel} below).
\end{enumerate}
Beyond this, the anchors of Section~\ref{sec:anchors} and the geometric-sum
arguments of Section~\ref{sec:compression} are the arguments that the research
notes ascribe to this development itself. This document makes no assessment of
anyone's priority or of the correctness of \cite{XiaoLi2026}.

\subsection{Conventions and reading guide}

Throughout, $\Phi$ and $\varphi$ are the standard normal distribution function
and density, $f(t)=\E e^{itX}$, $g(t)=e^{-t^2/2}$,
$\sinc(u)=\sin(u)/u$ for $u\ne0$ and $\sinc(0)=1$, and
\[
  q(x) := \frac1x - \cot x \qquad (0<x<\pi).
\]
All named constants are exact rationals or explicit algebraic numbers; no
floating-point number occurs in any statement.

\begin{convention}[Kernel normalization]\label{conv:kernel}
For $\sigma\in\{+1,-1\}$ and $0<s<1$ set
\[
  B(s) := (1-s)\cot(\pi s) + \frac1\pi,
  \qquad
  K_\sigma(s) := \sigma(1-s) + i\,B(s),
  \qquad K := K_{+1}.
\]
This $K$ is \emph{twice} the kernel that is usually called the Prawitz kernel;
the compensating factor $2$ that normally stands outside the integral is
absent. A reader comparing formulas with a source that uses the
half-normalized kernel must insert factors of two in the first three of the
four smoothing terms (the tail term is unchanged). The Lean definition is
\lean{BerryEsseen.prawitzKernel} in \file{BerryEsseen/PrawitzFourier.lean}.
\end{convention}

Each mathematical statement below is followed by a \textsc{Lean} note giving
the name of the corresponding declaration and its file, so that a reader can
locate the formal counterpart of every step. Where a proof is long or purely
computational we give a faithful sketch and the pointer; where a proof is short
we give it.

%=====================================================================
\section{The formal statement}
\label{sec:statement}
%=====================================================================

The definitions below are reproduced from \file{BerryEsseen/Defs.lean}; the
comparator's trusted file \file{Challenge.lean} repeats them verbatim and
imports only mathlib \cite{Mathlib}, so that nothing about the meaning of the
statement depends on the project's own modules. In particular the normal law is
mathlib's \lean{ProbabilityTheory.gaussianReal 0 1}, the product measure is
mathlib's \lean{MeasureTheory.Measure.pi}, and integrability, measurability and
the supremum are mathlib's.

\subsection{Admissible laws}

\begin{definition}\label{def:admissible}
For a Borel probability measure $\mu$ on $\R$ put
\[
  \beta(\mu) := \int_\R |x|^3\,d\mu(x) \in[0,\infty] .
\]
The measure $\mu$ is \emph{admissible} if
\begin{enumerate}[label=(\alph*),leftmargin=2.4em,itemsep=0pt]
\item $x\mapsto x$, $x\mapsto x^2$ and $x\mapsto|x|^3$ are $\mu$-integrable;
\item $\int x\,d\mu = 0$;
\item $\int x^2\,d\mu = 1$.
\end{enumerate}
\end{definition}

In Lean, $\beta(\mu)$ is \lean{BerryEsseen.thirdAbsMoment} and
Definition~\ref{def:admissible} is the structure
\lean{BerryEsseen.Admissible}, whose five fields are exactly (a)--(c). Note
that integrability is part of the definition rather than an afterthought: the
class of laws is not restricted to bounded or discrete laws, and no moment
beyond the third is assumed.

\subsection{Normalized sums via a product measure}

Independence is realized by a finite product measure, so that the object whose
distribution function is estimated is literally the law of a sum of $n$
independent copies.

\begin{definition}\label{def:sumlaw}
For $n\in\N$ define $\Sigma_n:\R^{\mathrm{Fin}\,n}\to\R$ by
$\Sigma_n(x) = \bigl(\sum_{i}x_i\bigr)/\sqrt n$, and let
\[
  \mu_n := (\Sigma_n)_{*}\Bigl(\textstyle\bigotimes_{i=1}^{n}\mu\Bigr)
\]
be the pushforward of the $n$-fold product measure. Let
$\Phi(x) := \gamma(({-\infty},x])$ where $\gamma$ is mathlib's
\lean{ProbabilityTheory.gaussianReal 0 1}.
\end{definition}

\begin{leanref}
\lean{normalizedSum}, \lean{normalizedSumLaw}, \lean{normalCDF} in
\file{BerryEsseen/Defs.lean}.
\end{leanref}

\subsection{The constant as a supremum}

\begin{definition}\label{def:constant}
For admissible $\mu$, $n\ge1$ and $x\in\R$ put
\[
  N(\mu,n,x) := \Bigl(
    \frac{\sqrt n}{\beta(\mu)}\,\bigl|\mu_n(({-\infty},x]) - \Phi(x)\bigr|
  \Bigr)_{+} \in[0,\infty] ,
\]
the value being taken in the extended nonnegative reals, and define
\[
  \Cii := \sup_{\mu\ \mathrm{admissible}} \ \sup_{n\ge1}\ \sup_{x\in\R}
  N(\mu,n,x) \ \in[0,\infty] .
\]
\end{definition}

\begin{leanref}
\lean{normalizedError} and \lean{berryEsseenConstant} in
\file{BerryEsseen/Defs.lean}. In Lean, $\Cii$ is an
\lean{ENNReal}-valued iterated supremum
$\bigsqcup_\mu\bigsqcup_{h:\,\mathrm{Admissible}\,\mu}\bigsqcup_n
\bigsqcup_{h:\,0<n}\bigsqcup_x \lean{normalizedError}\ \mu\ n\ x$, and
$N(\mu,n,x)=\lean{ENNReal.ofReal}(\cdots)$.
\end{leanref}

Three features of Definition~\ref{def:constant} deserve emphasis, since they
are what make the upper bound a substantive statement.

\begin{enumerate}[leftmargin=2.2em]
\item The value is \emph{a priori} allowed to be $+\infty$. Finiteness of
  $\Cii$ is therefore a theorem (a consequence of
  Theorem~\ref{thm:main-upper}) and not a side condition imposed on the
  definition.
\item The supremum is over \emph{all} admissible laws. There is no restriction
  to a finite family, to laws with bounded support, or to numerically sampled
  distributions.
\item The inner supremum over $x$ is over all real thresholds, and
  $\mu_n(({-\infty},x])$ is a genuine right-continuous distribution function;
  atoms are permitted and are handled explicitly in the smoothing step
  (Section~\ref{sec:prawitz}).
\end{enumerate}

\subsection{The reduction used by both bounds}

The following triviality is the interface between
Definition~\ref{def:constant} and everything that follows.

\begin{lemma}\label{lem:reduction}
For every $C\in[0,\infty]$,
\[
  \Cii \le C
  \iff
  \bigl(\forall\,\mu\ \mathrm{admissible},\ \forall n\ge1,\ \forall x\in\R:\
    N(\mu,n,x)\le C\bigr).
\]
Moreover $N(\mu,n,x)\le\Cii$ for every admissible $\mu$, $n\ge1$, $x$.
\end{lemma}

\begin{proof}
Both directions are the universal property of an iterated supremum in the
complete lattice $[0,\infty]$.
\end{proof}

\begin{leanref}
\lean{BerryEsseen.constant\_le\_iff} and
\lean{BerryEsseen.normalizedError\_le\_constant} in
\file{BerryEsseen/Constant.lean}.
\end{leanref}

Thus the lower bound is proved by exhibiting one triple $(\mu,n,x)$, and the
upper bound is proved by a pointwise estimate valid for all triples. The
translation from a bound on the Kolmogorov distance to a bound on
$N(\mu,n,x)$ is the following elementary step, used by every region of the
upper-bound proof.

\begin{lemma}\label{lem:cdf-to-normalized}
Let $\mu$ be admissible, $n\ge1$, $C\ge0$, $a\in\R$ with
$a\le\beta(\mu)/\sqrt n$, and suppose
$\bigl|\mu_n(({-\infty},x])-\Phi(x)\bigr|\le C\,a$. Then
$N(\mu,n,x)\le C$.
\end{lemma}

\begin{proof}
Since $\beta(\mu)>0$ for admissible $\mu$ (Lemma~\ref{lem:beta-ge-one}),
multiply the hypothesis by $\sqrt n/\beta(\mu)>0$ and use
$a\le\beta(\mu)/\sqrt n$ to get
$\frac{\sqrt n}{\beta}|{\cdots}|\le \frac{\sqrt n}{\beta}\,C\,\frac{\beta}{\sqrt n}=C$.
\end{proof}

\begin{leanref}
\lean{BerryEsseen.normalizedError\_le\_of\_cdf\_bound} in
\file{BerryEsseen/ParameterCellBounds.lean}; the variant used by the
finite-sample cells, which replaces $\sqrt n/\beta$ by $\sqrt H/L$ for a cell
with $n\le H$ and $L\le\beta$, is
\lean{BerryEsseen.normalizedError\_le\_finite\_budget} in
\file{BerryEsseen/FiniteSmoothing.lean}.
\end{leanref}

The asymmetry between $a$ and $\beta/\sqrt n$ in
Lemma~\ref{lem:cdf-to-normalized} is the mechanism by which finitely many
cells cover a continuum of parameters: on a cell whose Lyapunov fraction
ranges over $[a,b]$, the analytic majorants are evaluated at the \emph{upper}
endpoint $b$ (so that they dominate every law in the cell) while the numerical
budget is compared against the \emph{lower} endpoint $a$ (so that division by
$\ell\ge a$ is legitimate).

%=====================================================================
\section{The lower bound $0.40$}
\label{sec:lower}
%=====================================================================

By Lemma~\ref{lem:reduction} it suffices to exhibit one admissible law, one
sample size and one threshold with normalized error at least $2/5$. The
witness is a standardized Bernoulli variable with success probability $2/5$ and
six summands. Every quantity in the argument is either an exact rational
number, an exact multiple of $\sqrt6$, or a normal probability for which one
rational upper bound is proved.

\subsection{The witness law and its moments}

\begin{definition}\label{def:witness}
Let $B$ be Bernoulli with $\PR(B=1)=2/5$ and set
\[
  X := \frac{5B-2}{\sqrt6},
  \qquad\text{i.e.}\qquad
  X = \frac{\sqrt6}{2} \ \text{w.p.}\ \frac25,
  \qquad
  X = -\frac{\sqrt6}{3} \ \text{w.p.}\ \frac35 .
\]
Write $\mu^\star$ for the law of $X$.
\end{definition}

\begin{lemma}\label{lem:witness-moments}
$\mu^\star$ is admissible, and
\[
  \E X = 0,\qquad \E X^2 = 1,\qquad
  \beta^\star := \E|X|^3 = \frac{13\sqrt6}{30},
  \qquad \frac{\sqrt6}{\beta^\star} = \frac{30}{13}.
\]
\end{lemma}

\begin{proof}
Integrability is automatic for a two-point law. With $p=2/5$, $q=3/5$,
$a=\sqrt6/2$, $b=-\sqrt6/3$ we have $pa+qb=\frac25\cdot\frac{\sqrt6}{2}
-\frac35\cdot\frac{\sqrt6}{3}=\frac{\sqrt6}{5}-\frac{\sqrt6}{5}=0$ and
$pa^2+qb^2=\frac25\cdot\frac64+\frac35\cdot\frac69=\frac35+\frac25=1$.
For the third absolute moment,
$p|a|^3+q|b|^3=\frac25\cdot\frac{6\sqrt6}{8}+\frac35\cdot\frac{6\sqrt6}{27}
=\frac{3\sqrt6}{10}+\frac{2\sqrt6}{15}=\frac{9\sqrt6+4\sqrt6}{30}
=\frac{13\sqrt6}{30}$. Finally
$\sqrt6/\beta^\star=30/13$.
\end{proof}

\begin{leanref}
\lean{witnessLaw}, \lean{witness\_admissible}, \lean{witness\_thirdAbsMoment},
\lean{witness\_normalizing\_ratio} in \file{BerryEsseen/BernoulliWitness.lean}.
The law is built as mathlib's \lean{bernoulliMeasure} on the two atoms, so the
product measure of Definition~\ref{def:sumlaw} is the six-fold product of the
actual coin.
\end{leanref}

\subsection{The event and its exact probability}

\begin{lemma}\label{lem:witness-event}
Let $X_1,\dots,X_6$ be independent copies of $X$ and
$W_6=(X_1+\dots+X_6)/\sqrt6$. If $K=B_1+\dots+B_6$ denotes the number of
successes, then
\[
  W_6 = \frac{5K-12}{6},
  \qquad\text{so}\qquad
  W_6\le-\tfrac13 \iff K\le2,
\]
and
\[
  \PR\bigl(W_6\le-\tfrac13\bigr) = \sum_{k=0}^{2}\binom6k
  \Bigl(\frac25\Bigr)^{k}\Bigl(\frac35\Bigr)^{6-k}
  = \frac{1701}{3125}.
\]
\end{lemma}

\begin{proof}
$X_j=\frac{\sqrt6}{6}(5B_j-2)$, so $\sum_j X_j=\frac{\sqrt6}{6}(5K-12)$ and
$W_6=(5K-12)/6$. Then $W_6\le-1/3$ iff $5K-12\le-2$ iff $K\le2$. The
probability is $3^6/5^6 + 6\cdot2\cdot3^5/5^6 + 15\cdot4\cdot3^4/5^6
=(729+2916+4860)/15625=8505/15625=1701/3125$.
\end{proof}

\begin{leanref}
\lean{normalizedSum\_witness}, \lean{normalizedSum\_witness\_le\_iff},
\lean{witness\_sumLaw\_eq\_map}, \lean{witness\_six\_probability} in
\file{BerryEsseen/BernoulliWitness.lean}. Rather than invoking a binomial
formula, the Lean proof evaluates the probability by summing over all
$2^6=64$ Boolean outcomes of the product measure:
\lean{witness\_finite\_sum\_rational} is a single exact rational identity
\[
  \sum_{b\in\{0,1\}^{6}} \mathbf 1\{\#b\le2\}\prod_{i=1}^{6}
  \bigl(b_i\,\tfrac25+(1-b_i)\tfrac35\bigr) = \frac{1701}{3125}
\]
discharged by \lean{decide +kernel}. The passage from the pushforward of the
product measure on $\R^6$ to the product measure on $\{0,1\}^6$ uses
mathlib's \lean{Measure.pi\_map\_pi} and \lean{Measure.pi\_singleton}.
\end{leanref}

\subsection{One rational estimate of a normal probability}

\begin{lemma}\label{lem:normal-third}
$\displaystyle \Phi(-\tfrac13) \le \frac12 - \frac{1325}{10206}$.
\end{lemma}

\begin{proof}
By symmetry $\Phi(0)=\tfrac12$, so
$\Phi(-\tfrac13)=\tfrac12-\int_{-1/3}^{0}\varphi$. On $[-\tfrac13,0]$ we have
$e^{-x^2/2}\ge1-\tfrac{x^2}{2}$, and $\pi<3.15$ gives
$(2\pi)^{-1/2}\ge\tfrac{25}{63}$. Hence
\[
  \int_{-1/3}^{0}\varphi
  \ \ge\ \frac{25}{63}\int_{-1/3}^{0}\Bigl(1-\frac{x^2}{2}\Bigr)dx
  = \frac{25}{63}\Bigl(\frac13-\frac{1}{162}\Bigr)
  = \frac{25}{63}\cdot\frac{53}{162}
  = \frac{1325}{10206}. \qedhere
\]
\end{proof}

\begin{leanref}
\lean{normal\_density\_coefficient\_lower},
\lean{normal\_density\_lower\_on\_small\_interval},
\lean{normalCDF\_neg\_third\_upper} in \file{BerryEsseen/GaussianBounds.lean};
the bound $(2\pi)^{-1/2}\ge25/63$ is derived from mathlib's decimal bounds on
$\pi$.
\end{leanref}

\subsection{Proof of Theorem~\ref{thm:main-lower}}

\begin{proof}[Proof of Theorem~\ref{thm:main-lower}]
By Lemmas~\ref{lem:witness-moments}--\ref{lem:normal-third},
\[
  \frac{\sqrt6}{\beta^\star}
  \Bigl|\PR\bigl(W_6\le-\tfrac13\bigr)-\Phi(-\tfrac13)\Bigr|
  \ \ge\ \frac{30}{13}\Bigl(\frac{1701}{3125}-\frac12+\frac{1325}{10206}\Bigr)
  \ =\ 0.401874\ldots \ >\ \frac25 .
\]
(The displayed rational number is
$\frac{30}{13}\cdot\frac{1701\cdot2\cdot10206 - 3125\cdot10206+1325\cdot6250}
{2\cdot3125\cdot10206}$, and the inequality with $2/5$ is a rational
inequality.) Lemma~\ref{lem:reduction} applied to $\mu^\star$, $n=6$ and
$x=-1/3$ gives $\Cii\ge2/5$.
\end{proof}

\begin{leanref}
\lean{lower\_witness\_normal\_gap} in \file{BerryEsseen/GaussianBounds.lean}
and \lean{BerryEsseen.berryEsseenConstant\_lower\_bound} in
\file{BerryEsseen/LowerBound.lean}. The Lean proof also carries out the
conversion $(0.40:\lean{ENNReal})=\lean{ENNReal.ofReal}(2/5)$ explicitly.
\end{leanref}

\begin{remark}
The witness is chosen so that the lattice jump of $W_6$ straddles the
threshold: $W_6$ takes the values $(5k-12)/6$ for $k=0,\dots,6$, i.e.\
$-2,-\tfrac76,-\tfrac13,\tfrac12,\tfrac43,\tfrac{13}{6},3$, and $-1/3$ is an
atom. Half of that atom's mass is the leading source of the discrepancy; the
skewness of the asymmetric Bernoulli law contributes in the same direction.
Section~\ref{sec:esseen-lower} makes both contributions quantitative in the
limit $n\to\infty$.
\end{remark}

\subsection{Esseen's stronger lower bound (mathematics only, not formalized)}
\label{sec:esseen-lower}

The repository also contains a complete mathematical proof of Esseen's lower
bound \cite{Esseen1956}
\begin{equation}\label{eq:esseen-lower}
  \Cii \ \ge\ c_E = \frac{\sqrt{10}+3}{6\sqrt{2\pi}} = 0.4097321837\ldots,
\end{equation}
in \file{docs/ESSEEN\_LOWER\_BOUND.md}. \textbf{This asymptotic argument is not
formalized in Lean}; the formalized lower bound is the weaker $0.40$ of
Theorem~\ref{thm:main-lower}, and the upper-bound formalization does not use
\eqref{eq:esseen-lower} in any way. We summarize the argument because it
explains the numerical target and shows which features of the Bernoulli family
produce it. The Bernoulli coefficient used below is the one recorded in the
introduction and equation~(7) of \cite{Zolotukhin2018}.

Fix a rational $p\in(0,\tfrac12)$, put $q=1-p$, $\sigma=\sqrt{pq}$, let $B$ be
Bernoulli$(p)$ and $X=(B-p)/\sigma$. Then $\E X=0$, $\E X^2=1$ and
\[
  \beta(p) = \frac{pq^3+qp^3}{(pq)^{3/2}} = \frac{p^2+q^2}{\sqrt{pq}} .
\]
Take $n$ through multiples of the denominator of $p$, so that $m:=np$ is an
integer; at the threshold $0$ we have
$\PR\bigl((X_1+\dots+X_n)/\sqrt n\le0\bigr)=\PR(K_n\le m)$ with
$K_n=\sum_j B_j$. The key asymptotic is
\begin{equation}\label{eq:esseen-key}
  \PR(K_n\le np) = \frac12 + \frac{2-p}{3\sqrt{2\pi npq}} + \bigO_p(n^{-1}).
\end{equation}

\paragraph{An exact integral for the binomial tail.}
For integers $0\le m<n$ let $G(t)=\sum_{j=0}^{m}\binom nj t^j(1-t)^{n-j}$.
Differentiating and using $j\binom nj=n\binom{n-1}{j-1}$ and
$(n-j)\binom nj=n\binom{n-1}{j}$, all terms cancel except the last:
\[
  G'(t) = -(n-m)\binom nm t^m(1-t)^{n-m-1}.
\]
Since $G(0)=1$ and $G(1)=0$, integration gives the exact identity
\begin{equation}\label{eq:beta-ratio}
  \PR(K_n\le m) = G(p)
  = \frac{\int_p^1 t^m(1-t)^{n-m-1}\,dt}{\int_0^1 t^m(1-t)^{n-m-1}\,dt} .
\end{equation}
No Stirling approximation is needed.

\paragraph{Laplace expansion at the common maximum.}
Let $L(t)=p\log t+q\log(1-t)$, so the integrand in \eqref{eq:beta-ratio} is
$e^{nL(t)}/(1-t)$. Then $L$ has its unique maximum at $t=p$, with
\[
  L'(p)=0,\qquad L''(p)=-\frac1{pq},\qquad L'''(p)=\frac{2(q-p)}{p^2q^2}.
\]
Substituting $t=p+\sigma y/\sqrt n$, Taylor's theorem gives
\[
  n\bigl(L(p+\sigma y/\sqrt n)-L(p)\bigr)
  = -\frac{y^2}{2} + \frac{q-p}{3\sigma\sqrt n}y^3 + \bigO_p(y^4/n),
  \qquad
  \frac{q}{1-t} = 1 + \frac{\sigma}{q\sqrt n}y + \bigO_p(y^2/n).
\]
Writing
$A_n(y)=\frac{q}{q-\sigma y/\sqrt n}\exp\bigl(n[L(p+\sigma y/\sqrt n)-L(p)]\bigr)$
on $I_n=(-p\sqrt n/\sigma,\,q\sqrt n/\sigma)$, the common factor
$e^{nL(p)}\sigma/(q\sqrt n)$ cancels in the ratio \eqref{eq:beta-ratio} and
\begin{equation}\label{eq:An-expansion}
  A_n(y) = e^{-y^2/2}\Bigl(1+\frac{P_p(y)}{\sqrt n}\Bigr)
  + (\text{integrable error}),
  \qquad
  P_p(y) = \frac{\sigma}{q}y + \frac{q-p}{3\sigma}y^3 .
\end{equation}
The polynomial $P_p$ is odd: its contribution vanishes over the whole line but
is positive over the positive half-line. That asymmetry is the entire source of
the improvement over the symmetric example.

\paragraph{Why integrating \eqref{eq:An-expansion} is legitimate.}
Pointwise Taylor expansion does not suffice; the error must be integrated,
including the omitted tails. Choose $\delta>0$ with
$[p-\delta,p+\delta]\subset(0,1)$ small enough that
$L(t)-L(p)\le-c_p(t-p)^2$ on that interval and that $L^{(4)}$ and the first two
derivatives of $q/(1-t)$ are bounded there. Split $I_n$ at $R_n=n^{1/12}$ and
at $\delta\sqrt n/\sigma$.
\begin{itemize}[leftmargin=2.2em]
\item \emph{Central region $|y|\le R_n$.} With $|e^z-1-z|\le Cz^2$ for
  $|z|\le1$, the two expansions give
  \[
    \Bigl|A_n(y)-e^{-y^2/2}\bigl(1+P_p(y)/\sqrt n\bigr)\Bigr|
    \le \frac{C_p}{n}e^{-y^2/2}\bigl(y^2+y^4+y^6\bigr),
  \]
  where $y^4/n$ comes from the cubic remainder, $y^6/n$ from squaring the cubic
  term and $y^2/n$ from the prefactor. The Gaussian moments on the right are
  integrable over all of $\R$, so the total contribution is $\bigO_p(n^{-1})$.
\item \emph{Intermediate region $R_n<|y|\le\delta\sqrt n/\sigma$.} Here
  $A_n(y)\le C_pe^{-c_p\sigma^2y^2}$, whose integral over $|y|>R_n$ is smaller
  than every inverse power of $n$ because $R_n^2=n^{1/6}$; the tails of the
  Gaussian polynomial in \eqref{eq:An-expansion} have the same property.
\item \emph{Outer region, including the singular endpoint $t\to1$.} Since $L$
  tends to $-\infty$ at $0$ and $1$ and has its unique maximum at $p$, there is
  $\eta_p>0$ with $L(t)-L(p)\le-\eta_p$ for $|t-p|\ge\delta$. Hence for
  $n\ge n_0$,
  \[
    \int_{|t-p|\ge\delta}\frac{e^{n(L(t)-L(p))}}{1-t}\,dt
    \le e^{-n\eta_p/2}\int_0^1\frac{e^{(n_0/2)(L(t)-L(p))}}{1-t}\,dt,
  \]
  and the last integral converges because near $0$ the integrand behaves like
  $t^{n_0p/2}$ and near $1$ like $(1-t)^{n_0q/2-1}$ with exponent
  $>-1$. Changing variables multiplies the bound by $q\sqrt n/\sigma$, which
  leaves exponential decay.
\end{itemize}

\paragraph{Evaluating the coefficient.}
Since $P_p$ is odd, $\int_{I_n}A_n = \sqrt{2\pi}+\bigO_p(n^{-1})$. On the
positive half-line $\int_0^\infty ye^{-y^2/2}dy=1$ and
$\int_0^\infty y^3e^{-y^2/2}dy=2$, so
\[
  \int_{I_n\cap[0,\infty)}A_n
  = \frac{\sqrt{2\pi}}{2}
    + \frac1{\sqrt n}\Bigl(\frac{\sigma}{q}+\frac{2(q-p)}{3\sigma}\Bigr)
    + \bigO_p(n^{-1})
  = \frac{\sqrt{2\pi}}{2} + \frac{2-p}{3\sigma\sqrt n} + \bigO_p(n^{-1}),
\]
using $\sigma/q=p/\sigma$. Dividing by the denominator of
\eqref{eq:beta-ratio}, which tends to $\sqrt{2\pi}>0$, proves
\eqref{eq:esseen-key}.

\paragraph{Optimizing.}
Since $\Phi(0)=\tfrac12$, Lemma~\ref{lem:reduction} and
\eqref{eq:esseen-key} give, for each fixed rational $p\in(0,\tfrac12)$,
\[
  \Cii \ \ge\ \mathcal E(p) := \frac{2-p}{3\sqrt{2\pi}\,[p^2+(1-p)^2]} .
\]
The function $\mathcal E$ is continuous and the rationals are dense, so the
bound extends to all real $p\in(0,\tfrac12)$; in particular to
$p_*=2-\tfrac{\sqrt{10}}{2}$. Note the order of limits: for each fixed
rational $p$ one first lets $n\to\infty$ through the allowed multiples, and only
afterwards uses continuity of the explicit function $\mathcal E$; no
nonuniform limits are interchanged and no integrality of $np_*$ is needed.
Differentiating $(2-p)/(2p^2-2p+1)$ gives a derivative with the sign of
$2p^2-8p+3$, positive before $p_*$ and negative after it on $(0,\tfrac12)$, so
$p_*$ maximizes. At $p_*$ one has $2-p_*=\sqrt{10}/2$ and
$p_*^2+(1-p_*)^2=10-3\sqrt{10}$, whence
\[
  \Cii \ \ge\ \mathcal E(p_*)
  = \frac{\sqrt{10}}{6\sqrt{2\pi}\,(10-3\sqrt{10})}
  = \frac{\sqrt{10}+3}{6\sqrt{2\pi}} = c_E .
\]

\paragraph{Where the improvement comes from.}
The coefficient in \eqref{eq:esseen-key} can be rewritten
\[
  \frac1{\sqrt{2\pi pq}}\Bigl(\frac12 + \frac{q-p}{6}\Bigr):
\]
the term $\tfrac12$ is the half-jump of the lattice distribution at its mean,
and $(q-p)/6$ is the skewness correction. For $p<\tfrac12$ the two reinforce
each other at the inclusive threshold $K_n\le np$. Dividing by $\beta(p)$ and
balancing the two contributions produces $p_*$ and the $\sqrt{10}$.

%=====================================================================
\section{Cosine anchors: retaining the two-point law}
\label{sec:anchors}
%=====================================================================

This section and the next two develop the analytic estimates on which every
region of the upper bound rests. Throughout, $X$ satisfies
\eqref{eq:admissible}, $Y=|X|$, $\delta=\beta-1$ and $f(t)=\E e^{itX}$.

The point of the estimates is stated first. The classical third-order bound
$|f(t)-(1-t^2/2)|\le\beta|t|^3/6$ is insensitive to the structure of the law.
But $\beta=1$ forces $|X|=1$ a.s., and then the zero mean forces
$\PR(X=1)=\PR(X=-1)=\tfrac12$ and $f(t)=\cos t$ exactly. An estimate that
compares $f$ with $\cos$ instead of with a quadratic therefore retains
information that is lost otherwise, and every error term below carries a factor
vanishing as $\delta\downarrow0$.

\subsection{The moment excess}

\begin{lemma}\label{lem:beta-ge-one}
Let $w(y)=(y-1)^2(2y+1)=2y^3-3y^2+1$. Then $w\ge0$ on $[0,\infty)$ and
\[
  \boxed{\ \E\,w(Y) = 2\delta = 2(\beta-1).\ }
\]
Consequently $\beta\ge1$ (so $\beta>0$), with $\beta=1$ if and only if
$Y=1$ almost surely. Moreover $\E Y\le1$.
\end{lemma}

\begin{proof}
Expanding, $\E w(Y)=2\E Y^3-3\E Y^2+1=2\beta-3+1=2(\beta-1)$, using
$\E Y^2=\E X^2=1$ and $\E Y^3=\beta$. Nonnegativity of $w$ on $[0,\infty)$ is
clear from the factored form. If $\delta=0$ then $\E w(Y)=0$ with $w(Y)\ge0$,
so $w(Y)=0$ a.s., i.e.\ $Y=1$ a.s. Finally $2Y\le Y^2+1$ integrates to
$2\E Y\le 2$.
\end{proof}

\begin{leanref}
\lean{momentExcessWeight}, \lean{integral\_momentExcessWeight},
\lean{one\_le\_thirdAbsMoment}, \lean{thirdAbsMoment\_pos},
\lean{integral\_abs\_le\_one},
\lean{ae\_abs\_eq\_one\_of\_thirdAbsMoment\_eq\_one} in
\file{BerryEsseen/MomentExcess.lean}.
\end{leanref}

\subsection{Two weighted Cauchy--Schwarz estimates}

\begin{lemma}\label{lem:weighted-cs}
\[
  \boxed{\ \E\bigl[Y|Y^2-1|\bigr] \le \sqrt{\delta\,(\beta+\tfrac53)},
  \qquad
  \E\bigl[Y|Y-1|\bigr] \le \sqrt{\delta}.\ }
\]
\end{lemma}

\begin{proof}
For the first, put $\rho(y)=\tfrac{y^3}{2}+\tfrac{7y^2}{9}+\tfrac{y}{18}$, so
$\rho\ge0$ on $[0,\infty)$. The polynomial identity
\[
  w(y)\rho(y) - \bigl(y|y^2-1|\bigr)^2 = \frac{y(y-1)^4}{18} \ \ge 0
  \qquad (y\ge0)
\]
holds pointwise. Writing $U=\sqrt{w(Y)}$, $V=\sqrt{\rho(Y)}$ and using
$Y|Y^2-1|\le UV$ pointwise, Cauchy--Schwarz gives
\[
  \bigl(\E[Y|Y^2-1|]\bigr)^2 \le \bigl(\E[UV]\bigr)^2
  \le \E U^2\cdot\E V^2 = \E w(Y)\cdot\E\rho(Y)
  = 2\delta\Bigl(\frac\beta2+\frac56\Bigr) = \delta\Bigl(\beta+\frac53\Bigr),
\]
using
$\E\rho(Y)=\tfrac\beta2+\tfrac79+\tfrac{\E Y}{18}\le\tfrac\beta2+\tfrac56$ by
$\E Y\le1$ and $\tfrac79+\tfrac1{18}=\tfrac{15}{18}=\tfrac56$.

For the second, the pointwise identity
$w(y)\tfrac y2-\bigl(y|y-1|\bigr)^2=\tfrac{y(y-1)^2}{2}\ge0$ and
Cauchy--Schwarz give
$\bigl(\E[Y|Y-1|]\bigr)^2\le\E w(Y)\cdot\tfrac{\E Y}{2}\le\delta$.
\end{proof}

\begin{leanref}
\lean{cubicWeight}, \lean{cubicWeight\_identity}, \lean{weighted\_cubic\_sq\_le},
\lean{integral\_abs\_weighted\_sq\_bound}, \lean{weighted\_linear\_sq\_le},
\lean{integral\_abs\_weighted\_linear\_bound} in
\file{BerryEsseen/WeightedMoments.lean}; the abstract Cauchy--Schwarz step in
the form ``a pointwise quadratic inequality gives an integrated one'' is
\lean{integral\_sq\_le\_mul\_of\_quadratic\_nonneg} in
\file{BerryEsseen/IntegralQuadratic.lean}. Only a third moment is used: all
majorants have degree at most three, and the products are handled pointwise
before integration.
\end{leanref}

\subsection{One integral identity gives both sine estimates}

\begin{lemma}\label{lem:sine-identity}
For all $t,x\in\R$,
\begin{equation}\label{eq:sine-fundamental}
  \sin(tx) - x\sin t = tx\int_0^1\bigl[\cos(txv)-\cos(tv)\bigr]dv .
\end{equation}
Consequently
\begin{equation}\label{eq:sine-cubic}
  \bigl|\sin(tx)-x\sin t\bigr| \le \frac{|t|^3}{6}\,|x|\,|x^2-1| ,
\end{equation}
\begin{equation}\label{eq:sine-quadratic}
  \bigl|\sin(tx)-x\sin t\bigr| \le \frac{t^2}{2}\,|x|\,\bigl||x|-1\bigr| .
\end{equation}
\end{lemma}

\begin{proof}
Both sides of \eqref{eq:sine-fundamental} vanish at $t=0$ and have the same
$t$-derivative, by the fundamental theorem of calculus applied to
$v\mapsto\sin(txv)/ (tx)$ and $v\mapsto\sin(tv)$. For
\eqref{eq:sine-cubic} use $|\cos a-\cos b|\le\tfrac12|a^2-b^2|$, which follows
from the product-to-sum formula
$\cos a-\cos b=-2\sin\frac{a+b}{2}\sin\frac{a-b}{2}$ and $|\sin u|\le|u|$;
this makes the integrand at most $\tfrac12 t^2v^2|x^2-1|$, and
$\int_0^1v^2dv=\tfrac13$. For \eqref{eq:sine-quadratic} use
$|\cos a-\cos b|\le|a-b|$, giving integrand at most $|t|v\,|x-1|$ for
$x\ge0$ and $\int_0^1v\,dv=\tfrac12$; the case $x<0$ follows by oddness of both
sides in $x$ combined with $|{-x}|=|x|$.
\end{proof}

\begin{leanref}
\lean{abs\_cos\_sub\_cos\_le\_sq}, \lean{sin\_sub\_mul\_sin\_eq\_integral},
\lean{abs\_sin\_sub\_mul\_sin\_cubic},
\lean{abs\_sin\_sub\_mul\_sin\_quadratic} in
\file{BerryEsseen/TrigonometricAnchors.lean}.
\end{leanref}

\begin{proposition}[Imaginary anchor]\label{prop:imag}
\begin{equation}\label{eq:imag}
  \boxed{\ \bigl|\Ima f(t)\bigr| \le
  \min\Bigl\{
    \frac{|t|^3}{6}\sqrt{\delta(\beta+\tfrac53)},\
    \frac{t^2}{2}\sqrt{\delta},\
    \frac{\beta|t|^3}{6}
  \Bigr\}.\ }
\end{equation}
\end{proposition}

\begin{proof}
$\Ima f(t)=\E\sin(tX)$, and since $\E X=0$ we may subtract $X\sin t$ without
changing the expectation:
$\Ima f(t)=\E[\sin(tX)-X\sin t]$. Taking absolute values inside and applying
\eqref{eq:sine-cubic}, \eqref{eq:sine-quadratic} and
Lemma~\ref{lem:weighted-cs} gives the first two entries. The third is the
classical $|\sin u-u|\le|u|^3/6$ applied with $u=tX$, again using $\E X=0$.
\end{proof}

\begin{leanref}
\lean{charFun\_im\_cubic\_anchor}, \lean{charFun\_im\_quadratic\_anchor},
\lean{charFun\_im\_classical\_bound} in \file{BerryEsseen/ImaginaryPart.lean};
their minimum is packaged as \lean{BerryEsseen.imaginaryBound} with
\lean{charFun\_im\_le\_imaginaryBound} in
\file{BerryEsseen/CharacteristicBounds.lean}.
\end{leanref}

\subsection{Integrating the sine estimate gives the cosine anchor}

\begin{proposition}[Real anchor]\label{prop:real}
\begin{equation}\label{eq:real}
  \boxed{\ \bigl|\Rea f(t) - \cos t\bigr| \le \frac{\delta\,|t|^3}{6}.\ }
\end{equation}
\end{proposition}

\begin{proof}
For $y\ge0$ define the remainder with a quadratic correction
\[
  R_t(y) := \cos(ty)-\cos t + \frac{t\sin t}{2}\,(y^2-1),
\]
so that $R_t(1)=0$ and $R_t'(y)=-t\bigl[\sin(ty)-y\sin t\bigr]$. By
\eqref{eq:sine-quadratic}, $|R_t'(y)|\le\tfrac{|t|^3}{2}y|y-1|$. Integrating
from $1$ to $y$ in the appropriate direction and using
\[
  \int_{\min(1,y)}^{\max(1,y)} s|s-1|\,ds = \frac{2y^3-3y^2+1}{6}
  = \frac{w(y)}{6}
\]
gives $|R_t(y)|\le\tfrac{|t|^3}{12}w(y)$ for all $y\ge0$. Now substitute
$y=Y=|X|$ and take expectations. Since $\cos(tY)=\cos(tX)$ and
$\Rea f(t)=\E\cos(tX)$, and since $\E Y^2=1$ kills the quadratic correction,
\[
  \bigl|\Rea f(t)-\cos t\bigr| = \bigl|\E R_t(Y)\bigr|
  \le \frac{|t|^3}{12}\,\E w(Y) = \frac{|t|^3}{12}\cdot2\delta
  = \frac{\delta|t|^3}{6}. \qedhere
\]
\end{proof}

\begin{leanref}
\lean{cosineRemainder}, \lean{abs\_cosineRemainder\_le},
\lean{integral\_cosineRemainder}, \lean{charFun\_re\_anchor} in
\file{BerryEsseen/RealPart.lean}; the weight integral is
\lean{integral\_linear\_weight}. A pleasant consequence, proved as
\lean{charFun\_eq\_cos\_of\_thirdAbsMoment\_eq\_one} in
\file{BerryEsseen/CharacteristicBounds.lean}, is that $\beta=1$ implies
$f(t)=\cos t$ for all $t$.
\end{leanref}

\subsection{Combining the coordinates}

Propositions~\ref{prop:imag} and \ref{prop:real} bound the two coordinates of
$f(t)$ separately; combining them in the Euclidean norm gives bounds on
$|f(t)|$ and on $|f(t)-g(t)|$ that are better near $\beta=1$ than either
coordinate alone would suggest.

\begin{lemma}\label{lem:coordinate-combination}
If $|\Rea z|\le a$ and $|\Ima z|\le b$ then $|z|\le\sqrt{a^2+b^2}$.
Consequently, writing $I(\beta,t)$ for the right-hand side of
\eqref{eq:imag},
\begin{align}
  |f(t)| &\le \min\Bigl\{1,\ \sqrt{\bigl(|\cos t|+\tfrac{\delta|t|^3}{6}\bigr)^2
    + I(\beta,t)^2}\Bigr\}, \label{eq:mod-anchor}\\
  |f(t)-c| &\le \sqrt{\bigl(|\cos t-c|+\tfrac{\delta|t|^3}{6}\bigr)^2
    + I(\beta,t)^2} \qquad (c\in\R). \label{eq:err-anchor}
\end{align}
\end{lemma}

\begin{leanref}
\lean{norm\_le\_sqrt\_of\_coordinate\_bounds}, \lean{norm\_charFun\_anchor},
\lean{norm\_charFun\_sub\_real\_anchor} in
\file{BerryEsseen/CharacteristicBounds.lean}.
\end{leanref}

Two comparison functions are used for $c$ in \eqref{eq:err-anchor}: the
Gaussian $c=g(t)=e^{-t^2/2}$, for which one needs the Gaussian--cosine
discrepancy of Lemma~\ref{lem:gauss-cos} below, and, in the finite-sample
region, nothing else.

\begin{lemma}[Gaussian versus cosine]\label{lem:gauss-cos}
For all real $u$,
\[
  e^{-u^2/2}-\cos u = \int_0^u \sin(u-v)\,v^2\,e^{-v^2/2}\,dv,
\]
and hence for $0\le u\le\pi$,
\[
  0 \ \le\ e^{-u^2/2}-\cos u \ \le\ \frac{u^4}{12}
  \ \le\ \frac{u^4}{8}.
\]
There is also an explicit polynomial refinement
$0\le e^{-u^2/2}-\cos u\le u^4Q(u^2)$ on $[0,\pi]$ with $Q$ rational and
nonnegative.
\end{lemma}

\begin{proof}
The function $F(v)=\bigl(\cos(u-v)-v\sin(u-v)\bigr)e^{-v^2/2}$ has derivative
$F'(v)=\sin(u-v)\,v^2e^{-v^2/2}$, while $F(0)=\cos u$ and $F(u)=e^{-u^2/2}$;
the identity is $\int_0^uF'=F(u)-F(0)$.
For $0\le v\le u\le\pi$ we have $\sin(u-v)\ge0$, whence nonnegativity. For the
upper bound use $\sin(u-v)\le u-v$ and $e^{-v^2/2}\le1$ and
$\int_0^u(u-v)v^2dv=u^4/12$.
\end{proof}

\begin{leanref}
\lean{gaussian\_cosine\_integral}, \lean{gaussian\_cosine\_nonneg},
\lean{gaussian\_cosine\_le\_quartic} in \file{BerryEsseen/GaussianCosine.lean};
the polynomial refinement is \lean{gaussianCosineQ} with
\lean{gaussian\_cosine\_le\_Q} in \file{BerryEsseen/GaussianPolynomial.lean}.
\end{leanref}

\begin{definition}[The two all-frequency majorants]\label{def:bounds}
For $\beta\ge1$ and $t\in\R$ set
\begin{align*}
  M(\beta,t) &:= \min\Bigl\{
     \min\bigl\{1,\ \sqrt{(|\cos t|+\tfrac{(\beta-1)|t|^3}{6})^2+I(\beta,t)^2}\bigr\},\
     \sqrt{1-2t^2h((\beta+1)|t|)}\Bigr\},\\
  E(\beta,t) &:= \min\Bigl\{
     \sqrt{\bigl(|\cos t-e^{-t^2/2}|+\tfrac{(\beta-1)|t|^3}{6}\bigr)^2+I(\beta,t)^2},\
     \frac{\beta|t|^3}{6}+\frac{t^4}{8}\Bigr\},
\end{align*}
where $h$ is the convex minorant of Section~\ref{sec:modulus}. Then
$|f(t)|\le M(\beta,t)$ and $|f(t)-g(t)|\le E(\beta,t)$ for every admissible law
with third absolute moment $\beta$, at every real frequency.
\end{definition}

\begin{leanref}
\lean{BerryEsseen.modulusBound}, \lean{BerryEsseen.errorBound},
\lean{norm\_charFun\_le\_modulusBound},
\lean{norm\_charFun\_sub\_gaussian\_le\_errorBound} in
\file{BerryEsseen/SumBounds.lean}. The second entry of $E$ is the classical
estimate, proved in \file{BerryEsseen/ClassicalError.lean} as
\lean{norm\_charFun\_sub\_gaussian\_classical}. Each minimum combines
independently valid inequalities and no cutoff restriction on $t$ is imposed.
\end{leanref}

%=====================================================================
\section{A global exponential modulus bound}
\label{sec:modulus}
%=====================================================================

The estimate of this section is what makes the argument valid at \emph{all}
frequencies, including for lattice laws whose characteristic functions return
to modulus one. It is the reason no ``Gaussian decay outside a small interval''
assumption appears anywhere.

\begin{theorem}\label{thm:modulus}
Let $X$ satisfy \eqref{eq:admissible} and let $h$ be the function of
Definition~\ref{def:minorant}. Then for all $t\in\R$,
\[
  \boxed{\ |f(t)|^2 \le 1-2t^2h\bigl((\beta+1)|t|\bigr),
  \qquad
  |f(t)| \le \exp\bigl\{-t^2h((\beta+1)|t|)\bigr\}.\ }
\]
\end{theorem}

\begin{leanref}
\lean{norm\_charFun\_sq\_minorant} and \lean{norm\_charFun\_le\_minorant\_exp}
in \file{BerryEsseen/FullModulus.lean}.
\end{leanref}

\subsection{Symmetrization}

\begin{lemma}\label{lem:symm}
Let $X'$ be an independent copy of $X$ and $D=X-X'$. Then
\[
  |f(t)|^2 = \E\cos(tD), \qquad \E D^2 = 2, \qquad
  \E|D|^3 \le 2(\beta+1) .
\]
\end{lemma}

\begin{proof}
$|f(t)|^2=f(t)\overline{f(t)}=\E e^{itX}\E e^{-itX'}=\E e^{itD}=\E\cos(tD)$,
the imaginary part vanishing by symmetry of $D$. Also
$\E D^2=\E X^2+\E X'^2-2(\E X)^2=2$. For the third moment, the pointwise
inequality $|X-X'|\le|X|+|X'|$ gives
$|D|^3\le(X-X')^2(|X|+|X'|)$, and expanding
\[
  \E\bigl[(X-X')^2(|X|+|X'|)\bigr]
  = 2\E|X|^3 + 2\E|X| = 2\beta + 2\E|X| \le 2(\beta+1),
\]
where the cross terms containing $\E X$ or $\E X'$ vanish, the remaining ones
use $\E X^2=1$, and $\E|X|\le1$ by Cauchy--Schwarz (or by
Lemma~\ref{lem:beta-ge-one}).
\end{proof}

\begin{leanref}
\lean{integral\_cos\_symm}, \lean{integral\_symm\_sq},
\lean{integral\_symm\_abs\_cube\_le} in
\file{BerryEsseen/Symmetrization.lean}; the integrability of the symmetrized
cube uses the elementary majorant
$|x-y|^3\le4(|x|^3+|y|^3)$.
\end{leanref}

Symmetrization converts a complex modulus into a real cosine expectation while
retaining second- and third-moment information --- exactly the two pieces of
data the minorant argument consumes.

\subsection{The convex minorant}

\begin{definition}\label{def:minorant}
Let $r=\tfrac92$, $G(s)=\dfrac{1-\cos s}{s^2}$ for $s>0$, and
\[
  a := \frac{81}{80\,(1-\cos r)},
  \qquad
  h(s) := \begin{cases}
    \dfrac12-\dfrac{s}{10}, & s\le r,\\[2mm]
    a\,G(s), & r< s\le 2\pi,\\[1mm]
    0, & s>2\pi .
  \end{cases}
\]
The linear branch is used also for $s<0$, so that $h$ is defined and convex on
all of $\R$; only nonnegative arguments enter the probabilistic estimate.
\end{definition}

\begin{proposition}\label{prop:minorant}
The function $h$ is convex on $\R$, nonnegative, nonincreasing,
$h(s)\le\tfrac12$ for $s\ge0$, $h(s)\le G(s)$ for $s>0$, and
$\tfrac56<a\le1$. Consequently, for all real $v$,
\begin{equation}\label{eq:cos-minorant}
  \cos v \le 1 - v^2 h(|v|).
\end{equation}
\end{proposition}

\begin{proof}[Proof sketch]
The pieces are checked as follows.
\begin{enumerate}[leftmargin=2.2em]
\item \emph{Linear branch below $G$.} The global inequality
  $\cos s-1+\tfrac{s^2}{2}\le\tfrac{|s|^3}{10}$ rearranges to
  $G(s)\ge\tfrac12-\tfrac{s}{10}$ for $s>0$.
\item \emph{Scale.} Elementary sine bounds at $\tfrac{3\pi}{2}-r$ give
  $-\tfrac3{14}\le\cos r\le-\tfrac18$, so $\tfrac56<a\le1$; hence
  $aG\le G$, and the zero branch is trivially below $G$.
\item \emph{Continuity at the joins.} $aG(r)=\tfrac1{20}=\tfrac12-\tfrac r{10}$
  and $G(2\pi)=0$.
\item \emph{Monotonicity and convexity of $G$ on $[r,2\pi]$.} Differentiating,
  \[
    G'(s) = \frac{s\sin s-2(1-\cos s)}{s^3}\le0,
    \qquad
    G''(s) = \frac{(s^2-6)\cos s-4s\sin s+6}{s^4}\ge0 .
  \]
  For $G''$, split at $\tfrac{3\pi}{2}$: on $[r,\tfrac{3\pi}{2}]$ the distance
  to $\tfrac{3\pi}{2}$ is at most $\tfrac12$, so $\sin s\le-\tfrac34$ by
  $\cos v\ge1-\tfrac{v^2}{2}$, and with $r\le s\le5$, $\cos s\ge-1$, the
  numerator is at least $12-s^2+3s\ge\tfrac12$; on $[\tfrac{3\pi}{2},2\pi]$
  every nonconstant term of the numerator is nonnegative.
\item \emph{No downward jump of slope.} At the first join
  $G'(r)\ge-2G(r)$, hence $aG'(r)\ge-\tfrac1{10}$, the slope of the linear
  branch; at the second join $G'(2\pi)=0$. The Lean proof does not assume
  differentiability at a join: it constructs an explicit supporting line at
  every point and applies the criterion ``a function with a supporting line
  everywhere is convex''.
\end{enumerate}
Finally \eqref{eq:cos-minorant} is exactly $h(|v|)\le G(|v|)$ for $v\ne0$, and
trivial at $v=0$.
\end{proof}

\begin{leanref}
\lean{cosineProfile} and its two derivative lemmas in
\file{BerryEsseen/CosineProfile.lean}; the splicing criterion
\lean{convexOn\_of\_supporting\_lines} and \lean{convexOn\_convexSplice} in
\file{BerryEsseen/ConvexSplice.lean}; \lean{minorantScale},
\lean{cosineMinorant}, \lean{cos\_splice\_bounds}, \lean{minorantScale\_bounds},
\lean{convexOn\_cosineMinorant}, \lean{antitone\_cosineMinorant},
\lean{cosineMinorant\_le\_half}, \lean{cosineMinorant\_supporting\_line},
\lean{cos\_le\_minorant} in \file{BerryEsseen/FullMinorant.lean}; the cubic
cosine inequality \lean{cosine\_tenth} in
\file{BerryEsseen/CosineMinorant.lean}.
\end{leanref}

\subsection{One weighted Jensen step}

\begin{proof}[Proof of Theorem~\ref{thm:modulus}]
Put $z=(\beta+1)|t|$ and choose, by convexity, a supporting line to $h$ at
$z$: $h(y)\ge h(z)+m(y-z)$ for all $y$, with $m\le0$ because $h$ is
nonincreasing. By Lemma~\ref{lem:symm} and \eqref{eq:cos-minorant} applied with
$v=tD$,
\[
  1-|f(t)|^2 = \E\bigl[1-\cos(tD)\bigr]
  \ \ge\ t^2\,\E\bigl[D^2h(|t|\,|D|)\bigr] .
\]
Using the supporting line with $y=|t||D|$,
\[
  \E\bigl[D^2h(|t||D|)\bigr]
  \ \ge\ h(z)\,\E D^2 + m\bigl(|t|\,\E|D|^3 - z\,\E D^2\bigr)
  = 2h(z) + m\bigl(|t|\E|D|^3-2z\bigr) \ \ge\ 2h(z),
\]
since $|t|\E|D|^3\le 2(\beta+1)|t|=2z$ and $m\le0$. This is the first
assertion. All integrals exist: $h$ is bounded on $[0,\infty)$ and the
supporting-line expression uses only the second and third moments already
established. The second assertion follows from $1-v\le e^{-v}$ and taking
square roots.
\end{proof}

\begin{remark}
The step above is Jensen's inequality with weights $D^2/2$, written so that no
new probability measure has to be constructed. It is also where the specific
combination $(\beta+1)|t|$ arises: it is precisely $\tfrac12|t|\E|D|^3/\E D^2$
up to the bound of Lemma~\ref{lem:symm}.
\end{remark}

Since $h$ vanishes for $s>2\pi$, Theorem~\ref{thm:modulus} gives no decay for
$(\beta+1)|t|>2\pi$. This is unavoidable --- for the symmetric two-point law
$|f(2\pi k)|=1$ --- and it is the reason the smoothing cutoff $T$ is always
chosen so that the largest frequency $T\tau/\sqrt n$ entering the bound stays
inside the useful range.

%=====================================================================
\section{Prawitz smoothing, proved from scratch}
\label{sec:prawitz}
%=====================================================================

Every region of the upper bound uses the same smoothing inequality: for any
cutoff $T>0$ and any split $0<\tau\le1$, the Kolmogorov distance between a law
with a finite first moment and the standard normal law is bounded by four
explicit positive integrals in which the law appears only through its
characteristic function on the bounded band $[0,T]$. This is Prawitz's
inequality \cite{Prawitz1972}; the development proves it, including all
integrability and endpoint issues, without importing it.

\subsection{A pointwise majorant of the step function}

\begin{definition}\label{def:H}
For $0<s<1$ let $B(s)=(1-s)\cot(\pi s)+\tfrac1\pi$ as in
Convention~\ref{conv:kernel} and set
\[
  H(x) := 2\int_0^1 B(s)\sin(2\pi xs)\,ds,
  \qquad
  S(x) := \frac{1+H(x)+\sinc(\pi x)^2}{2}.
\]
\end{definition}

Although $B$ has a pole at $s=0$, the product $B(s)\sin(2\pi xs)$ is
integrable: Jordan's inequality on the two half-intervals $(0,\tfrac12]$ and
$[\tfrac12,1)$ gives
\begin{equation}\label{eq:B-bound}
  |B(s)\sin(2\pi xs)| \le \pi|x|+2 \qquad (0<s<1),
\end{equation}
so $H$ is well defined, odd, and $H(0)=0$; endpoint values of the integrand are
irrelevant to the integral.

\begin{theorem}\label{thm:majorant}
For all $x\in\R$,
\[
  \boxed{\ \mathbf 1_{[0,\infty)}(x) \ \le\ S(x).\ }
\]
Moreover $S(0)=1$, $S$ is measurable, and $|S(x)|\le\pi|x|+3$.
\end{theorem}

The proof has three steps, each elementary.

\paragraph{Step 1: a shift recurrence.}
The sine difference identity cancels the cotangent denominator:
\[
  \cot(\pi s)\bigl[\sin(2\pi(x+1)s)-\sin(2\pi xs)\bigr]
  = \cos(2\pi(x+1)s)+\cos(2\pi xs) .
\]
Only two elementary integrals then remain,
\[
  \int_0^1\sin(ks)\,ds = \frac{1-\cos k}{k},
  \qquad
  \int_0^1(1-s)\cos(ks)\,ds = \frac{1-\cos k}{k^2}
  \qquad (k\ne0),
\]
and substitution and cancellation give, for $x>0$,
\begin{equation}\label{eq:shift}
  \boxed{\ H(x+1)-H(x) = \frac{c(x)}{x^2(x+1)^2},
  \qquad c(x) := \Bigl(\frac{\sin\pi x}{\pi}\Bigr)^2 \ge 0.\ }
\end{equation}
The same identity at $x=0$ gives $H(1)=1$, and since $c(n)=0$ for every
integer, $H(n)=1$ for every positive integer $n$.

\begin{leanref}
\lean{integral\_sin\_linear}, \lean{integral\_triangular\_cos},
\lean{prawitz\_kernel\_shift\_identity}, \lean{prawitz\_integrand\_bound},
\lean{prawitzSineApprox\_shift} in \file{BerryEsseen/PrawitzKernel.lean};
\lean{prawitzSineApprox\_one}, \lean{prawitzSineApprox\_nat\_succ} in
\file{BerryEsseen/PrawitzMajorant.lean}. The Lean name for $H$ is
\lean{prawitzSineApprox} and for $B$ is \lean{prawitzSineKernel}.
\end{leanref}

\paragraph{Step 2: the limit needs only Riemann--Lebesgue.}
Fix $x\in\R$. Iterating the difference formula, or rather applying the
integrand identity directly,
\[
  H(x+n)-H(n)
  = 4\int_0^1 B(s)\sin(\pi xs)\,\cos\bigl(2\pi s(n+x/2)\bigr)\,ds .
\]
The amplitude $s\mapsto B(s)\sin(\pi xs)$ is integrable on $(0,1)$ by
\eqref{eq:B-bound}, so the Riemann--Lebesgue lemma gives
$H(x+n)-H(n)\to0$; since $H(n)=1$,
\[
  H(x+n) \longrightarrow 1 \qquad (n\to\infty).
\]
No improper sine integral is evaluated and no cotangent Fourier series is
summed.

\begin{leanref}
\lean{tendsto\_integral\_mul\_cos}, \lean{tendsto\_interval\_mul\_cos} in
\file{BerryEsseen/OscillatoryIntegral.lean} (the real interval form of
mathlib's Riemann--Lebesgue theorem);
\lean{tendsto\_prawitzSineApprox\_nat\_shift} in
\file{BerryEsseen/PrawitzMajorant.lean}.
\end{leanref}

\paragraph{Step 3: two monotone sequences.}
Fix $x>0$ and write $c=c(x)\ge0$; integer shifts leave $c$ unchanged. By
\eqref{eq:shift} the sequence $n\mapsto H(x+n)$ is nondecreasing with limit
$1$, hence
\[
  H(x)\le 1 .
\]
Adding the correction $c/(x+n)^2$ produces a nonincreasing sequence: a direct
computation from \eqref{eq:shift} gives
\[
  \Bigl[H(x+n+1)+\frac{c}{(x+n+1)^2}\Bigr]
  -\Bigl[H(x+n)+\frac{c}{(x+n)^2}\Bigr]
  = -\frac{2c}{(x+n)(x+n+1)^2} \le 0 .
\]
This corrected sequence dominates $H(x+n)$ and has the same limit $1$, so
\[
  1 \le H(x) + \frac{c(x)}{x^2} = H(x)+\sinc(\pi x)^2 ,
\]
that is, $S(x)\ge1$ for $x>0$. For $x<0$ oddness of $H$ together with
$H(-x)\le1$ gives $H(x)\ge-1$, hence $S(x)\ge0$; and $S(0)=1$ because
$H(0)=0$ and $\sinc(0)=1$. This proves Theorem~\ref{thm:majorant}; the growth
bound $|S|\le\pi|x|+3$ follows from \eqref{eq:B-bound} and
$|\sinc|\le1$.

\begin{leanref}
The abstract two-sequence argument is
\lean{bounds\_of\_kernel\_recurrence} in
\file{BerryEsseen/KernelRecurrence.lean}; \lean{prawitzStepApprox},
\lean{step\_le\_prawitzStepApprox}, \lean{abs\_prawitzStepApprox\_le},
\lean{measurable\_prawitzStepApprox} in
\file{BerryEsseen/PrawitzMajorant.lean}.
\end{leanref}

\subsection{From the majorant to two-sided bounds on a distribution function}

\begin{proposition}\label{prop:prawitz-prob}
Let $\nu$ be a probability measure on $\R$ with $\int|y|\,d\nu<\infty$, let
$F(a)=\nu(({-\infty},a])$ and let $k>0$. Then
\[
  1 - \int S\bigl(k(y-a)\bigr)d\nu(y)
  \ \le\ F(a) \ \le\
  \int S\bigl(k(a-y)\bigr)d\nu(y).
\]
\end{proposition}

\begin{proof}
Apply Theorem~\ref{thm:majorant} pointwise at $x=k(a-y)$: the left side is
$\mathbf 1\{y\le a\}$, and integrating gives the upper bound. The lower bound is
the reflected inequality. Both expectations exist by the growth bound
$|S(x)|\le\pi|x|+3$ and the assumed first moment. An atom at $a$ is handled
explicitly: the upper pointwise bound equals $1$ at $x=0$ and the reflected
lower bound equals $0$ there, which is exactly what right-continuity of $F$
requires.
\end{proof}

\begin{leanref}
\lean{cdf\_le\_integral\_prawitzStepApprox},
\lean{integral\_prawitzStepApprox\_reflected\_le\_cdf} in
\file{BerryEsseen/PrawitzProbability.lean}.
\end{leanref}

\subsection{The Fourier form}

Interchanging the probability integral with the $s$-integral in
Definition~\ref{def:H} turns Proposition~\ref{prop:prawitz-prob} into a
statement about the characteristic function. The interchange is justified by
the bound \eqref{eq:B-bound}, which supplies an integrable majorant linear in
$|y|$; in addition
\[
  \sinc(\pi z)^2 = 2\int_0^1(1-s)\cos(2\pi zs)\,ds .
\]

\begin{proposition}\label{prop:prawitz-fourier}
Let $\nu$ be a probability measure with finite first moment and characteristic
function $\hat\nu$, and let $T>0$. With $K_\sigma$ as in
Convention~\ref{conv:kernel},
\[
  \frac12+\int_0^1\Rea\Bigl\{K_{-1}(s)e^{-iTsx}\hat\nu(Ts)\Bigr\}ds
  \ \le\ F(x) \ \le\
  \frac12+\int_0^1\Rea\Bigl\{K_{+1}(s)e^{-iTsx}\hat\nu(Ts)\Bigr\}ds .
\]
\end{proposition}

\begin{proof}[Proof sketch]
Take $k=T/(2\pi)$ in Proposition~\ref{prop:prawitz-prob}, expand
$S(k(x-y))$ using Definition~\ref{def:H} and the $\sinc^2$ representation, and
evaluate the resulting sine and cosine expectations as real parts of
$\hat\nu$. The two kernels differ only in the sign of the real part, which is
exactly the effect of reflecting $x-y\mapsto y-x$.
\end{proof}

\begin{leanref}
\lean{prawitzIntegrand}, \lean{integral\_prawitzIntegrand\_prod},
\lean{integral\_integral\_prawitzIntegrand\_swap},
\lean{cdf\_le\_prawitz\_doubleIntegral},
\lean{prawitz\_doubleIntegral\_le\_cdf} in
\file{BerryEsseen/PrawitzFubini.lean};
\lean{prawitzKernel}, \lean{cdf\_le\_prawitzFourierIntegral},
\lean{prawitzFourierIntegral\_le\_cdf} in
\file{BerryEsseen/PrawitzFourier.lean}.
\end{leanref}

\subsection{The normal distribution function as one Fourier integral}

\begin{lemma}\label{lem:normal-fourier}
For all $x\in\R$,
\[
  \Phi(x)-\frac12 = \frac1\pi\int_0^\infty e^{-t^2/2}\,\frac{\sin(xt)}{t}\,dt
  = \frac{T}{\pi}\int_0^\infty e^{-T^2s^2/2}\,\frac{\sin(xTs)}{Ts}\,ds
  \quad (T>0).
\]
\end{lemma}

\begin{proof}
Let $\gamma$ be the standard Gaussian measure. Fubini's theorem on the finite
rectangle $[0,x]\times\R$, legitimate because $|\cos(ut)|\le1$, gives
\[
  \int_\R\frac{\sin(xt)}{t}\,d\gamma(t)
  = \int_0^x\Bigl(\int_\R\cos(ut)\,d\gamma(t)\Bigr)du
  = \int_0^x e^{-u^2/2}\,du ,
\]
the frequency $t=0$ being a $\gamma$-null set. Substituting the Gaussian
density and using evenness gives the first identity; the integral converges
absolutely since $|\sin(xt)/t|\le|x|$. The substitution $t=Ts$ gives the
second.
\end{proof}

\begin{leanref}
\lean{integral\_cos\_standardGaussian},
\lean{integral\_sine\_quotient\_standardGaussian},
\lean{normalCDF\_fourier\_full}, \lean{normalCDF\_fourier\_positive},
\lean{normalCDF\_fourier\_scaled} in
\file{BerryEsseen/GaussianInversion.lean}.
\end{leanref}

\subsection{The four-term smoothing inequality}

\begin{theorem}[Prawitz smoothing]\label{thm:smoothing}
Let $\nu$ be a probability measure on $\R$ with finite first moment and
characteristic function $\hat\nu$, let $T>0$ and $0<\tau\le1$, and write
$g_T(s)=e^{-T^2s^2/2}$. Then for every $x\in\R$,
\[
  \boxed{
  \begin{aligned}
  \bigl|F(x)-\Phi(x)\bigr| \le{}&
    \int_0^\tau |K(s)|\,\bigl|\hat\nu(Ts)-g_T(s)\bigr|\,ds
    &&\text{(low-frequency error)}\\
  &+ \int_\tau^1 |K(s)|\,\bigl|\hat\nu(Ts)\bigr|\,ds
    &&\text{(high-frequency magnitude)}\\
  &+ \int_0^\tau \Bigl|K(s)-\frac{i}{\pi s}\Bigr|\,g_T(s)\,ds
    &&\text{(normal correction)}\\
  &+ \frac1\pi\int_\tau^\infty \frac{g_T(s)}{s}\,ds
    &&\text{(Gaussian tail).}
  \end{aligned}}
\]
\end{theorem}

\begin{proof}[Proof sketch]
Subtract Lemma~\ref{lem:normal-fourier} from
Proposition~\ref{prop:prawitz-fourier}. On $(0,\tau)$ use the algebraic
identity
\[
  K(s)\hat\nu(Ts) - \frac{i}{\pi s}g_T(s)
  = K(s)\bigl(\hat\nu(Ts)-g_T(s)\bigr)
    + \Bigl(K(s)-\frac{i}{\pi s}\Bigr)g_T(s),
\]
which separates the distributional error from the correction to the normal
inversion kernel; multiplication by $e^{-iTsx}$ preserves norms, and a real
part is bounded by a modulus. On $(\tau,1)$ bound the distributional term by
$|K(s)||\hat\nu(Ts)|$. The remaining normal integral over $[\tau,\infty)$ is
bounded using $|\sin(Tsx)|\le1$. The reflected kernel gives the same two norms
because $K_{-1}=-\overline{K_{+1}}$ and
$K_{-1}-\tfrac{i}{\pi s}=-\overline{K_{+1}-\tfrac{i}{\pi s}}$, so both
inequalities of Proposition~\ref{prop:prawitz-fourier} yield the same bound.

The apparent singularity at $s=0$ is removable. Jordan's inequality and
elementary Taylor bounds give
$|\cot(\pi s)-\tfrac1{\pi s}|\le\tfrac{\pi^2s}{3}$ for $0<s\le\tfrac12$, so
$K(s)-\tfrac{i}{\pi s}$ stays bounded and $s|K(s)|$ stays bounded; and a
finite first moment gives
\[
  \bigl|\hat\nu(Ts)-g_T(s)\bigr|
  \le \Bigl(|T|\textstyle\int|y|d\nu + \frac{T^2}{2}\Bigr)s
  \qquad (0\le s\le1),
\]
which cancels the $1/s$ of the kernel. Away from zero the kernel norms are
bounded and the last integral is dominated by a Gaussian divided by $\pi\tau$.
\end{proof}

\begin{leanref}
\lean{BerryEsseen.smoothingIntegral} and
\lean{BerryEsseen.prawitz\_smoothing} in
\file{BerryEsseen/PrawitzSmoothing.lean}; the four integrands are
\lean{prawitzLowError}, \lean{prawitzHighMagnitude},
\lean{prawitzNormalCorrection}, \lean{prawitzNormalTail} in
\file{BerryEsseen/SmoothingIntegrability.lean}, and the splitting lemma is in
\file{BerryEsseen/SmoothingSplit.lean}.
\end{leanref}

For use with explicit majorants there is a corollary that substitutes any
integrable upper bounds and requires them only on \emph{open} intervals, so
that endpoint values never matter.

\begin{corollary}\label{cor:majorant-form}
In the situation of Theorem~\ref{thm:smoothing}, let $R,M:\R\to\R$ satisfy
$|\hat\nu(Ts)-g_T(s)|\le R(s)$ for $s\in(0,\tau)$ and
$|\hat\nu(Ts)|\le M(s)$ for $s\in(\tau,1)$, with
$s\mapsto|K(s)|R(s)$ integrable on $(0,\tau)$ and $s\mapsto|K(s)|M(s)$
integrable on $(\tau,1)$. Then $|F(x)-\Phi(x)|$ is at most the same four-term
expression with $|\hat\nu(Ts)-g_T(s)|$ and $|\hat\nu(Ts)|$ replaced by
$R(s)$ and $M(s)$.
\end{corollary}

\begin{leanref}
\lean{prawitz\_smoothing\_of\_majorants} and, specialized to the normalized
i.i.d.\ sum, \lean{normalizedSum\_prawitz\_smoothing} in
\file{BerryEsseen/SmoothingMajorants.lean}; the required first moment of
$\mu_n$ comes from \lean{integrable\_id\_normalizedSumLaw}.
\end{leanref}

\subsection{Exact kernel norms and the function $q$}
\label{sec:kernel-norms}

The two kernel norms appearing in Theorem~\ref{thm:smoothing} have closed forms
in terms of $q(x)=\tfrac1x-\cot x$.

\begin{lemma}\label{lem:kernel-norms}
For $0<s<1$,
\[
  \boxed{\
  \Bigl|K(s)-\frac{i}{\pi s}\Bigr| = (1-s)\sqrt{1+q(\pi s)^2},
  \qquad
  |K(s)| = (1-s)\sqrt{1+q(\pi(1-s))^2}. \ }
\]
Moreover $q$ is nonnegative and increasing on $(0,\pi)$, so $|K|$ is
decreasing on $(0,1)$.
\end{lemma}

\begin{proof}
$\Ima K(s)=(1-s)\cot(\pi s)+\tfrac1\pi$. Using
$\cot(\pi(1-s))=-\cot(\pi s)$,
\[
  (1-s)q(\pi(1-s)) = \frac1\pi+(1-s)\cot(\pi s) = \Ima K(s),
\]
which with $\Rea K(s)=1-s$ gives the second identity. For the first,
$\Ima\bigl(K(s)-\tfrac{i}{\pi s}\bigr)
=(1-s)\cot(\pi s)+\tfrac1\pi-\tfrac1{\pi s}=-(1-s)q(\pi s)$, as one checks by
expanding $q(\pi s)=\tfrac1{\pi s}-\cot(\pi s)$.
Nonnegativity and monotonicity of $q$ follow from
\[
  q'(x) = \frac1{\sin^2x}-\frac1{x^2}\ \ge\ 0,
  \qquad
  \sin x - x\cos x = \int_0^x t\sin t\,dt \ \ge\ 0 . \qedhere
\]
\end{proof}

Rational two-sided bounds for $q$ are also needed, both for the analytic
regions and as the meaning of the certificate tables. They come from signed
Taylor bounds for sine, integrated against $t$.

\begin{lemma}\label{lem:q-rational}
Put
\begin{align*}
  A_-(x) &= \tfrac13-\tfrac{x^2}{30}+\tfrac{x^4}{840}-\tfrac{x^6}{45360},
  & A_+(x) &= A_-(x)+\tfrac{x^8}{3991680},\\
  D_+(x) &= 1-\tfrac{x^2}{6}+\tfrac{x^4}{120}-\tfrac{x^6}{5040}
            +\tfrac{x^8}{362880},
  & D_-(x) &= D_+(x)-\tfrac{x^{10}}{39916800}.
\end{align*}
Then $x^3A_-(x)\le\sin x-x\cos x\le x^3A_+(x)$ and
$xD_-(x)\le\sin x\le xD_+(x)$, and consequently
\begin{equation}\label{eq:q-sandwich}
  q_-(x) := \frac{xA_-(x)}{D_+(x)} \ \le\ q(x) \ \le\
  \frac{xA_+(x)}{D_-(x)} =: q_+(x),
\end{equation}
the lower comparison requiring $A_-\ge0$ and the upper one $D_->0$. Both
denominators and the lower numerator are strictly positive on $[0,\tfrac\pi2]$.
\end{lemma}

\begin{leanref}
The four polynomials and the real inequalities are in
\file{BerryEsseen/CotangentPolynomials.lean}; the positivity of the
denominators on the closed half-period is in
\file{BerryEsseen/FiniteKernels.lean}; the passage from the saved integer
tables to \eqref{eq:q-sandwich}, including the conversion from rational bounds
on $\pi$ to the table argument, is \lean{qEntry\_lower\_sound} and
\lean{qEntry\_upper\_sound} in \file{BerryEsseen/Numerics/CotangentTable.lean};
the panel inequalities derived from monotonicity of $q$ and of $|K|$ are in
\file{BerryEsseen/CotangentCorrection.lean}.
\end{leanref}

A coarse bound that is convenient in the small-fraction region is: for
$0<x<\pi$ with $1-x^2/6>0$,
\begin{equation}\label{eq:q-coarse}
  0\le q(x)\le\frac{x}{3(1-x^2/6)},
\end{equation}
which follows from $\sin x-x\cos x=\int_0^xt\sin t\,dt\le x^3/3$ and
$\sin x\ge x(1-x^2/6)$.

Finally, the following reflected estimate is used for the high-frequency
endpoint panels of Section~\ref{sec:finite}.

\begin{lemma}\label{lem:q-reflected}
For $0<s<1$,
\begin{equation}\label{eq:q-refl}
  q\bigl(\pi(1-s)\bigr) \ \le\ \frac{1+s}{\pi s},
\end{equation}
and the right-hand side is decreasing in $s>0$.
\end{lemma}

\begin{proof}
Jordan's inequality $\sin t\ge2t/\pi$ on $[0,\tfrac\pi2]$ gives
$\sin x-x\cos x=\int_0^xt\sin t\,dt\ge\tfrac{2x^3}{3\pi}$; dividing by
$x\sin x\le x^2$ yields $q(x)\ge\tfrac{2x}{3\pi}$ on that half-period. For
$s\le\tfrac12$ this gives
$q(\pi s)\ge\tfrac{2s}{3}\ge\tfrac{s}{\pi(1-s)}$ (using $\pi\ge3$), and
substituting into the exact reflection identity
\[
  q\bigl(\pi(1-s)\bigr) = \frac1{\pi s}+\frac1\pi+\frac{s}{\pi(1-s)}-q(\pi s)
\]
gives \eqref{eq:q-refl}. For $s\ge\tfrac12$, monotonicity of $q$ gives
$q(\pi(1-s))\le q(\pi/2)=2/\pi\le(1+s)/(\pi s)$.
\end{proof}

\begin{leanref}
\file{BerryEsseen/FiniteDarboux.lean}.
\end{leanref}

%=====================================================================
\section{Power comparison and parameter compression}
\label{sec:compression}
%=====================================================================

Independence gives
\[
  \hat\mu_n(t) = f\bigl(t/\sqrt n\bigr)^n ,
\]
so the low-frequency term of Theorem~\ref{thm:smoothing} requires an estimate
of $|z^n-w^n|$ where $z=f(t/\sqrt n)$ and $w=e^{-t^2/(2n)}$. This section
explains how a single inequality is made valid for infinitely many $n$ at once,
which is what allows $467$ cells to cover all $n\ge20$.

\begin{leanref}
\lean{charFun\_normalizedSumLaw} in \file{BerryEsseen/IndependentSums.lean}
proves $\hat\mu_n(t)=f(t/\sqrt n)^n$ for the product-measure definition of
$\mu_n$, and \lean{gaussian\_normalized\_power} in
\file{BerryEsseen/SumBounds.lean} proves
$\bigl(e^{-(t/\sqrt n)^2/2}\bigr)^n=e^{-t^2/2}$.
\end{leanref}

\subsection{Keeping the damping in the geometric sum}

The elementary factorization is
\begin{equation}\label{eq:power-diff}
  |z^n-w^n| \ \le\ |z-w|\sum_{j=0}^{n-1}|z|^{n-1-j}\,|w|^{j} .
\end{equation}
The two bases decay at different rates, and bounding every term of the sum by
its largest value would throw away the extra Gaussian decay of $w$. The
following device retains it.

\begin{definition}\label{def:geom}
Put
\[
  A(x) := \int_0^1 e^{-xs}\,ds,
  \qquad
  G_q(x) := \frac{A(x)}{A(qx)} .
\]
\end{definition}

$A$ is positive, decreasing, continuous, and for $x\ne0$ equals
$(1-e^{-x})/x$; hence for $q,x\ne0$
\[
  G_q(x) = \frac{q\,(1-e^{-x})}{1-e^{-qx}},
  \qquad G_q(0)=1 .
\]
The definition through $A$ avoids the removable singularity at $x=0$, which
matters because the certificates evaluate $G$ at rational arguments that may
be zero.

\begin{lemma}\label{lem:geom-monotone}
For $x\ge0$: $q\mapsto G_q(x)$ is nondecreasing; $0<G_q(x)\le1$ when
$0<q\le1$; and for fixed $0<q\le1$, $x\mapsto G_q(x)$ is nonincreasing on
$[0,\infty)$.
\end{lemma}

\begin{proof}
The first two assertions are immediate from positivity and monotonicity of $A$.
For the third, differentiate the quotient formula: the sign of the derivative is
the sign of
\[
  e^{-x}\bigl(1-e^{-qx}\bigr) - q\,e^{-qx}\bigl(1-e^{-x}\bigr),
\]
and multiplying by $e^{(1+q)x}>0$ reduces this to
$e^{qx}-1-q(e^x-1)\le0$, which is exactly the convexity inequality
$e^{qx}\le1-q+q\,e^x$ for $0\le q\le1$. Continuity extends the conclusion to
$x=0$.
\end{proof}

\begin{leanref}
\lean{expAverage}, \lean{geometricDamping}, \lean{expAverage\_formula},
\lean{geometricDamping\_le\_one}, \lean{geometricDamping\_mono\_q},
\lean{geometricDamping\_deriv\_nonpos},
\lean{antitoneOn\_geometricDamping},
\lean{geometricDamping\_geometric\_sum} in
\file{BerryEsseen/GeometricDamping.lean}.
\end{leanref}

\begin{proposition}\label{prop:compression}
Let $\gamma=1/n$, let $d>0$ satisfy $(\beta+1)/\sqrt n\le d$, and put
\[
  \Psi(t) := t^2h(d|t|), \qquad D(t) := \frac{t^2}{2}-\Psi(t)\ \ge0 .
\]
Then $|z|\le e^{-\gamma\Psi(t)}$ and $w=e^{-\gamma(\Psi(t)+D(t))}$, so
$|z^n|\le e^{-\Psi(t)}$ at every frequency. If moreover
$|z-w|\le\gamma E(t)$, then
\begin{equation}\label{eq:compress-raw}
  |z^n-w^n| \ \le\ E(t)\,e^{-(1-\gamma)\Psi(t)}\,G_\gamma\bigl(D(t)\bigr),
\end{equation}
and consequently, for any cap $\gamma\le v\le1$,
\begin{equation}\label{eq:compress}
  \boxed{\ |z^n-w^n| \ \le\ E(t)\,e^{-(1-v)\Psi(t)}\,G_v\bigl(D(t)\bigr).\ }
\end{equation}
\end{proposition}

\begin{proof}
The modulus bound is Theorem~\ref{thm:modulus} rescaled, using that $h$ is
nonincreasing and $(\beta+1)|t|/\sqrt n\le d|t|$; the nonnegativity
$D\ge0$ is $h\le\tfrac12$. Substituting the two modulus bounds into
\eqref{eq:power-diff},
\[
  \sum_{j=0}^{n-1}|z|^{n-1-j}w^j
  \le e^{-(1-\gamma)\Psi(t)}\sum_{j=0}^{n-1}e^{-j\gamma D(t)}
  = n\,e^{-(1-\gamma)\Psi(t)}\,G_\gamma\bigl(D(t)\bigr),
\]
the last equality being the exact evaluation of the geometric sum in terms of
$A$. Combining with $|z-w|\le\gamma E(t)=E(t)/n$ gives
\eqref{eq:compress-raw}. Since $\Psi\ge0$, replacing $\gamma$ by the larger $v$
increases both $e^{-(1-\gamma)\Psi}$ and, by Lemma~\ref{lem:geom-monotone},
$G_\gamma(D)$; this gives \eqref{eq:compress}.
\end{proof}

\begin{leanref}
\lean{dampingExponent}, \lean{norm\_charFun\_scaled\_le\_damping},
\lean{norm\_charFun\_normalized\_le\_damping},
\lean{norm\_charFun\_normalized\_sub\_gaussian\_le\_scalar\_cap} in
\file{BerryEsseen/ScalarCompression.lean};
\lean{norm\_pow\_sub\_pow\_le\_damped},
\lean{norm\_pow\_sub\_pow\_le\_damped\_cap} in
\file{BerryEsseen/DampedPower.lean}.
\end{leanref}

For numerical evaluation one further needs an upper bound for $G_v$ at a
rational argument. Integrating the elementary Taylor bounds for the exponential
gives, for $x\ge0$, $q\le1$ and $qx<2$,
\begin{equation}\label{eq:G-poly}
  G_q(x) \le \min\Bigl\{1,\
  \frac{1-\tfrac x2+\tfrac{x^2}{6}-\tfrac{x^3}{24}+\tfrac{x^4}{120}}
       {1-\tfrac{qx}{2}}\Bigr\},
\end{equation}
and $G_q(x)\le1$ whenever $qx\ge2$.

\begin{leanref}
\lean{expAverage\_lower}, \lean{expAverage\_upper},
\lean{geometricDamping\_le\_polynomial},
\lean{geometricDamping\_le\_min} in
\file{BerryEsseen/GeometricDamping.lean}; the integer guard that implies
\eqref{eq:G-poly} for a saved certificate is \lean{geometricGuard} with
\lean{geometricGuard\_sound} in
\file{BerryEsseen/Numerics/GeometricWitness.lean}.
\end{leanref}

\subsection{Two error prefactors}

Two choices of $E$ in Proposition~\ref{prop:compression} are used. Write
$b$ for an upper cap on $\ell=\beta/\sqrt n$ and $v$ for an upper cap on
$1/n$.

\paragraph{Classical prefactor.}
With $u=t/\sqrt n$, Definition~\ref{def:bounds} gives
\[
  n\bigl|f(u)-e^{-u^2/2}\bigr|
  \le \frac{\ell|t|^3}{6}+\frac{t^4}{8n}
  \le E_c(t) := \frac{b|t|^3}{6}+\frac{v\,t^4}{8} .
\]

\paragraph{Vector prefactor.}
On a moment band $1\le\beta\le U\le2$, and under the frequency restriction
$|t|/\sqrt n\le\pi$, set
\[
  \delta_U := 1-\frac1U,
  \qquad
  a_U^2 := \Bigl(1-\frac1U\Bigr)\Bigl(1+\frac{5}{3U}\Bigr).
\]
Both normalized moment coefficients are nondecreasing in $U$ on this band, so
Propositions~\ref{prop:real}, \ref{prop:imag} and
Lemma~\ref{lem:gauss-cos} give
\[
  n\bigl|\Rea(f(u)-e^{-u^2/2})\bigr| \le \frac{b\,\delta_U|t|^3}{6}
    + \frac{v\,t^4}{12},
  \qquad
  n\bigl|\Ima f(u)\bigr| \le \frac{b\,a_U|t|^3}{6},
\]
whence, taking the Euclidean norm,
\[
  E_v(t) := \sqrt{\Bigl(\frac{b\delta_U|t|^3}{6}+\frac{v t^4}{12}\Bigr)^2
    + \Bigl(\frac{b\,a_U|t|^3}{6}\Bigr)^2}.
\]
The benefit near $\beta=1$ is visible: both moment-dependent coordinates vanish
as $U\downarrow1$. Either prefactor is valid on its stated domain, so the
smaller may be used, and the certificates select per panel.

\begin{leanref}
\lean{realMomentRatio}, \lean{imaginaryMomentRatioSq},
\lean{realMomentRatio\_mono}, \lean{imaginaryMomentRatioSq\_mono} in
\file{BerryEsseen/VectorRemainder.lean}; \lean{vectorPrefactor} and
\lean{norm\_charFun\_normalized\_sub\_gaussian\_le\_vector\_cap} in
\file{BerryEsseen/VectorCompression.lean}. The packaged envelopes are
\lean{scalarErrorEnvelope}, \lean{vectorErrorEnvelope},
\lean{modulusEnvelope}, \lean{errorDamping} and
\lean{envelopeSmoothingBound} in
\file{BerryEsseen/CompressedSmoothing.lean}, whose two main theorems
\lean{normalizedSum\_cdf\_le\_scalarSmoothing} and
\lean{normalizedSum\_cdf\_le\_vectorSmoothing} bound the actual Kolmogorov
distance by an integral expression containing no unknown law. That file also
proves the integrability of each envelope against the kernel, which is a
genuine hypothesis of Corollary~\ref{cor:majorant-form} and not something the
numerics may assume. The pointwise minimum of the two envelopes is handled in
\file{BerryEsseen/MixedSmoothing.lean}.
\end{leanref}

\subsection{Choosing the caps}
\label{sec:caps}

It remains to produce, from the cell's parameters, admissible values of $v$ and
$d$.

\begin{lemma}\label{lem:caps}
Let $\ell\le b$.
\begin{enumerate}[label=(\alph*),leftmargin=2.4em]
\item If $\beta\ge B>0$ then $1/\sqrt n\le b/B$, hence $1/n\le(b/B)^2$. Taking
  $B=2$ treats all positive sample sizes in the region $\beta\ge2$.
\item If $n\ge20$ and $L\le\beta$, put
  \[
    \rho := \min\Bigl\{\frac{223606798}{10^9},\ \frac bL\Bigr\},
    \qquad
    v := \min\Bigl\{\frac1{20},\ \Bigl(\frac bL\Bigr)^2\Bigr\},
    \qquad
    d := b+\rho .
  \]
  Then $1/\sqrt n\le\rho$, $1/n\le v$ and $(\beta+1)/\sqrt n\le d$; the
  first of these rests on the rational inequality
  $20\,(223606798/10^9)^2\ge1$. Furthermore the guard $\rho\,T\tau\le\pi$
  implies the frequency restriction $T\tau/\sqrt n\le\pi$ required by the
  vector prefactor.
\end{enumerate}
\end{lemma}

\begin{leanref}
\lean{large\_moment\_caps}, \lean{moment\_band\_caps},
\lean{minorant\_argument\_cap} and
\lean{reciprocal\_sqrt\_le\_of\_sample\_bound} in
\file{BerryEsseen/ParameterCaps.lean}; their use is
\lean{normalizedError\_le\_of\_scalar\_cell} and
\lean{normalizedError\_le\_of\_moment\_cell} in
\file{BerryEsseen/ParameterCellBounds.lean}. Note
$223606798/10^9=111803399/(5\cdot10^8)$; the same rational appears in
Section~\ref{sec:scalar} as the endpoint of each moment band.
\end{leanref}

\subsection{Finite sample ranges: a geometric average}
\label{sec:finite-geom}

For the small sample sizes $1\le n\le19$ a different compression is used, in
which the cutoff is $T_n=c\sqrt n$ so that the single-summand frequency
$T_ns/\sqrt n=cs$ is independent of $n$. The comparison across a range
$N\le n\le H$ then rests on the following identity.

\begin{lemma}\label{lem:geom-average}
For $m,g\in\R$ and $k\ge1$ put $G_k(m,g)=\sum_{j=0}^{k-1}m^jg^{k-1-j}$. Then
\begin{equation}\label{eq:geom-integral}
  \frac{G_k(m,g)}{k} = \int_0^1\bigl((1-v)g+vm\bigr)^{k-1}dv .
\end{equation}
If moreover $0\le m,g\le1$, then for $N\le n\le H$,
\begin{equation}\label{eq:geom-compare}
  G_n(m,g) \ \le\ \frac nN\,G_N(m,g) \ \le\ \frac HN\,G_N(m,g).
\end{equation}
\end{lemma}

\begin{proof}
For $m\ne g$, integrate the derivative of $v\mapsto((1-v)g+vm)^k$ and use
$(m-g)G_k(m,g)=m^k-g^k$; for $m=g$ both sides of
\eqref{eq:geom-integral} equal $m^{k-1}$. No case is excluded, including equal
or zero bases. If $0\le m,g\le1$ the integrand of
\eqref{eq:geom-integral} lies in $[0,1]$ and decreases as $k$ increases, which
gives \eqref{eq:geom-compare}.
\end{proof}

\begin{leanref}
\file{BerryEsseen/FiniteGeometricAverage.lean}.
\end{leanref}

Consequently, with $u=cs$, $g=e^{-u^2/2}$, and $M,E$ as in
Section~\ref{sec:finite-bounds},
\begin{equation}\label{eq:finite-power}
  \bigl|f(u)^n-g^n\bigr| \le \frac HN\,u^3E(u)\,G_N\bigl(M(u),g\bigr),
  \qquad |f(u)|^n \le M(u)^N ,
\end{equation}
and $g^n=e^{-T_n^2s^2/2}$ is exactly the comparison function that
Theorem~\ref{thm:smoothing} needs.

%=====================================================================
\section{Region I: small Lyapunov fractions, analytically}
\label{sec:small}
%=====================================================================

This region uses no numerical optimization and no certificate. It proves a
\emph{regional} constant strictly better than the global target.

\begin{theorem}\label{thm:small}
Let $\mu$ be admissible with $\beta=\beta(\mu)$, let $n\ge1$ and suppose
$\ell=\beta/\sqrt n\le\tfrac1{20}$. Then for every $x\in\R$,
\[
  \boxed{\ \bigl|F_n(x)-\Phi(x)\bigr| \ \le\ \frac{183}{400}\,\frac{\beta}{\sqrt n}
  \qquad\Bigl(\frac{183}{400}=0.4575\Bigr).\ }
\]
Consequently $N(\mu,n,x)\le183/400\le293/625$ on this region.
\end{theorem}

We emphasize that this is a regional statement. It is \emph{not} a bound of
$0.4575$ on $\Cii$: the other three regions are what force the global
constant up to $0.4688$.

\begin{leanref}
\lean{BerryEsseen.normalizedError\_small\_fraction} and
\lean{BerryEsseen.normalizedError\_le\_target\_of\_small\_fraction} in
\file{BerryEsseen/SmallFraction.lean}.
\end{leanref}

\subsection{Parameters}

Put
\[
  r := n^{-1/2},\quad \ell := \beta r,\quad
  \nu := 1+\beta^{-1},\quad
  T := \frac{2\pi}{(\beta+1)r} = \frac{2\pi\sqrt n}{\beta+1},\quad
  \tau := \frac{3}{20},\quad
  \Gamma := \sqrt{(\beta-1)\bigl(2\beta+\tfrac23\bigr)} .
\]
Three consequences are recorded once and used throughout.

\begin{lemma}\label{lem:small-params}
If $\beta\ge1$ and $\ell\le\tfrac1{20}$, then $n\ge400$,
\[
  (\beta+1)rT = 2\pi, \qquad \pi\le\ell T, \qquad T\ge60 .
\]
\end{lemma}

\begin{proof}
From $\beta\ge1$ and $\beta/\sqrt n\le\tfrac1{20}$ we get $\sqrt n\ge20\beta\ge20$,
so $n\ge400$. The first identity is the definition of $T$. Next
$\ell T=\beta rT=2\pi\beta/(\beta+1)\ge\pi$ because $\beta\ge1$. Finally
$r\le\beta r\le\tfrac1{20}$ gives $(\beta+1)r=\beta r+r\le\tfrac1{10}$, hence
$T=2\pi/((\beta+1)r)\ge20\pi\ge60$.
\end{proof}

\begin{leanref}
\lean{small\_fraction\_cutoff\_lower} in
\file{BerryEsseen/SmallFractionBudget.lean}.
\end{leanref}

\subsection{The cubic error coefficient and an ellipse inequality}

\begin{lemma}\label{lem:small-error}
For $0\le u\le\pi$,
\begin{equation}\label{eq:small-error}
  \bigl|f(u)-e^{-u^2/2}\bigr| \ \le\ \frac{\Gamma u^3}{6}+\frac{u^4}{8}.
\end{equation}
\end{lemma}

\begin{proof}
Start from \eqref{eq:err-anchor} with $c=e^{-u^2/2}$. By
Lemma~\ref{lem:gauss-cos} the real coordinate is at most
$\tfrac{(\beta-1)u^3}{6}+\tfrac{u^4}{12}$, and by
Proposition~\ref{prop:imag} the imaginary coordinate is at most
$\tfrac{u^3}{6}\sqrt{(\beta-1)(\beta+\tfrac53)}$. The elementary inequality
$\sqrt{(a+b)^2+c^2}\le\sqrt{a^2+c^2}+b$ ($a,b\ge0$) peels off the quartic term,
and
\[
  \sqrt{\Bigl(\frac{(\beta-1)u^3}{6}\Bigr)^2
    + \Bigl(\frac{u^3}{6}\sqrt{(\beta-1)(\beta+\tfrac53)}\Bigr)^2}
  = \frac{u^3}{6}\sqrt{(\beta-1)^2+(\beta-1)(\beta+\tfrac53)}
  = \frac{\Gamma u^3}{6},
\]
because $(\beta-1)+(\beta+\tfrac53)=2\beta+\tfrac23$. Finally
$u^4/12\le u^4/8$.
\end{proof}

\begin{lemma}[Ellipse relation]\label{lem:ellipse}
For $\beta\ge1$,
\begin{equation}\label{eq:ellipse}
  2\nu^2+3\Bigl(\frac\Gamma\beta\Bigr)^2 = 8,
  \qquad\text{hence}\qquad
  \frac{\nu}{5}+\frac{5}{48}\,\frac\Gamma\beta \ \le\ \frac{7}{16}.
\end{equation}
\end{lemma}

\begin{proof}
Writing $\nu=(\beta+1)/\beta$ and $\Gamma^2=(\beta-1)(2\beta+\tfrac23)$,
\[
  2\frac{(\beta+1)^2}{\beta^2}+3\frac{(\beta-1)(2\beta+\frac23)}{\beta^2}
  = \frac{2\beta^2+4\beta+2+6\beta^2-4\beta-2}{\beta^2} = 8 .
\]
For the second assertion, weighted Cauchy--Schwarz gives
\[
  \Bigl(\frac\nu5+\frac5{48}\frac\Gamma\beta\Bigr)^2
  \le \Bigl(2\nu^2+3\Bigl(\frac\Gamma\beta\Bigr)^2\Bigr)
      \Bigl(\frac1{2\cdot25}+\frac{1}{3}\Bigl(\frac{5}{48}\Bigr)^2\Bigr)
  = 8\Bigl(\frac1{50}+\frac{25}{6912}\Bigr)
  = 0.18893\ldots < \Bigl(\frac7{16}\Bigr)^2 = 0.19140625 . \qedhere
\]
\end{proof}

\begin{leanref}
\lean{smallCubicCoefficient}, \lean{sqrt\_shift\_square\_le},
\lean{errorBound\_le\_smallCubic},
\lean{smallCubicCoefficient\_ellipse},
\lean{small\_fraction\_leading\_coefficient} in
\file{BerryEsseen/SmallFractionCoefficients.lean}.
\end{leanref}

\subsection{Low frequencies}

For $0\le s\le\tau$ put $u=Ts/\sqrt n=Tsr$; then $(\beta+1)u=2\pi s\le3\pi/10$,
which is below the splice $9/2$, so the convex minorant is on its linear branch
$h(v)=\tfrac12-\tfrac v{10}$.

Apply Proposition~\ref{prop:compression} in the following concrete form. Let
$m=\exp[-u^2h((\beta+1)u)]$; then $|f(u)|\le m$ by
Theorem~\ref{thm:modulus}, and also $e^{-u^2/2}\le m$ because $h\le\tfrac12$.
Hence by \eqref{eq:power-diff} the geometric sum is at most $nm^{n-1}$, and
\[
  (n-1)\,u^2h\bigl((\beta+1)u\bigr) \ \ge\ \frac{2}{5}\,n\,u^2
  \qquad\text{i.e.}\qquad
  (n-1)h\bigl((\beta+1)u\bigr) \ \ge\ \frac25 n ,
\]
because $(\beta+1)u\le3\pi/10\le\tfrac{19}{20}$ gives
$h\ge\tfrac12-\tfrac{19}{200}=\tfrac{81}{200}$, and $81(n-1)\ge80n$ for
$n\ge81$, which holds since $n\ge400$. Since $nu^2=T^2s^2$, this yields
\begin{equation}\label{eq:small-power}
  \bigl|f(u)^n-e^{-T^2s^2/2}\bigr|
  \ \le\ n\Bigl(\frac{\Gamma u^3}{6}+\frac{u^4}{8}\Bigr)e^{-2T^2s^2/5} .
\end{equation}

\begin{lemma}[Kernel bound on the low band]\label{lem:small-kernel}
For $0<s\le\tau=\tfrac3{20}$,
\[
  s\,|K(s)| \ \le\ \frac{10}{9\pi}.
\]
\end{lemma}

\begin{proof}
The two coordinates of $sK(s)$ are $s(1-s)$ and
$\tfrac1\pi-s(1-s)q(\pi s)$. The first is increasing on $[0,\tfrac12]$, so it
is at most $\tau(1-\tau)=\tfrac{51}{400}$; the second is nonnegative and at
most $\tfrac1\pi$ because $q\ge0$. Hence
$s^2|K(s)|^2\le(\tfrac{51}{400})^2+\tfrac1{\pi^2}$, and
$(\tfrac{51}{400})^2\le\tfrac{19}{81\pi^2}$ follows from $\pi^2\le10$; adding
$\tfrac1{\pi^2}$ gives $\tfrac{100}{81\pi^2}=(\tfrac{10}{9\pi})^2$.
\end{proof}

\begin{lemma}[Finite Gaussian moments]\label{lem:gauss-moments}
For $b>0$ and $R>0$,
\[
  \int_0^R e^{-bs^2}ds \le \frac{\sqrt{\pi/b}}{2},
  \qquad
  \int_0^R s^2e^{-bs^2}ds \le \frac{\sqrt{\pi/b}}{4b},
  \qquad
  \int_0^R s^3e^{-bs^2}ds \le \frac{1}{2b^2}.
\]
\end{lemma}

\begin{proof}
The first is the half-line Gaussian integral. The other two follow by
integrating the derivatives of $se^{-bs^2}$ and of
$-(bs^2+1)e^{-bs^2}/(2b^2)$ and discarding nonnegative endpoint terms.
\end{proof}

\begin{proposition}\label{prop:small-low}
The low-frequency term satisfies
\[
  L := \int_0^\tau |K(s)|\,\bigl|\hat\mu_n(Ts)-e^{-T^2s^2/2}\bigr|\,ds
  \ \le\ \frac{5}{48}\,\Gamma\,r + \frac{7}{50}\,r^2 .
\]
\end{proposition}

\begin{proof}
Insert \eqref{eq:small-error} into \eqref{eq:small-power}, then pair the factor
$s$ with the kernel using Lemma~\ref{lem:small-kernel}. With $u=Tsr$ and
$nr^2=1$, the cubic part of the integrand becomes
$\tfrac{10}{9\pi}\cdot\tfrac{\Gamma T^3r}{6}\,s^2e^{-bs^2}$ and the quartic part
$\tfrac{10}{9\pi}\cdot\tfrac{T^4r^2}{8}\,s^3e^{-bs^2}$, where $b=2T^2/5$. Apply
Lemma~\ref{lem:gauss-moments}. Since
$\sqrt{\pi/b}=\sqrt{5\pi/2}/T\le 9\pi/(10T)$ --- equivalent to $\pi\ge250/81$
--- we get
\[
  \int_0^\tau s^2e^{-bs^2}ds \le \frac{\sqrt{\pi/b}}{4b}\le\frac{9\pi}{16T^3},
  \qquad
  \int_0^\tau s^3e^{-bs^2}ds \le \frac1{2b^2} = \frac{25}{8T^4},
\]
whence the cubic contribution is at most
$\tfrac{10}{9\pi}\cdot\tfrac{\Gamma T^3r}{6}\cdot\tfrac{9\pi}{16T^3}
=\tfrac{5}{48}\Gamma r$ and the quartic contribution at most
$\tfrac{10}{9\pi}\cdot\tfrac{T^4r^2}{8}\cdot\tfrac{25}{8T^4}
=\tfrac{125}{288\pi}r^2\le\tfrac{7}{50}r^2$.
\end{proof}

\begin{leanref}
\lean{small\_fraction\_low\_integral} in
\file{BerryEsseen/SmallFractionLow.lean}; the three ingredients are in
\file{BerryEsseen/SmallFractionPower.lean} (power comparison),
\file{BerryEsseen/SmallFractionKernels.lean} (Lemma~\ref{lem:small-kernel}),
and \file{BerryEsseen/SmallFractionGaussian.lean}
(Lemma~\ref{lem:gauss-moments}).
\end{leanref}

\subsection{Normal correction}

\begin{proposition}\label{prop:small-normal}
\[
  N := \int_0^\tau\Bigl|K(s)-\frac{i}{\pi s}\Bigr|e^{-T^2s^2/2}\,ds
  \ \le\ \frac{2\pi}{5T}+\frac{\pi^3}{40T^3}.
\]
\end{proposition}

\begin{proof}
By Lemma~\ref{lem:kernel-norms} and \eqref{eq:q-coarse}, on $s\le\tfrac3{20}$
we get $q(\pi s)^2\le\pi^2s^2/8$, hence
\[
  \Bigl|K(s)-\frac{i}{\pi s}\Bigr| = (1-s)\sqrt{1+q(\pi s)^2}
  \le 1+\frac{\pi^2s^2}{16},
\]
using $\sqrt{1+y}\le1+y/2$. Now apply Lemma~\ref{lem:gauss-moments} with
$b=T^2/2$, for which $\sqrt{\pi/b}=\sqrt{2\pi}/T\le4\pi/(5T)$ (equivalent to
$\pi\ge25/8$). The constant term contributes at most $2\pi/(5T)$ and the
quadratic term at most
$\tfrac{\pi^2}{16}\cdot\tfrac{4\pi}{5T}\cdot\tfrac1{2T^2}=\tfrac{\pi^3}{40T^3}$.
\end{proof}

\begin{leanref}
\file{BerryEsseen/SmallFractionNormal.lean}.
\end{leanref}

\subsection{High frequencies and the tail}

Split the high band at $R=9/(4\pi)$.

\begin{lemma}\label{lem:small-decay}
For $\tau\le s\le R$ one has $s^2h(2\pi s)\ge\tfrac9{1000}$, and for
$R\le s\le1$, writing $y=1-s$, one has $s^2h(2\pi s)\ge\tfrac3{10}y^2$.
\end{lemma}

\begin{proof}[Proof sketch]
Writing $J(v)=v^2h(v)$ one has $s^2h(2\pi s)=J(2\pi s)/(4\pi^2)$. On
$[\tau,R]$ the argument $2\pi s$ ranges over $[3\pi/10,\,9/2]$, where $J$ is the
cubic branch $v^2/2-v^3/10$; this increases up to $v=10/3$ and decreases
afterwards, so its minimum on the interval is at an endpoint. Evaluating at
$v=3\pi/10$ and $v=9/2$ gives $J\ge0.3604\ldots$, and dividing by $4\pi^2$
gives $\ge9/1000$. On $[R,1]$ the argument lies in $[9/2,2\pi]$, where $h$ is
the cosine branch $aG$ with $a>\tfrac56$; using $\sin(\pi y)\ge\tfrac{17}{20}\pi y$
one obtains $s^2h(2\pi s)\ge\tfrac3{10}y^2$.
\end{proof}

\begin{proposition}\label{prop:small-high}
With $H_1$ and $H_2$ the parts of the high-frequency term over $[\tau,R]$ and
$[R,1]$,
\[
  H_1 \le 100\,e^{-9T^2/1000},
  \qquad
  H_2 \le \frac{50}{27T^2},
\]
and moreover, for $T\ge60$,
\[
  100\,e^{-9T^2/1000} \le \frac{9}{100T^2},
  \qquad
  G := \frac1\pi\int_\tau^\infty\frac{e^{-T^2s^2/2}}{s}\,ds
  \le \frac{9}{100T^2}.
\]
\end{proposition}

\begin{proof}
By Theorem~\ref{thm:modulus}, $|\hat\mu_n(Ts)|\le\exp(-T^2s^2h(2\pi s))$,
since $(\beta+1)Ts/\sqrt n=2\pi s$. On $[\tau,R]$ the coarse kernel bound
$|K(s)|\le(4+\pi^2)/s\le100$ (valid for $s\ge\tfrac3{20}$) and
Lemma~\ref{lem:small-decay} give $H_1\le100e^{-9T^2/1000}$. On $[R,1]$,
Lemma~\ref{lem:kernel-norms} together with the same cotangent estimate gives
$|K(s)|\le\tfrac{10}{9}y$, so
\[
  H_2 \le \frac{10}{9}\int_0^{1-R}y\,e^{-3T^2y^2/10}\,dy
  \le \frac{10}{9}\cdot\frac{10}{6T^2} = \frac{50}{27T^2}.
\]
For the first remainder bound, put $x=9T^2/1000\ge32$ (using $T\ge60$). From
$e^x\ge x^7/7!$ we get $xe^{-x}\le 7!/x^6\le5040/32^6\le81/10^7$, hence
$100e^{-x}=100\,(xe^{-x})/x\le100\cdot\tfrac{81}{10^7}\cdot\tfrac{1000}{9T^2}
=\tfrac9{100T^2}$. For the tail, $1/s\le s/\tau^2$ on $[\tau,\infty)$ gives
$G\le e^{-T^2\tau^2/2}/(\pi T^2\tau^2)$; here $T^2\tau^2/2\ge40$, and
$e^{-x}\le\tfrac1{400}$ for $x\ge40$ (from $e^x\ge x^2/2$), so with $\pi\ge3$,
$G\le\tfrac1{400}\cdot\tfrac{400}{9\pi T^2}=\tfrac1{9\pi T^2}\le\tfrac9{100T^2}$.
\end{proof}

\begin{leanref}
\file{BerryEsseen/SmallFractionDecay.lean} (Lemma~\ref{lem:small-decay}),
\file{BerryEsseen/SmallFractionHigh.lean} ($H_1,H_2$),
\file{BerryEsseen/SmallFractionRemainders.lean} (the two exponential
remainders).
\end{leanref}

\subsection{The rational budget}

\begin{proof}[Proof of Theorem~\ref{thm:small}]
Theorem~\ref{thm:smoothing} applied to $\mu_n$ with cutoff $T$ and split
$\tau=\tfrac3{20}$ gives $|F_n(x)-\Phi(x)|\le L+H_1+H_2+N+G$. Collect the
first-order terms: because $(\beta+1)rT=2\pi$ we have
$\tfrac{2\pi}{5T}=\tfrac{(\beta+1)r}{5}=\tfrac\nu5\ell$, so by
Lemma~\ref{lem:ellipse}
\[
  \frac5{48}\Gamma r+\frac{2\pi}{5T}
  = \Bigl(\frac5{48}\frac\Gamma\beta+\frac\nu5\Bigr)\ell
  \ \le\ \frac7{16}\,\ell .
\]
For the remaining terms use $r\le\ell$, $\pi\le\ell T$ and
$\ell\le\tfrac1{20}$:
\[
  \frac7{50}r^2\le\frac7{50}\ell^2,
  \qquad
  \frac{\pi^3}{40T^3}\le\frac{\ell^3}{40}\le\frac{\ell^2}{100},
  \qquad
  \frac{18}{100T^2}+\frac{50}{27T^2}\le\frac{\ell^2}{4},
\]
the last because $1/T^2\le\ell^2/\pi^2\le\ell^2/9$ and
$(0.18+50/27)/9\le\tfrac14$. The three coefficients sum to
$\tfrac7{50}+\tfrac1{100}+\tfrac14=\tfrac25$. Therefore
\[
  \bigl|F_n(x)-\Phi(x)\bigr| \le \frac7{16}\ell+\frac25\ell^2
  \le \Bigl(\frac7{16}+\frac25\cdot\frac1{20}\Bigr)\ell
  = \Bigl(\frac7{16}+\frac1{50}\Bigr)\ell = \frac{183}{400}\,\ell .
\]
Lemma~\ref{lem:cdf-to-normalized} with $a=\ell$ and $C=183/400$ finishes the
proof, and $183/400<293/625$.
\end{proof}

\begin{leanref}
\lean{small\_fraction\_total\_budget} in
\file{BerryEsseen/SmallFractionBudget.lean} is exactly the displayed rational
inequality, and \lean{normalizedError\_small\_fraction} in
\file{BerryEsseen/SmallFraction.lean} assembles the five integral bounds with
it.
\end{leanref}

%=====================================================================
\section{Region II: large Lyapunov fractions, by Cantelli}
\label{sec:large}
%=====================================================================

When $\ell$ is large the normalization $\sqrt n/\beta=1/\ell$ is small and no
Fourier calculation is needed: a universal bound of $\tfrac{11}{20}$ on the
Kolmogorov distance between \emph{any} centred, variance-one law and the normal
law suffices.

\begin{theorem}\label{thm:cantelli}
Let $Y$ have mean zero and variance one, and let $F(x)=\PR(Y\le x)$. Then
\[
  \boxed{\ \sup_{x\in\R}\bigl|F(x)-\Phi(x)\bigr| \ \le\ \frac{11}{20}. \ }
\]
Laws with atoms are included.
\end{theorem}

\begin{proof}
\emph{One-sided Chebyshev.} For $a>0$, Markov's inequality applied to
$(Y+1/a)^2$ gives
\[
  \PR(Y\ge a) \le \frac{\E(Y+1/a)^2}{(a+1/a)^2}
  = \frac{1+a^{-2}}{(a+a^{-1})^2} = \frac1{1+a^2},
\]
using $\E(Y+1/a)^2=1+a^{-2}$ (mean zero, variance one) and
$\{Y\ge a\}\subseteq\{(Y+1/a)^2\ge(a+1/a)^2\}$. Applying the same to $-Y$
bounds $\PR(Y\le-a)$.

\emph{Upper side.} For $x>0$ this gives $F(x)\ge1-\tfrac1{1+x^2}
=\tfrac{x^2}{1+x^2}$, hence
$\Phi(x)-F(x)\le\Phi(x)-\tfrac{x^2}{1+x^2}$. If $x\ge1$ the right side is at
most $1-\tfrac12=\tfrac12$. If $0\le x\le1$, then $\varphi\le\tfrac25$ gives
$\Phi(x)\le\tfrac12+\tfrac{2x}{5}$ and it remains to check
\[
  \frac12+\frac{2x}{5}-\frac{x^2}{1+x^2}\ \le\ \frac{11}{20}
  \qquad (0\le x\le1).
\]
Multiplying by $20(1+x^2)>0$, this is equivalent to $P(x)\ge0$ where
\[
  P(x) := 1-8x+21x^2-8x^3 .
\]
On $[0,\tfrac12]$ we have $21x^2-8x^3\ge17x^2$, so
$P(x)\ge1-8x+17x^2=17(x-\tfrac4{17})^2+\tfrac1{17}>0$. On $[\tfrac12,1]$ we
have $21x^2-8x^3\ge13x^2$, so
$P(x)\ge1-8x+13x^2=13(x-\tfrac12)^2+5(x-\tfrac12)+\tfrac14>0$.

\emph{Lower side and negative thresholds.} $F(x)-\Phi(x)\le1-\Phi(x)\le\tfrac12$
for $x\ge0$. For $x<0$ use the lower-tail bound together with
$\Phi(-x)=1-\Phi(x)$. At $x=0$ use $0\le F(0)\le1$ and $\Phi(0)=\tfrac12$.
\end{proof}

\begin{leanref}
\lean{cantelli\_upper}, \lean{cantelli\_lower},
\lean{cdf\_error\_le\_eleven\_twentieths},
\lean{normalizedSum\_cdf\_error\_le\_eleven\_twentieths} in
\file{BerryEsseen/Cantelli.lean}; the normal density and distribution function
facts ($\varphi\le\tfrac25$, monotonicity, symmetry, and the polynomial gap
lemma) are in \file{BerryEsseen/GaussianBounds.lean}; the moment identities for
the normalized sum ($\E S_n=0$, $\operatorname{Var}S_n=1$, $S_n\in L^2$) are in
\file{BerryEsseen/IndependentSums.lean}.
\end{leanref}

\begin{corollary}\label{cor:large}
If $\mu$ is admissible, $n\ge1$ and $\ell=\beta/\sqrt n\ge\tfrac65$, then for
every $x$,
\[
  N(\mu,n,x) \ \le\ \frac{5}{6}\cdot\frac{11}{20} = \frac{11}{24}
  = 0.4583\overline{3} \ <\ \frac{293}{625}.
\]
\end{corollary}

\begin{proof}
$\sqrt n/\beta=1/\ell\le\tfrac56$; multiply Theorem~\ref{thm:cantelli}
by it.
\end{proof}

\begin{leanref}
\lean{normalizedError\_large\_fraction} and
\lean{normalizedError\_le\_target\_of\_large\_fraction} in
\file{BerryEsseen/LargeFraction.lean}; the same file contains the conversion
\lean{upperTarget\_eq\_ofReal} proving
$(0.4688:\lean{ENNReal})=\lean{ENNReal.ofReal}(293/625)$, which every region
uses at the end.
\end{leanref}

%=====================================================================
\section{Region III: 467 certified interval cells}
\label{sec:scalar}
%=====================================================================

This region covers
\[
  \ell\ge\tfrac1{20}
  \quad\text{together with}\quad
  \bigl(\beta\ge2,\ n\ge1\bigr)
  \ \text{ or }\
  \bigl(n\ge20,\ \beta\ \text{arbitrary}\bigr),
\]
that is, everything except small fractions and small sample sizes with small
moments.

\begin{theorem}\label{thm:scalar}
Let $\mu$ be admissible and $x\in\R$.
\begin{enumerate}[label=(\alph*),leftmargin=2.4em]
\item If $n\ge1$, $\beta\ge2$ and $\ell\ge\tfrac1{20}$, then
  $N(\mu,n,x)\le\tfrac{293}{625}$.
\item If $n\ge20$ and $\ell\ge\tfrac1{20}$, then
  $N(\mu,n,x)\le\tfrac{293}{625}$.
\end{enumerate}
\end{theorem}

\begin{leanref}
\lean{ScalarFullHAll.large\_moment\_normalizedError} in
\file{.../ScalarFullH/LargeMomentCover.lean} and
\lean{ScalarFullHAll.large\_sample\_normalizedError} in
\file{.../ScalarFullH/MomentBands.lean}.
\end{leanref}

\subsection{What a cell is}

Fix a closed Lyapunov-fraction interval $[a,b]\subseteq[\tfrac1{20},\tfrac65]$,
a third-moment floor $L\ge1$, a sample-size floor $N_{\min}\ge1$, a cutoff
$T>0$, a split $0<\tau<\tfrac12$ and a panel count $P\ge1$. From these the
caps $v$ and $d$ of Lemma~\ref{lem:caps} are computed, together with a moment
ceiling $U$ when the cell enables the anchored (``vector'') error prefactor.
Corollary~\ref{cor:majorant-form} and Proposition~\ref{prop:compression} then
bound the Kolmogorov distance of \emph{every} admissible law in the cell by
\begin{equation}\label{eq:cell-functional}
  \mathcal B :=
  \int_0^\tau |K(s)|\,\mathcal E(Ts)\,ds
  + \int_\tau^1 |K(s)|\,e^{-\Psi(Ts)}\,ds
  + \int_0^\tau \Bigl|K(s)-\frac{i}{\pi s}\Bigr| e^{-T^2s^2/2}ds
  + \frac1\pi\int_\tau^\infty \frac{e^{-T^2s^2/2}}{s}\,ds ,
\end{equation}
where $\Psi(t)=t^2h(d|t|)$ and the error envelope is
\[
  \mathcal E(t) := \mathcal P(t)\,
  e^{-(1-v)\Psi(t)}\,G_v\bigl(t^2/2-\Psi(t)\bigr),
\]
$\mathcal P$ being the scalar prefactor $\tfrac{b|t|^3}{6}+\tfrac{v t^4}{8}$,
the vector prefactor $E_v$ of Section~\ref{sec:compression}, or their pointwise
minimum. Note that \eqref{eq:cell-functional} contains no unknown probability
law: only the rational cell parameters.

\begin{leanref}
\lean{ScalarFullH.GeneralCellData} and its Boolean validator
\lean{generalCellData} in
\file{BerryEsseen/Numerics/GeneralCellCore.lean}, built on
\lean{FixedCellData}/\lean{fixedCellData} in
\file{BerryEsseen/Numerics/FixedCellCore.lean}. A \lean{GeneralCell} is a pair
of a data record and a proof that its validator returns \lean{true}. The
executable arithmetic of these validators was copied unchanged from the saved
certificate core; only module and import wrappers were changed.
\end{leanref}

The validator checks, among other things,
\[
  \tfrac1{20}\le a<b\le\tfrac65,\quad 1\le L,\quad 0<N_{\min},\quad
  1\le N_{\min}\rho^2,\quad 0<\tau<\tfrac12,\quad
  d\,T\tau\le\tfrac92,\quad 0<v<1,
\]
that the integer grid step and denominator represent exactly $\tau/P$, and that
each stored fixed-point integer really dominates (or is dominated by) the
corresponding real quantity: e.g.\ $bT^3/6\le\lean{cubicError}/2^{64}$,
$vT^4/8\le\lean{quarticError}/2^{64}$,
$\lean{cubicDecayLo}/2^{64}\le dT^3/10\le\lean{cubicDecayHi}/2^{64}$,
$\lean{quadraticDecayLo}/2^{64}\le T^2/2$, $v\le\lean{varianceHi}/2^{64}$. When
the anchor is enabled it additionally checks $L\le U\le2$,
$\rho T\tau\le\pi_{\mathrm{low}}$, and three fixed-point bounds for
$b\delta_UT^3/6$, $vT^4/12$ and $b^2a_U^2T^6/36$. Here
$\pi_{\mathrm{low}}=3141592653589793238/10^{18}\le\pi$.

\begin{leanref}
The real consequences of these Boolean checks are proved in
\file{BerryEsseen/Numerics/CellBasicFacts.lean} (cutoff and split
inequalities, grid positivity, the exact relation between the rational split and
the integer grid, the Gaussian quadratic coefficient),
\file{.../CellDecayFacts.lean} (positivity of $d$, its relation to the stored
argument scale, cubic and cosine coefficient bounds),
\file{.../CellErrorFacts.lean} (variance cap, cubic and quartic scalar error
coefficients) and \file{.../CellAnchorFacts.lean} (moment range, frequency
restriction, separate real and imaginary coefficients). These are derived from
the record checks, not supplied as extra assumptions by the generator.
\end{leanref}

\subsection{Panel bounds: how endpoint data control whole integrals}
\label{sec:panel-bounds}

Each of the three finite integrals in \eqref{eq:cell-functional} is split into
$P$ panels $[l,r]$ with $r-l=\tau/P$ (low and normal) or into a grid on
$[\tau,1]$ (high), and each panel is bounded using only data at its endpoints.
The mechanism is the following monotonicity structure.

\begin{lemma}[The decay profile has one maximum]\label{lem:decay-profile}
Let $J(x)=x^2h(x)$, i.e.
\[
  J(x) =
  \begin{cases}
    \tfrac{x^2}{2}-\tfrac{x^3}{10}, & 0\le x\le\tfrac92,\\
    a\,(1-\cos x), & \tfrac92\le x\le2\pi,\\
    0, & x\ge2\pi.
  \end{cases}
\]
Then $J$ increases on $[0,\tfrac{10}{3}]$ and decreases on
$[\tfrac{10}{3},\infty)$, and the branches agree at both joins. Consequently,
since $\Psi(t)=J(dt)/d^2$, for $0\le l\le s\le r$,
\[
  \Psi(Ts) \ \ge\ \min\{\Psi(Tl),\Psi(Tr)\} ,
\]
including for panels crossing either splice or extending into the zero branch.
\end{lemma}

\begin{leanref}
\file{BerryEsseen/DecayPanels.lean}; the three admissible endpoint witness
modes --- the value $0$, the cubic formula, and a cosine-table entry --- are
justified in \file{BerryEsseen/DecayWitnesses.lean}. In the cosine mode the
table argument is rounded \emph{upwards} inside $[\tfrac92,2\pi]$, which
increases the cosine and therefore decreases the certified decay, i.e.\ errs
safely.
\end{leanref}

\begin{proposition}[High-frequency and normal-correction panels]
\label{prop:high-normal-panel}
Let $0<l\le r<1$.
\begin{enumerate}[label=(\alph*),leftmargin=2.4em]
\item If $\Lambda\le\Psi(Tl)$ and $\mathrm{R}\le\Psi(Tr)$ then
  \[
    \int_l^r |K(s)|\,e^{-\Psi(Ts)}\,ds
    \ \le\ (r-l)\,|K(l)|\,e^{-\min(\Lambda,\mathrm R)} ,
  \]
  using that $|K|$ is decreasing (Lemma~\ref{lem:kernel-norms}).
\item If $q(\pi r)\le Q$ then
  \[
    \int_l^r \Bigl|K(s)-\frac{i}{\pi s}\Bigr|\,e^{-T^2s^2/2}\,ds
    \ \le\ (r-l)(1-l)\sqrt{1+Q^2}\;e^{-T^2l^2/2} ,
  \]
  using that $q$ is nonnegative and increasing. The case $l=0$ is included; no
  value at the singular endpoint is used.
\end{enumerate}
\end{proposition}

\begin{leanref}
\file{BerryEsseen/SmoothingPanels.lean}.
\end{leanref}

\begin{proposition}[Low-frequency panels]\label{prop:low-panel}
Let $0\le l\le s\le r\le\tfrac12$, let $0<p\le\pi$ and $0\le Q\le q(\pi l)$.
Then
\begin{equation}\label{eq:Kstar}
  s|K(s)| \ \le\ K_* := \sqrt{[r(1-r)]^2+\bigl[\tfrac1p-l(1-l)Q\bigr]^2} .
\end{equation}
If in addition $dTr\le\tfrac92$ --- so that the panel stays in the cubic branch,
where $(Ts)^2/2-\Psi(Ts)=dT^3s^3/10$ exactly --- then for any
$0\le\eta\le dT^3l^3/10$, any $0<v\le1$, and $\Lambda,\mathrm R$ as above,
\[
  \int_l^r |K(s)|\,\mathcal E(Ts)\,ds
  \ \le\ (r-l)\,K_*\,\mathcal P^\sharp(r)\,
    e^{-(1-v)\min(\Lambda,\mathrm R)}\,G_v(\eta),
\]
where $\mathcal P^\sharp$ is the increasing majorant of $s\mapsto s\mathcal P(Ts)$
on the panel:
\[
  \mathcal P^\sharp_{\mathrm{scalar}}(s) = \frac{bT^3s^2}{6}+\frac{vT^4s^3}{8},
  \qquad
  \mathcal P^\sharp_{\mathrm{vector}}(s) = s^2\sqrt{
    \Bigl(\frac{b\delta_UT^3}{6}+\frac{vT^4s}{12}\Bigr)^2
    + \Bigl(\frac{b a_U T^3}{6}\Bigr)^2 } .
\]
\end{proposition}

\begin{proof}[Proof sketch of \eqref{eq:Kstar}]
The two coordinates of $sK(s)$ are $s(1-s)$, increasing on $[0,\tfrac12]$, and
$\tfrac1\pi-s(1-s)q(\pi s)$, which is nonnegative and, being decreasing in $s$,
is at most its value at $l$ with the certified lower bounds $p\le\pi$ and
$Q\le q(\pi l)$ substituted.
\end{proof}

\begin{leanref}
\file{BerryEsseen/LowSmoothingPanels.lean} (scalar) and
\file{BerryEsseen/VectorSmoothingPanels.lean} (vector);
\lean{PanelSums.lean} proves that adjacent panel integrals telescope to the
complete integral, with the first and last endpoints identified explicitly.
\end{leanref}

\begin{proposition}[Infinite normal tail]\label{prop:tail}
For $\tau>0$ and $P=T^2\tau^2/2$,
\[
  \frac1\pi\int_\tau^\infty\frac{e^{-T^2s^2/2}}{s}\,ds
  \ \le\ \frac{e^{-P}}{2\pi P} ,
\]
using $1/s\le s/\tau^2$ on $[\tau,\infty)$ and the elementary antiderivative of
$se^{-bs^2}$.
\end{proposition}

\begin{leanref}
\file{BerryEsseen/GaussianTailBound.lean}, which also proves integrability and
the form in which certified lower bounds on $P$ and $\pi$ and an upper bound
on the exponential are substituted.
\end{leanref}

\subsection{From a passing record to a real integral}

The saved certificate stores, for each panel, a tuple of natural numbers and
table references, and a Boolean check. For each of the five record kinds the
development proves a theorem of the form ``if the check returns \lean{true},
then the corresponding real integral over the corresponding exact panel is at
most $\lean{termUpper}/2^{64}$''. Table~\ref{tab:soundness} lists them.

\begin{table}[ht]
\centering
\small
\begin{tabular}{@{}lll@{}}
\toprule
Saved check & Real conclusion & Lean theorem \\
\midrule
\lean{generalTail} & the infinite normal tail
  & \lean{generalTail\_sound} \\
\lean{generalNormalPanel} & the normal-correction integral on that panel
  & \lean{generalNormalPanel\_sound} \\
\lean{generalHighPanel} & the high-frequency modulus integral there
  & \lean{generalHighPanel\_sound} \\
\lean{generalLowPanel} (scalar) & the scalar-error integral on that panel
  & \lean{generalLowPanel\_scalar\_sound} \\
\lean{generalLowPanel} (vector) & the vector-error integral on that panel
  & \lean{generalLowPanel\_vector\_sound} \\
\lean{generalLowPanel} (either) & the cell's error-envelope integral
  & \lean{generalLowPanel\_sound} \\
\bottomrule
\end{tabular}
\caption{Record-to-integral soundness theorems; each applies to \emph{every}
passing record, and every table-entry proof and integer-to-real conversion is in
its dependency closure.}
\label{tab:soundness}
\end{table}

Several details are worth recording, because they are exactly the places where a
naive ``the numbers check out'' argument would be insufficient.
\begin{itemize}[leftmargin=2.2em]
\item The normal-panel proof uses the \emph{upper} cotangent table entry at the
  \emph{right} endpoint and the Gaussian exponential at the \emph{left}
  endpoint; the high-panel proof uses both stored kernel formulas, both
  endpoint decay witnesses and a rounded exponential bound.
\item The low-panel proof checks the lower cotangent bound, the scaled kernel
  norm, both decay endpoints, a lower bound for the damping difference
  $dT^3l^3/10$, the exponential and geometric factors, and the final product.
  Natural-number subtraction is converted to real subtraction only after its
  nonnegativity has been proved, and rounding the variance cap \emph{upwards}
  weakens both damping bounds in the safe direction.
\item For anchor-enabled cells the low envelope is the pointwise \emph{minimum}
  of the scalar and vector envelopes; that minimum is proved integrable against
  the kernel and proved to bound the actual characteristic-function error, so
  the panel records may select either branch independently.
\item \lean{generalLowPanels\_sound} requires a passing record at \emph{every}
  index and proves that the grid covers exactly $[0,\tau]$;
  \lean{HighGrid.lean} proves that the high grid starts at $\tau$ and ends at
  $1$; \lean{CellPanelSums.lean} assembles the three complete panel lists and
  the tail into one bound for the four-term expression
  \eqref{eq:cell-functional}.
\end{itemize}

\begin{leanref}
\file{BerryEsseen/Numerics/CellApplication.lean},
\file{.../CellPanelSums.lean}, \file{.../HighGrid.lean},
\file{.../IndexedPanelLists.lean}, \file{.../CellCertificate.lean};
\file{BerryEsseen/MixedSmoothing.lean} for the minimum envelope.
\end{leanref}

\subsection{The budget and the covering}

\begin{definition}\label{def:certified-cell}
A \emph{certified cell} is a \lean{GeneralCell} together with complete lists of
passing low, high and normal panel records indexed by $0,\dots,P-1$, a passing
tail record, and a proof of the rational strict inequality
\begin{equation}\label{eq:cell-budget}
  \frac{S}{2^{64}} \ <\ \frac{293}{625}\cdot a,
  \qquad
  S := \sum_{\text{low}}\lean{termUpper}+\sum_{\text{high}}\lean{termUpper}
       +\sum_{\text{normal}}\lean{termUpper}+\lean{termUpper}_{\text{tail}} ,
\end{equation}
where $a$ is the cell's \emph{lower} Lyapunov-fraction endpoint.
\end{definition}

\begin{theorem}\label{thm:cell-sound}
Let $c$ be a certified cell with parameters $(a,b,L,N_{\min},U)$, and let $\mu$
be admissible with $n\ge N_{\min}$, $L\le\beta$, $a\le\ell\le b$, and
$\beta\le U$ if the cell enables the anchor. Then for every $x\in\R$,
\[
  N(\mu,n,x) \ \le\ \frac{293}{625}.
\]
\end{theorem}

\begin{proof}
By the panel soundness theorems and \eqref{eq:cell-functional},
$|F_n(x)-\Phi(x)|\le\mathcal B\le S/2^{64}$. By \eqref{eq:cell-budget} this is
$<\tfrac{293}{625}a\le\tfrac{293}{625}\ell$, and
Lemma~\ref{lem:cdf-to-normalized} with $C=293/625$ finishes.
\end{proof}

\begin{leanref}
\lean{ScalarFullH.CertifiedCell} with
\lean{CertifiedCell.normalizedError\_le} in
\file{BerryEsseen/Numerics/CertifiedCell.lean}, on top of
\lean{generalCell\_certificate\_normalizedError}. The hypotheses on the law are
exactly its sample-size lower bound, moment range and fraction range; cutoff
validity, the sample-size caps and the anchor's frequency restriction all follow
from the checked cell data rather than being assumed.
\end{leanref}

Finitely many cells cover a continuum because their endpoints chain.

\begin{lemma}[Coverage]\label{lem:coverage}
Let $c_0,\dots,c_{m}$ be cells whose stored adjacency check
$\lean{adjacentCells}$ returns \lean{true}, i.e.\ $b_{j}=a_{j+1}$ for all $j$.
Then every real $\ell$ with $a_0\le\ell\le b_m$ lies in some $[a_j,b_j]$.
\end{lemma}

\begin{proof}
Induct on the list: if $\ell\le b_0$ take $j=0$; otherwise $\ell>b_0=a_1$ and
apply the inductive hypothesis to the tail.
\end{proof}

\begin{leanref}
\lean{ScalarFullHAll.adjacentCells\_cover} and
\lean{certified\_cover\_normalizedError} in
\file{.../ScalarFullH/Coverage.lean}. This proves coverage
of the continuous parameter interval, not merely of the stored endpoints.
\end{leanref}

\subsection{The five coverings}

The saved data consist of $467$ cells, \lean{cell0} through \lean{cell466},
grouped into five chains. Their parameters are summarized in
Table~\ref{tab:scalar-cells}.

\begin{table}[ht]
\centering
\begin{tabular}{@{}llrl@{}}
\toprule
Third-moment range & Sample sizes & Cells & Lean indices \\
\midrule
$\beta\ge2$ & $n\ge1$ & $137$ & \lean{cell0}--\lean{cell136} \\
$1\le\beta\le\tfrac{11}{10}$ & $n\ge20$ & $76$ & \lean{cell137}--\lean{cell212} \\
$\tfrac{11}{10}\le\beta\le\tfrac54$ & $n\ge20$ & $83$ & \lean{cell213}--\lean{cell295} \\
$\tfrac54\le\beta\le\tfrac32$ & $n\ge20$ & $82$ & \lean{cell296}--\lean{cell377} \\
$\tfrac32\le\beta\le2$ & $n\ge20$ & $89$ & \lean{cell378}--\lean{cell466} \\
\midrule
\multicolumn{2}{l}{\emph{total}} & $467$ & \\
\bottomrule
\end{tabular}
\caption{The five chains of certified interval cells. The counts were verified
by listing \file{BerryEsseen/Certificates/ScalarFullH/}, which contains $467$
files \file{Cell$k$.lean} together with \file{Data.lean},
\file{Complete.lean}, \file{Coverage.lean}, \file{LargeMomentCover.lean},
\file{MomentBands.lean} and \file{import\_manifest.json}.}
\label{tab:scalar-cells}
\end{table}

\paragraph{The large-moment chain.}
The first chain has $a_0=\tfrac1{20}$, $b_{136}=\tfrac65$, moment floor
$L=2$, $N_{\min}=1$ and the anchor disabled. By Lemma~\ref{lem:coverage} it
covers every $\ell\in[\tfrac1{20},\tfrac65]$ for laws with $\beta\ge2$; and
$\ell>\tfrac65$ is Corollary~\ref{cor:large}. This proves
Theorem~\ref{thm:scalar}(a).

\paragraph{The four moment bands.}
Each band has $N_{\min}=20$, a common moment floor $L$ and ceiling $U$ taken
from Table~\ref{tab:scalar-cells}, the anchor enabled, and first endpoint
$a\le\tfrac1{20}$. The crucial point is the \emph{last} endpoint: because
$n\ge20$ implies
\begin{equation}\label{eq:sqrt-cap}
  \frac1{\sqrt n} \ \le\ \frac{111803399}{500000000} = 0.223606798 ,
\end{equation}
we have $\ell=\beta/\sqrt n\le U\cdot\tfrac{111803399}{500000000}$, and the Lean
band lemma requires exactly the hypothesis
$U\cdot\tfrac{111803399}{500000000}\le b_{\text{last}}$, i.e.\ that the chain's
last endpoint reaches that rational bound. (The project's notes record that each
band's final endpoint \emph{equals} it.) So no large-sample tail is omitted: for
$n\ge20$ and $L\le\beta\le U$ the actual fraction always lies inside the chain.
Splitting $1\le\beta\le2$ at $\tfrac{11}{10}$,
$\tfrac54$, $\tfrac32$ and appending the large-moment chain for $\beta>2$
proves Theorem~\ref{thm:scalar}(b).

\begin{leanref}
\lean{moment\_cover\_normalizedError} in
\file{.../ScalarFullH/Coverage.lean} (the band argument, including
\eqref{eq:sqrt-cap} through
\lean{reciprocal\_sqrt\_le\_of\_sample\_bound});
\lean{moment\_band1\_normalizedError} through
\lean{moment\_band4\_normalizedError},
\lean{bounded\_moment\_large\_sample\_normalizedError} and
\lean{large\_sample\_normalizedError} in \file{.../MomentBands.lean};
\lean{large\_moment\_middle\_normalizedError} and
\lean{large\_moment\_normalizedError} in \file{.../LargeMomentCover.lean}.
\end{leanref}

\begin{remark}[Scale and size]
The scalar cells use the fixed-point scale $2^{64}$
($\lean{s64}=18446744073709551616$). A single cell file is a few hundred lines
of integer data with one \lean{decide +kernel} per record; for example
\file{Cell0.lean} has $795$ lines, $128$ low panels, $128$ high panels, $128$
normal panels, one tail, and the final budget line
\begin{verbatim}
theorem cell0Bound :
  (sum0 : Rat)/(s64 : Rat) < (293/625 : Rat) * cell0.data.ellLower :=
  by decide +kernel
\end{verbatim}
Its recorded parameters are $a=\tfrac1{20}$, $b=\tfrac{41}{800}=0.05125$,
$L=2$, $T=\tfrac{326460989}{4000000}=81.615\ldots$,
$\tau=\tfrac{1635671}{20000000}=0.0817835\ldots$, $P=128$, $N_{\min}=1$,
anchor disabled.
\end{remark}

%=====================================================================
\section{Region IV: 174 certified finite-sample cells}
\label{sec:finite}
%=====================================================================

The remaining parameters are $1\le n\le19$ with $1\le\beta<2$ and
$\ell<\tfrac65$. Here a different reduction is used, because for such small $n$
the compression of Section~\ref{sec:compression} is wasteful: instead the cutoff
is chosen proportional to $\sqrt n$, which freezes the single-summand frequency.

\begin{theorem}\label{thm:finite}
Let $\mu$ be admissible with $\beta<2$ and let $1\le n<20$. Then
$N(\mu,n,x)\le\tfrac{293}{625}$ for every $x\in\R$.
\end{theorem}

\begin{leanref}
\lean{FiniteParameterBlocks.finite\_normalizedError} in
\file{.../FiniteParameterBlocks/Coverage.lean}.
\end{leanref}

\subsection{The cell reduction}
\label{sec:finite-bounds}

A finite cell specifies integers $1\le N\le n\le H$, a moment interval
$1\le L\le\beta\le B$, a cutoff constant $c>0$ and a split $0<\tau\le\tfrac12$,
and sets
\[
  T_n := c\sqrt n,
  \qquad\text{so that}\qquad
  \frac{T_n s}{\sqrt n} = cs
  \quad\text{is independent of } n .
\]
Choose functions $M$ and $E$ with
\[
  |f(u)| \le M(u)\le1 \quad (0<u<c),
  \qquad
  |f(u)-e^{-u^2/2}| \le u^3E(u),\quad E\ge0 \quad (0<u<c\tau).
\]
All moment-dependent bounds are evaluated at the \emph{upper} endpoint $B$;
their monotonicity in $\beta$ justifies the substitution.

The two error coefficients used by the records are the classical
\[
  E_{\mathrm{classical}}(u) = \frac B6+\frac u8
\]
and, for $0\le u\le\pi$, the anchored
\[
  E_{\mathrm{anchor}}(u) =
  \sqrt{\Bigl(\frac{B-1}{6}+u\,Q(u^2)\Bigr)^2
    + \frac{(B-1)(B+\tfrac53)}{36}} ,
\]
where $Q$ is the explicit rational polynomial with
$0\le e^{-u^2/2}-\cos u\le u^4Q(u^2)$ of Lemma~\ref{lem:gauss-cos}. The
anchored form is exactly Propositions~\ref{prop:real}, \ref{prop:imag}
combined in the Euclidean norm, with the nonnegative factor $u^3$ taken outside
the square root; the case $u=0$ is included.

The available smooth modulus branches are
\[
  M_{\mathrm{line}}(u) = \sqrt{1-u^2+\tfrac{B+1}{5}u^3}
    \quad\bigl((B+1)u\le\tfrac92\bigr),
\]
\[
  M_{\cos}(u) = \sqrt{1-\frac{2a\bigl(1-\cos((B+1)u)\bigr)}{(B+1)^2}}
    \quad\bigl(\tfrac92\le(B+1)u\le2\pi\bigr),
  \qquad a = \frac{81}{80(1-\cos\tfrac92)},
\]
the constant bound $M\equiv1$, and an anchored branch
\[
  M_{\mathrm{anchor}}(u) = \sqrt{\bigl(|\cos u|+\tfrac{(B-1)u^3}{6}\bigr)^2
    + I(u)^2},
\]
with $I$ any of the three proved imaginary-part bounds of
Proposition~\ref{prop:imag}. To differentiate the anchored branch on a panel,
the sign of $\cos u$ is fixed and certified on the whole panel, the expression
under the square root is checked to be at most one, and strict positivity of
the radicand is checked (square-root derivatives require it).

\begin{leanref}
\file{BerryEsseen/FiniteSampleBounds.lean} (the two error coefficients and
their application to the actual i.i.d.\ law),
\file{BerryEsseen/FiniteModulus.lean} (the branch inequalities with all domain
conditions explicit), \file{BerryEsseen/GaussianPolynomial.lean} ($Q$).
\end{leanref}

Combining \eqref{eq:finite-power} with the rational kernel envelopes
\[
  K_0(s) = \sqrt{[s(1-s)]^2+\bigl[\tfrac1\pi-s(1-s)q_-(\pi s)\bigr]^2},
\]
\[
  K_R(s) = (1-s)\sqrt{1+q_+(\pi(1-s))^2},
  \qquad
  K_C(s) = (1-s)\sqrt{1+q_+(\pi s)^2},
\]
for which $s|K(s)|\le K_0(s)$ on $(0,\tfrac12]$, $|K(s)|\le K_R(s)$ on
$[\tfrac12,1)$ and $|K(s)-\tfrac{i}{\pi s}|\le K_C(s)$ on $(0,\tfrac12]$
(Lemmas~\ref{lem:kernel-norms} and \ref{lem:q-rational}), gives the following
theorem. Multiplication by $s$ in $K_0$ is what removes the low-frequency
singularity.

\begin{theorem}[Finite-sample smoothing bound]\label{thm:finite-smoothing}
With the above data define
\[
  \mathcal L(s) = \frac HN c^3s^2K_0(s)E(cs)\,
    G_N\bigl(M(cs),e^{-c^2s^2/2}\bigr),
  \qquad
  \mathcal V(s) = K_C(s)\,e^{-c^2Ns^2/2},
\]
\[
  \mathcal U(s) =
  \begin{cases}
    K_0(s)M(cs)^N/s, & s\le\tfrac12,\\
    K_R(s)M(cs)^N, & s>\tfrac12,
  \end{cases}
  \qquad
  \mathcal J = \frac{e^{-c^2N\tau^2/2}}{\pi c^2N\tau^2} .
\]
Then for every admissible law with $\beta\le B$, every $N\le n\le H$ and every
$x\in\R$,
\[
  \boxed{\ \bigl|F_n(x)-\Phi(x)\bigr| \ \le\
  \mathcal B := \int_0^\tau\mathcal L
    + \int_\tau^1\mathcal U
    + \int_0^\tau\mathcal V
    + \mathcal J . \ }
\]
Consequently, if $a\ge0$ satisfies $H\le(aL)^2$ and
$a\,\mathcal B\le\tfrac{293}{625}$, then $N(\mu,n,x)\le\tfrac{293}{625}$
throughout the cell, because
$\sqrt n/\beta\le\sqrt H/L\le a$.
\end{theorem}

\begin{leanref}
\lean{BerryEsseen.cdf\_le\_finiteSmoothingBound} and
\lean{normalizedError\_le\_finite\_budget} in
\file{BerryEsseen/FiniteSmoothing.lean}. The three integrability hypotheses are
explicit premises of the theorem: numerical data must discharge them, not merely
approximate the integrals. The Gaussian factor decreases in $n$, so evaluating
the normal correction and the tail at $N$ bounds them for all $n\ge N$; the
tail estimate itself is
\[
  \int_\tau^\infty\frac{e^{-c^2ns^2/2}}{\pi s}\,ds
  \le \frac{e^{-c^2N\tau^2/2}}{\pi c^2N\tau^2}
\]
from \file{BerryEsseen/FiniteTails} (see Section~\ref{sec:finite-tails}).
\end{leanref}

\subsection{What a normalized Taylor panel proves}
\label{sec:taylor-panel}

The three integrals of Theorem~\ref{thm:finite-smoothing} are bounded panel by
panel using a signed fourth-order Taylor estimate. This is the single most-used
real-analysis lemma of the numerical part.

\begin{theorem}[Taylor panel bound]\label{thm:taylor-panel}
Let $F$ be four times continuously differentiable on $[m-h,m+h]$, $h\ge0$, and
suppose
\[
  F(m)\le U_0, \qquad F''(m)\le U_2, \qquad
  F''''(s)\le U_4 \quad\text{for \emph{all} } s\in[m-h,m+h].
\]
Then
\begin{equation}\label{eq:taylor-panel}
  \int_{m-h}^{m+h}F(s)\,ds \ \le\ 2hU_0+\frac{h^3}{3}U_2+\frac{h^5}{60}U_4 .
\end{equation}
The bounds $U_0,U_2,U_4$ may be negative.
\end{theorem}

\begin{proof}[Proof sketch]
Taylor's theorem with the integral (or Lagrange) remainder at the midpoint
gives $F(s)\le\sum_{k\le3}\frac{F^{(k)}(m)}{k!}(s-m)^k+\frac{U_4}{24}(s-m)^4$
for $s$ in the panel. The odd powers integrate to zero over the symmetric
interval, $\int(s-m)^2=\tfrac23h^3$ and $\int(s-m)^4=\tfrac25h^5$, and the
fourth-order remainder has a nonnegative weight on both sides of $m$.
\end{proof}

In the certificates the rescaled function $G(t)=F(m+ht)$ is used, and the
stored quantities are enclosures of the normalized coefficients
$G^{(k)}(t)/k!$, which already include the factor $h^k$. In those terms
\eqref{eq:taylor-panel} reads
\begin{equation}\label{eq:taylor-normalized}
  \int_{m-h}^{m+h}F(s)\,ds \ \le\ 2h\Bigl(u_0+\frac{u_2}{3}+\frac{u_4}{5}\Bigr),
\end{equation}
where $u_0\ge G(0)$, $u_2\ge G''(0)/2!$, and --- crucially ---
$u_4\ge G^{(4)}(t)/4!$ for \emph{every} $t\in[-1,1]$, an enclosure on the whole
panel and not a sampling of the derivative.

\begin{leanref}
\lean{taylor\_remainder\_on\_segment}, \lean{cubic\_upper\_of\_fourth\_deriv\_le},
\lean{integral\_centered\_pow}, \lean{midpoint\_integral\_fourth\_bound} in
\file{BerryEsseen/TaylorPanel.lean};
\lean{normalized\_midpoint\_integral\_bound} and
\lean{integral\_le\_normalized\_jet\_bound} in
\file{BerryEsseen/Numerics/TaylorPanelSoundness.lean}, which also connects the
scale-$2^{32}$ integer interval endpoints to the real coefficient bounds.
\end{leanref}

\subsection{Endpoint (Darboux) panels for the high band}
\label{sec:darboux}

On part of the high band a cruder but cheaper estimate is used, in which only
the two endpoints of the panel matter. This works because the line-branch
radicand is a cubic with a single interior minimum.

\begin{lemma}\label{lem:radicand-range}
For $B\ge1$ let $P_B(u)=1-u^2+\tfrac{B+1}{5}u^3$ and
$m_B(u)=\sqrt{P_B(u)}$. Then
\[
  P_B'(u) = \frac u5\bigl(3(B+1)u-10\bigr),
\]
so the unique interior critical point $u_*=\tfrac{10}{3(B+1)}$ is a minimum, and
\[
  P_B(u) - \Bigl(1-\frac{100}{27(B+1)^2}\Bigr)
  = \frac{\bigl(3(B+1)u-10\bigr)^2\bigl(3(B+1)u+5\bigr)}{135(B+1)^2}
  \ \ge\ 0 \qquad (u\ge0).
\]
Consequently, on any panel $[l,r]$ with $(B+1)cr\le\tfrac92$,
\[
  m_B(cs) \ \le\ v := \max\{m_B(cl),\,m_B(cr)\} \qquad (l\le s\le r).
\]
\end{lemma}

\begin{leanref}
\file{BerryEsseen/FiniteRadicandRange.lean}.
\end{leanref}

\begin{proposition}[Endpoint panel bound]\label{prop:darboux}
Let $0<l\le s\le r\le1$, $(B+1)cr\le\tfrac92$, $1\le N\le n$, $T=c\sqrt n$.
Then
\begin{equation}\label{eq:darboux}
  \int_l^r |K(s)|\,\bigl|\hat\mu_n(Ts)\bigr|\,ds
  \ \le\ (r-l)\Bigl(1-\frac{l+r}{2}\Bigr)
    \sqrt{1+\Bigl(\frac{1+l}{\pi l}\Bigr)^2}\;v^N .
\end{equation}
\end{proposition}

\begin{proof}
By \eqref{eq:finite-power} and Lemma~\ref{lem:radicand-range},
$|\hat\mu_n(Ts)|\le m_B(cs)^N\le v^N$. By Lemma~\ref{lem:kernel-norms} and
Lemma~\ref{lem:q-reflected},
\[
  |K(s)| = (1-s)\sqrt{1+q(\pi(1-s))^2}
  \le (1-s)\sqrt{1+\Bigl(\frac{1+s}{\pi s}\Bigr)^2}
  \le (1-s)\sqrt{1+\Bigl(\frac{1+l}{\pi l}\Bigr)^2},
\]
because $(1+s)/(\pi s)$ decreases. Integrating the remaining affine factor
exactly, $\int_l^r(1-s)ds=(r-l)(1-\tfrac{l+r}{2})$, gives
\eqref{eq:darboux}.
\end{proof}

The certificate for such a panel is entirely rational. With
$p=13493037704/2^{32}<\pi$ it proposes nonnegative rationals
\[
  M \ge \max\{\sqrt{P_B(cl)},\sqrt{P_B(cr)}\},
  \qquad
  \Lambda \ge \sqrt{1+\bigl((1+l)/(pl)\bigr)^2},
\]
which Lean verifies by squaring, checks the parameter and branch conditions, and
checks the rational inequality
$(r-l)\bigl(1-\tfrac{l+r}{2}\bigr)\Lambda M^N\le U$. The conclusion is that the
actual integral is at most $U$, uniformly over every admissible law with
$\beta\le B$ and every $n\ge N$.

\begin{leanref}
\file{BerryEsseen/FiniteDarboux.lean} (the analytic statements) and
\file{BerryEsseen/Numerics/DarbouxCertificate.lean} (the checker and its
soundness theorems \lean{Panel.sound}, \lean{CertifiedPanel.sound}). The $458$
instances are collected in
\file{BerryEsseen/Certificates/FiniteDarbouxPanels/Complete.lean}; their
manifest records $174$ rows, $4034$ source panels and $458$ generated entries.
\end{leanref}

\subsection{Normal-correction panels}

For a cell with $N\le n\le H$, $T=c\sqrt n$ and $0<\tau\le\tfrac12$, the
normal-correction term is bounded by
\[
  \int_0^\tau (1-s)\sqrt{1+q_+(\pi s)^2}\;e^{-c^2Ns^2/2}\,ds ,
\]
with $q_+$ the rational upper bound of Lemma~\ref{lem:q-rational}, whose
denominator is positive on $[0,\tfrac\pi2]$. Increasing $n$ decreases the
Gaussian factor, so a bound at $N$ applies for all $n\ge N$.

Each panel proof supplies the hypotheses of
\eqref{eq:taylor-normalized} as follows.
\begin{enumerate}[leftmargin=2.2em]
\item Every arithmetic node carries an exact function of $t$, a proof that it is
  four times continuously differentiable at each point of its domain, and five
  interval enclosures of its normalized derivatives.
\item Constants use the $693$ exact scalar enclosures of
  Section~\ref{sec:scalar-constants}; arithmetic uses the proved Taylor
  recurrences; inverses and square roots carry checked domain guards;
  exponentials use certified seeds.
\item One evaluation encloses the derivatives at $t=0$ and a second encloses
  them throughout $[-1,1]$, and Lean checks that the two evaluations represent
  the \emph{same} function.
\item Lean proves an explicit identity between that function and the displayed
  normal envelope after the affine change of variable, unfolding the rational
  polynomials and simplifying the exact arithmetic.
\item \eqref{eq:taylor-normalized} yields a rational panel bound, and the
  analytic envelope inequality transfers it to the actual Prawitz correction
  integral.
\end{enumerate}

There are $669$ normal panels across the $174$ cells. A cell's list is checked
to begin at $0$, to have adjacent endpoints, to end exactly at its split, and to
use the same cutoff and sample-size floor throughout; the endpoint check also
verifies that all intervals lie in the required half-period.

\begin{leanref}
\file{BerryEsseen/FiniteNormalPanels.lean},
\file{BerryEsseen/FiniteNormalCover.lean},
\file{BerryEsseen/Certificates/FiniteNormalPanels/} ($669$ files
\file{Panel$i$\_$j$.lean} plus \file{Complete.lean} and \file{manifest.json}).
\end{leanref}

\subsection{High- and low-frequency Taylor panels}

The high-frequency Taylor panels bound
\[
  \int_l^r K_R(s)\,M_B(cs)^N\,ds
  \qquad\text{on } [l,r]\subseteq[\tfrac12,1],
\]
and on a left-half panel the kernel factor is instead $K_0(s)/s$ with
$0<l\le r\le\tfrac12$, whose denominator and sine-polynomial domain are checked
throughout the panel. The branch conditions are checked as inequalities valid on
the whole panel: the line branch checks $(B+1)cr\le\tfrac92$, the cosine branch
checks both $(B+1)cl\ge\tfrac92$ and $(B+1)cr\le2\pi$. Continuity of the kernel
envelope and of the modulus gives integrability of the target, including at the
endpoint $s=1$. Anchored panels additionally use whole-panel cosine-sign and
squared-modulus bounds to prove the sign and the bound $M\le1$ at every real
point of the panel; no sign or domain condition is inferred from a sampled
function value.

The low-frequency Taylor panels bound $\int\mathcal L$ using the previously
proved sample-interval comparison \eqref{eq:geom-compare} and either the
anchored or the classical error coefficient. The special case $N=1$ has
geometric factor one and therefore does not require an unused modulus branch.

Counts: $1806$ high-frequency Taylor panels and $1101$ low-frequency Taylor
panels, verified against the two manifests
(\lean{generated\_entries}=\lean{expected\_entries} with
\lean{complete\_cover}=\lean{true} in both).

\begin{leanref}
\file{BerryEsseen/FiniteHighPanels.lean},
\file{BerryEsseen/FiniteLowPanels.lean},
\file{BerryEsseen/FiniteHighCover.lean},
\file{BerryEsseen/FiniteLowCover.lean};
\file{BerryEsseen/Certificates/FiniteHighTaylorPanels/} ($91$ batch files) and
\file{.../FiniteLowTaylorPanels/} ($111$ batch files);
\lean{Numerics/JetPanelRange.lean} for the affine-inverse rule that turns a
proved whole-panel zeroth coefficient into a pointwise bound.
\end{leanref}

\subsection{Gaussian tails}
\label{sec:finite-tails}

For $n\ge N>0$, $c>0$, $\tau>0$, put $Q=c^2N\tau^2$, a positive rational in
every saved cell, so that the exponent is the exact rational $-Q/2$. The
checker verifies: the sample floor, cutoff and split are positive; the
exponential seed's input interval contains $-Q/2$; and its certified output
upper endpoint $\mathcal E$ satisfies
$\mathcal E/(\pi_{\mathrm{low}}Q)\le U$ with
$\pi_{\mathrm{low}}=13493037704/2^{32}\le\pi$. Monotonicity of the exponential
and division by positive denominators then prove the real tail bound. All $174$
cells have checked tail certificates, and each applies to every sample size
above the recorded floor.

\begin{leanref}
\file{BerryEsseen/Numerics/FiniteTailCertificate.lean}
(\lean{tailCheck}, \lean{CertifiedTail.envelope\_bound},
\lean{CertifiedTail.sound}) and
\file{BerryEsseen/Certificates/FiniteTails/Complete.lean}.
\end{leanref}

\subsection{Cell and block assembly}

\begin{definition}\label{def:finite-cell}
A \emph{certified finite cell} bundles a certified low cell, high cell, normal
cell and tail, all with the same $N$, $H$, $B$, $c$, $\tau$, together with
rationals $L$ (moment floor), $a$ (normalization factor) and the constant
$\tfrac{293}{625}$, and a Boolean check that includes
\[
  0<N,\quad 0<c,\quad 0<\tau\le1,\quad 0<L\le B,\quad a\ge0,\quad
  H\le(aL)^2,\quad 0\le \mathcal U,\quad a\,\mathcal U\le\frac{293}{625},
\]
where $\mathcal U$ is the sum of the four stored rational upper bounds.
\end{definition}

\begin{leanref}
\lean{finiteCellCheck}, \lean{CertifiedFiniteCell} and
\lean{CertifiedFiniteCell.sound} in
\file{BerryEsseen/FiniteCellCertificate.lean}. The soundness proof combines
Theorem~\ref{thm:finite-smoothing} with $\sqrt H/L\le a$ (obtained from
$H\le(aL)^2$ by taking square roots) and Lemma~\ref{lem:cdf-to-normalized}.
\end{leanref}

The $174$ certified finite cells are grouped into $17$ blocks by sample-size
range, each block chaining moment intervals so that adjacent cells meet exactly.
Table~\ref{tab:finite-blocks} lists them.

\begin{table}[ht]
\centering
\begin{tabular}{@{}lrr@{}}
\toprule
Sample sizes $[N,H]$ & Moment ceiling & Cells \\
\midrule
$[1,1]$   & $\tfrac{12001}{10000}$ & $2$ \\
$[2,2]$   & $\tfrac{16971}{10000}$ & $9$ \\
$[3,3]$   & $2$ & $13$ \\
$[4,4]$   & $2$ & $14$ \\
$[5,5]$   & $2$ & $14$ \\
$[6,6]$   & $2$ & $13$ \\
$[7,7]$   & $2$ & $11$ \\
$[8,8]$   & $2$ & $9$ \\
$[9,9]$   & $2$ & $9$ \\
$[10,10]$ & $2$ & $9$ \\
$[11,11]$ & $2$ & $9$ \\
$[12,12]$ & $2$ & $9$ \\
$[13,13]$ & $2$ & $9$ \\
$[14,14]$ & $2$ & $9$ \\
$[15,16]$ & $2$ & $14$ \\
$[17,18]$ & $2$ & $13$ \\
$[19,19]$ & $2$ & $8$ \\
\midrule
\emph{total} & & $174$ \\
\bottomrule
\end{tabular}
\caption{The $17$ finite-sample parameter blocks. The counts were verified by
reading \file{BerryEsseen/Certificates/FiniteParameterBlocks/Block$k$.lean} and
counting the \lean{FiniteCompleteCells.cell$j$} references in each.}
\label{tab:finite-blocks}
\end{table}

\begin{proof}[Proof of Theorem~\ref{thm:finite}]
If $\ell\ge\tfrac65$ apply Corollary~\ref{cor:large}. Otherwise $\beta<2$ and
$\ell<\tfrac65$; the elementary lemma
\[
  \beta<2 \ \wedge\ \frac{\beta}{\sqrt n}<\frac65
  \ \Longrightarrow\
  \beta < \begin{cases}
    \tfrac{12001}{10000}, & n=1,\\
    \tfrac{16971}{10000}, & n=2,\\
    2, & n\ge3,
  \end{cases}
\]
(for $n=1$ this is immediate; for $n=2$ it uses
$\sqrt2\le\tfrac{28285}{20000}$) reduces the moment range to the block's
ceiling. Since $1\le n\le19$, one of the $17$ blocks applies; its chained
moment intervals cover $[1,\text{ceiling})$ by the same argument as
Lemma~\ref{lem:coverage}, and Definition~\ref{def:finite-cell} gives the bound
$\tfrac{293}{625}$.
\end{proof}

\begin{leanref}
\lean{finiteMomentCap} and \lean{lt\_finiteMomentCap} in
\file{BerryEsseen/FiniteMomentCap.lean};
\lean{FiniteParameterCover.lean} for the adjacency of moment intervals;
\lean{FiniteParameterBlocks.finite\_small\_moment\_bound} and
\lean{finite\_normalizedError} in
\file{.../FiniteParameterBlocks/Coverage.lean}. The Lean proof performs an
\lean{interval\_cases n} over $1\le n\le19$, dispatching $n=16$ to the
$[15,16]$ block and $n=18$ to the $[17,18]$ block.
\end{leanref}

\begin{remark}[The tightest saved budget]
The manifest of \file{FiniteCompleteCells} records that the largest normalized
budget among the $174$ cells is at row $48$, namely
\[
  \frac{41323819126424383349609375}{88230205055486500383227904}
  = 0.4683636301245874\ldots \ <\ \frac{293}{625}=0.4688 .
\]
That cell has $N=H=5$, $B=\tfrac{3^{18}}{5^{12}}=1.58672\ldots$,
$c=\tfrac{239773}{100000}$, $\tau=\tfrac{36843}{100000}$,
$L=\tfrac{3^{15}}{5^{10}}=1.46932\ldots$,
$a=\tfrac{93787488623046875}{61628086298345472}=1.52185\ldots$, and nine
low-frequency panels. The margin to the target is about $0.0004$, so the
budgets are tight but not marginal.
\end{remark}

%=====================================================================
\section{From integers to real integrals}
\label{sec:numerics}
%=====================================================================

The numerical part of the proof consists of exact integer and rational
inequalities. Its correctness rests on a chain of soundness theorems, each of
which says that a particular \emph{passing check} implies a particular
\emph{real} statement. A list of passing integer inequalities is never treated
as a proof of a probability bound on its own.

\subsection{Interval arithmetic with outward rounding}

\begin{definition}\label{def:interval}
Fix a scale $S$. An integer pair $[\alpha,\beta]$ \emph{encloses} a real $x$
when $\alpha\le Sx\le\beta$. Two scales occur:
\[
  S = 2^{32} = 4294967296
  \quad\text{(the Taylor-jet records of the finite-sample cells)},
\]
\[
  S = 2^{64} = 18446744073709551616
  \quad\text{(the interval cells of Section~\ref{sec:scalar})}.
\]
\end{definition}

Every arithmetic operation rounds outwards. Addition and negation are exact on
endpoints; multiplication takes the smallest and largest of the four endpoint
products, rounding the former down and the latter up after division by $S$.
Square-root witnesses are verified by squaring integers; reciprocals carry a
nonzero-denominator guard.

\begin{leanref}
\file{BerryEsseen/Numerics/JetCore.lean} defines the interval type and the
executable operations at scale $2^{32}$;
\file{BerryEsseen/Numerics/IntervalSoundness.lean} proves that each rule
preserves real enclosures, including for intervals with endpoints of different
signs (\lean{Contains.add}, \lean{Contains.neg}, \lean{Contains.sub},
\lean{Contains.mul}, \lean{Contains.scaleNat}, \lean{Contains.divConst}, and the
outward-rounding lemmas \lean{floorScale\_bound}, \lean{ceilScale\_bound}).
\end{leanref}

\subsection{Taylor jets: from an expression to derivative bounds}

A \emph{jet} stores the five normalized coefficients
$c_j=F^{(j)}(x)/j!$, $j=0,\dots,4$, as intervals. Products use ordinary
convolution $(fg)_k=\sum_{j\le k}f_jg_{k-j}$, and the other operations follow
from identities that determine the next coefficient from earlier ones:
\begin{center}
\begin{tabular}{@{}ll@{}}
\toprule
Operation & Determining identity \\
\midrule
$y=1/f$ & $fy=1$ \\
$y=\sqrt f$ & $y^2=f$ \\
$y=e^{f}$ & $y'=f'y$ \\
$(s,c)=(\sin f,\cos f)$ & $s'=f'c$, $c'=-f's$ \\
\bottomrule
\end{tabular}
\end{center}
The Lean proofs connect these recurrences to \emph{actual derivatives} of the
functions involved, and they require a nonzero denominator or a positive
square-root argument where the identity needs one; the executable guards
establish these conditions. An exponential or trigonometric node additionally
needs a valid interval for its \emph{zeroth} coefficient, which the recurrence
cannot supply; that premise is explicit, and is discharged by the seed
certificates of Sections~\ref{sec:exp-seeds} and \ref{sec:trig-seeds}.

\begin{leanref}
\file{BerryEsseen/Numerics/DifferentialJet.lean},
\file{.../InverseJet.lean}, \file{.../SquareRootJet.lean},
\file{.../ExponentialJet.lean}, \file{.../TrigonometricJet.lean},
\file{.../JetSoundness.lean}, and the whole-domain form
\lean{EnclosesOn} in \file{.../EnclosesOn.lean}, whose constructors
(\lean{EnclosesOn.add}, \lean{.mul}, \lean{.inv}, \lean{.sqrt}, \lean{.exp},
\lean{.sincos}, \lean{.affine}, \lean{.constant}, \lean{.refine\_radicand})
carry the domain conditions as hypotheses. The affine input rule and the rule
for replacing a zeroth interval by a separately proved tighter enclosure are in
\file{.../JetInputs.lean}.
\end{leanref}

\subsection{Scalar constants keep their exact real meanings}
\label{sec:scalar-constants}

A rational interval alone does not identify which real constant it encloses: a
rounded interval may enclose several different rationals. The construction
therefore records the exact real value \emph{in the type}:
\begin{verbatim}
structure ScalarEnclosure (value : Real) where
  interval : Interval
  contains : interval.Contains value
\end{verbatim}
The parameter \lean{value} is part of the mathematical statement, so it cannot
change when an arithmetic operation or an interval widening is applied. Each
constructor produces an enclosure of the corresponding exact real operation, and
the theorem \lean{encloses\_constant} initializes a constant Taylor jet from
such an enclosure.

The finite Taylor formulas use $69{,}759$ constant occurrences across $3{,}576$
panels. A complete replay reduces them to $647$ distinct constant records plus
$46$ supporting scalar operations, i.e.\ $693$ scalar enclosures in all, namely
$506$ rational enclosures, one $\pi$ enclosure, one cosine, one subtraction,
two reciprocals, $152$ products and $30$ square roots. Among the constants
actually formed are
\[
  c\sqrt N, \qquad \sqrt{(B-1)(B+\tfrac53)}, \qquad \sqrt{B-1},
  \qquad \frac1\pi, \qquad \frac{81}{80\,(1-\cos\tfrac92)} ,
\]
together with the rational polynomial coefficients and the sample and moment
endpoints. The square-root rule permits the value zero, which is necessary for
$\sqrt{B-1}$ when $B=1$; trigonometric constants use the already proved
endpoint and range certificates.

\begin{leanref}
\file{BerryEsseen/Numerics/ScalarEnclosure.lean} (definitions and soundness
rules: \lean{rational}, \lean{pi}, \lean{add}, \lean{neg}, \lean{sub},
\lean{mul}, \lean{divNat}, \lean{inv}, \lean{sqrt}, \lean{exp}, \lean{sin},
\lean{cos}, \lean{widen}) and
\file{BerryEsseen/Certificates/FiniteScalarConstants/Complete.lean} (the $693$
checked records; the directory has seven batch files and a manifest recording
\lean{generated\_entries}=$693$). The public data file
\file{data/finite\_scalar\_meanings.json} records all scalar nodes, all
$69{,}759$ occurrence mappings and the $3{,}576$ sets of coefficient targets; it
is untrusted input to certificate construction.
\end{leanref}

\subsection{Exponential seeds}
\label{sec:exp-seeds}

The finite-sample programs contain $3{,}514$ exponential operations. For each,
the starting interval for $e^{x/S}$ ($S=2^{32}$) is certified as follows. Choose
$k\ge0$ and put $r=x/(S2^k)$, requiring $|r|\le1$. For any $n\ge1$, Taylor's
theorem gives
\[
  \Bigl|e^r-\sum_{j<n}\frac{r^j}{j!}\Bigr| \le R := \frac{2|r|^n}{n!} .
\]
Both the polynomial $P$ and $R$ are exact rationals, and a proposed seed
$[\alpha,\beta]$ is accepted only if $\alpha\le S(P-R)$ and $S(P+R)\le\beta$,
which implies $\alpha\le S e^r\le\beta$. The original argument is recovered by
squaring the interval outwards $k$ times using the proved integer multiplication
rule, which encloses $(e^r)^{2^k}=e^{x/S}$; the final proposed interval must
\emph{contain} the computed one, so harmless widening is allowed but unjustified
narrowing is not. Two shortcuts have separate proofs: $x=0$ gives $e^0=1$, and
$x/S\le-32$ gives $e^{x/S}\le e^{-32}\le2^{-32}$ (using $e\ge2$), so $[0,1]$ is
valid. Finally, for an input interval $[l,u]$ one certifies both endpoints and
uses monotonicity of $\exp$ to cover every real $y$ with $l\le Sy\le u$; the
lower bound may also be raised to $0$ since $\exp>0$. No floating-point
exponential occurs anywhere in the check.

\begin{leanref}
\file{BerryEsseen/Numerics/ExpSeedCore.lean} (executable check),
\file{.../ExpSeedSoundness.lean} (\lean{expIntervalCheck\_sound}),
\file{.../CertifiedExponentialJet.lean} (\lean{Encloses.exp\_certified},
which requires the checked input interval to \emph{equal} the actual zeroth jet
coefficient), and
\file{BerryEsseen/Certificates/FiniteExpSeeds/} ($36$ batch files plus
\file{Complete.lean}; manifest: $36$ sources, $3{,}576$ panels,
$3{,}514$ unique entries).
\end{leanref}

\subsection{Trigonometric seeds}
\label{sec:trig-seeds}

The frozen programs contain $4{,}248$ paired sine/cosine calls with input
intervals inside $[0,2\pi]$. Each certificate proves bounds for every real input
in its interval, including intervals crossing a critical point. The proof has two
parts.

\paragraph{Endpoint values.}
The fixed interval $[13493037704,\,13493037705]$ encloses $S\pi$ at
$S=2^{32}$, which follows from mathlib's proved decimal bounds on $\pi$.
Choosing a quadrant witness $k\in\{0,1,2,3,4\}$, outward integer arithmetic
encloses $z=x/S-k\pi/2$, and the checker verifies that this interval lies in
$[-1,1]$; it need not prove that $k$ is the nearest quadrant, since only the
inclusion and the trigonometric identities are used. For the reduced argument,
\[
  C(z) = 1-\frac{z^2}{2!}+\dots+\frac{z^{12}}{12!},
  \qquad
  S(z) = z-\frac{z^3}{3!}+\dots+\frac{z^{13}}{13!} ,
\]
and since the next Taylor coefficient vanishes in each case,
\[
  |\cos z-C(z)|\le\frac{|z|^{14}}{14!},
  \qquad
  |\sin z-S(z)|\le\frac{|z|^{15}}{15!},
\]
both smaller than $2^{-32}$ for $|z|\le1$. Each polynomial is evaluated by
outward Horner arithmetic and each endpoint widened by one integer unit. The
original values are reconstructed by the quadrant identities
\[
  \begin{array}{c|cc}
    k\bmod4 & \sin(x/S) & \cos(x/S)\\\hline
    0 & \sin z & \cos z\\
    1 & \cos z & -\sin z\\
    2 & -\sin z & -\cos z\\
    3 & -\cos z & \sin z
  \end{array}
\]
with a separate shortcut at $x=0$.

\paragraph{Values between the endpoints.}
A continuous differentiable function on a closed interval attains its extrema at
an endpoint or at an interior critical point; the Lean proof uses the
extreme-value theorem and Fermat's theorem, and does \emph{not} infer a range
from sampled values. Inside $(0,2\pi)$ the relevant critical data are
\[
  \sin:\ \pi/2\mapsto1,\quad 3\pi/2\mapsto-1;
  \qquad
  \cos:\ \pi\mapsto-1 .
\]
The checker uses outward intervals for these locations: if a critical-location
interval overlaps the input interval, the output must include the corresponding
value $\pm1$. An overlap may be conservative but cannot cause a critical point
to be missed. Endpoint enclosures may also be intersected with $[-1,1]$.

\begin{leanref}
\file{BerryEsseen/Numerics/TrigPolynomialSoundness.lean},
\file{.../TrigEndpointSoundness.lean},
\file{BerryEsseen/TrigonometricRange.lean} (critical-point classification and
range theorem), \file{.../TrigSeedCore.lean},
\file{.../TrigSeedSoundness.lean} (\lean{trigIntervalCheck\_sound}),
\file{.../CertifiedTrigonometricJet.lean} (\lean{Encloses.sincos\_certified}),
and \file{BerryEsseen/Certificates/FiniteTrigSeeds/} ($43$ batch files;
manifest: $4{,}248$ unique entries).
\end{leanref}

\subsection{Radicand refinements}

The finite Taylor programs contain $2{,}884$ operations that replace the
interval for a function's value by a smaller one while leaving its four
derivative coefficients unchanged; $1{,}540$ concern the cubic branch and
$1{,}344$ the cosine branch. This is a \emph{conditional composition} result:
the original interval must already contain the stated real radicand and its
argument must lie in the stated input interval.

\paragraph{Cubic branch.}
For a rational $B\ge1$ and $P_B(u)=1-u^2+\tfrac{B+1}{5}u^3$ on a nonnegative
interval $[a,b]$, the only possible interior critical point is
$u_*=\tfrac{10}{3(B+1)}$, so the range is controlled by the two endpoint values
together with, if $u_*\in[a,b]$, the additional rational value
$P_B(u_*)=1-\tfrac{100}{27(B+1)^2}$. The checker tests that the proposed range
contains the endpoint values and either excludes $u_*$ or contains its value;
the proof uses compactness to attain the extrema and Fermat's theorem at an
interior one.

\paragraph{Cosine branch.}
Here the radicand is
\[
  R_B(u) = 1-\frac{2A\bigl(1-\cos((B+1)u)\bigr)}{(B+1)^2},
  \qquad A = \frac{81}{80(1-\cos\tfrac92)} ,
\]
and $-1\le\cos\le1$ gives $1-\tfrac{4A}{(B+1)^2}\le R_B(u)\le1$ for every real
$u$. A separately checked sine/cosine endpoint certificate at $9/2$ proves
$A\le A_+:=\tfrac{3591567126}{2^{32}}$, so the rational interval
$\bigl[1-\tfrac{4A_+}{(B+1)^2},\,1\bigr]$ contains $R_B(u)$ for every real $u$;
in particular this stays valid if outward rounding pushes the numerical input
slightly past a branch endpoint. The branch conditions used to bound the
characteristic function are separate obligations.

\paragraph{Intersection.}
If $I\ni F(x)$ is the original interval and the analytic argument proves
$J\ni F(x)$, then $I\cap J\ni F(x)$, and the checker verifies
$I\cap J\subseteq I_{\mathrm{new}}$. Merely checking
$I_{\mathrm{new}}\subseteq I$ would \emph{not} prove $F(x)\in I_{\mathrm{new}}$.
Only the zeroth enclosure changes; the real function and all derivative
enclosures are the same.

\begin{leanref}
\file{BerryEsseen/Numerics/RadicandRefinement.lean}
(\lean{Jet.Encloses.refine\_radicand}, with the original interval and the
real-expression identity as explicit premises) and
\file{BerryEsseen/Certificates/FiniteRadicandRefinements/Complete.lean}
($2{,}884$ instances across $29$ batch files).
\end{leanref}

\subsection{Cotangent tables}

The kernel tables store integer bounds for $q(\pi k/4096)$ (both directions),
for $e^{-k/1024}$, and for $\cos(k/4096)$. Their meanings are proved as
follows.
\begin{itemize}[leftmargin=2.2em]
\item \emph{Exponential entries.} The integer polynomial is exactly the
  degree-$24$ Taylor polynomial for $e^{k/1024}$ with cleared denominators;
  every omitted term is positive, so taking reciprocals gives the required
  \emph{upper} bound for $e^{-k/1024}$.
\item \emph{Cosine entries.} The polynomial has degree $32$. On $[0,7]$, group
  the omitted terms of the absolutely convergent cosine series into
  negative--positive pairs starting at degrees $34$ and $36$; each pair is
  nonpositive because $x^2\le(m+1)(m+2)$ for $m\ge34$ and $0\le x\le7$. Hence
  the degree-$32$ polynomial is an upper bound throughout the table's argument
  range. The alternate branch using the trivial bound $\cos\le1$ is also proved.
\item \emph{Cotangent entries.} These are exactly
  \eqref{eq:q-sandwich}, with the positivity guards $A_-\ge0$ and $D_->0$ being
  the ones the integer tables check, plus the conversion from the rational
  bounds on $\pi$ to the intended table argument.
\end{itemize}

\begin{leanref}
\file{BerryEsseen/Numerics/ScalarTableCore.lean} (unchanged executable
arithmetic from the saved certificate, at scale $2^{64}$),
\file{.../ScalarTables.lean} (\lean{tableEntry\_exp},
\lean{tableEntry\_cos}, \lean{tableEntry\_q\_lower},
\lean{tableEntry\_q\_upper} and their ``at a point'' variants),
\file{.../CotangentTable.lean},
\file{BerryEsseen/CotangentPolynomials.lean},
\file{BerryEsseen/CotangentCorrection.lean};
\file{BerryEsseen/ScalarRemainders.lean} justifies replacing $\exp$, $\sin$ and
$\cos$ by rational polynomials with explicit error intervals, and
\file{BerryEsseen/ExponentialRemainder.lean} supplies the exponential
remainders.
\end{leanref}

\subsection{Inventory of the checked collection}

Table~\ref{tab:counts} gives the counts of the certified objects, each verified
against the corresponding generated manifest or by listing the directory.

\begin{table}[ht]
\centering
\begin{tabular}{@{}lrl@{}}
\toprule
Object & Count & Location \\
\midrule
Interval cells (scale $2^{64}$) & $467$ & \file{Certificates/ScalarFullH/} \\
Finite-sample complete cells & $174$ & \file{Certificates/FiniteCompleteCells/} \\
Finite-sample parameter blocks & $17$ & \file{Certificates/FiniteParameterBlocks/} \\
High-frequency Taylor panels & $1806$ & \file{Certificates/FiniteHighTaylorPanels/} \\
Low-frequency Taylor panels & $1101$ & \file{Certificates/FiniteLowTaylorPanels/} \\
Normal-correction Taylor panels & $669$ & \file{Certificates/FiniteNormalPanels/} \\
Endpoint (Darboux) panels & $458$ & \file{Certificates/FiniteDarbouxPanels/} \\
Gaussian tails & $174$ & \file{Certificates/FiniteTails/} \\
Exponential seeds & $3514$ & \file{Certificates/FiniteExpSeeds/} \\
Paired sine/cosine seeds & $4248$ & \file{Certificates/FiniteTrigSeeds/} \\
Radicand refinements & $2884$ & \file{Certificates/FiniteRadicandRefinements/} \\
Scalar constant enclosures & $693$ & \file{Certificates/FiniteScalarConstants/} \\
\midrule
\multicolumn{3}{@{}l}{\emph{Illustrative individual panel files (not extra
  certificates):}}\\
High-frequency example panels & $7$ & \file{Certificates/FiniteHighPanels/} \\
Low-frequency example panels & $6$ & \file{Certificates/FiniteLowPanels/} \\
\bottomrule
\end{tabular}
\caption{The checked certificate collection. The seven high-frequency example
panels --- $(\text{row},\text{panel})=(0,0),(0,1),(0,2),(0,7),(0,23),(166,1),
(173,30)$ --- cover all six branch/kernel combinations, both cosine signs, and
sample-size floor $19$; they are imported by the supporting library to
illustrate the general construction, while the complete collection lives in the
batch files.}
\label{tab:counts}
\end{table}

\begin{table}[ht]
\centering
\begin{tabular}{@{}lr@{}}
\toprule
Directory & \lean{decide +kernel} occurrences \\
\midrule
\file{ScalarFullH}                  & $349{,}824$ \\
\file{FiniteLowTaylorPanels}        & $263{,}610$ \\
\file{FiniteHighTaylorPanels}       & $256{,}223$ \\
\file{FiniteNormalPanels}           & $76{,}936$ \\
\file{FiniteTrigSeeds}              & $4{,}249$ \\
\file{FiniteExpSeeds}               & $3{,}515$ \\
\file{FiniteRadicandRefinements}    & $2{,}885$ \\
\file{FiniteLowPanels}              & $1{,}276$ \\
\file{FiniteHighPanels}             & $973$ \\
\file{FiniteScalarConstants}        & $694$ \\
\file{FiniteDarbouxPanels}          & $459$ \\
\file{FiniteTails}                  & $349$ \\
\file{FiniteCompleteCells}          & $175$ \\
\file{FiniteHighCells}              & $175$ \\
\file{FiniteLowCells}               & $175$ \\
\file{FiniteNormalCells}            & $175$ \\
\file{FiniteParameterBlocks}        & $17$ \\
\midrule
\textbf{Total} & $\mathbf{961{,}710}$ \\
\bottomrule
\end{tabular}
\caption{Kernel-checked finite decisions, counted over
\file{BerryEsseen/Certificates/} with
\texttt{grep -rc "decide +kernel" ... | awk -F: '\{s+=\$2\} END \{print s\}'}.
Every one of these is evaluated by Lean's kernel on exact arithmetic; none uses
\lean{native\_decide}.}
\label{tab:kernel}
\end{table}

\begin{remark}[Recorded rational margins]
The module \file{BerryEsseen/CertificateMargins.lean} records, as rational
inequalities only, the following strict comparisons with the target
$\tfrac{293}{625}=0.4688$:
\[
  \begin{array}{lll}
    \lean{small\_ell\_margin}: & \tfrac{4473}{10000} = 0.4473 & <\ \tfrac{293}{625},\\
    \lean{large\_ell\_margin}: & \tfrac{11}{24} = 0.4583\overline3 & <\ \tfrac{293}{625},\\
    \lean{all\_n\_margin}: & \tfrac{4302436550217957753}{2^{63}} = 0.4664711\ldots & <\ \tfrac{293}{625},\\
    \lean{band\_margin}: & \tfrac{8647540607797249413}{2^{64}} = 0.4687841\ldots & <\ \tfrac{293}{625},\\
    \lean{finite\_margin}: & \tfrac{2011606487}{2^{32}} = 0.4683636\ldots & <\ \tfrac{293}{625}.
  \end{array}
\]
The module's own docstring states that these are rational inequalities only and
that connecting the numbers to probability errors requires the analytic
soundness and complete certificate proofs, so this module alone is not an upper
bound for $\Cii$. The proved regional constant for the small-fraction region is
the $\tfrac{183}{400}$ of Theorem~\ref{thm:small}.
\end{remark}

\subsection{What is not trusted}

The Python generators are \emph{outside} the trusted proof. They propose
cutoffs, subdivisions, panel counts, rational bounds, seeds and refinement
witnesses; Lean then verifies the meaning and correctness of the resulting
certificates. Concretely:
\begin{itemize}[leftmargin=2.2em]
\item No floating-point result, external interval library, or producer success
  flag is a premise of any Lean theorem.
\item Source hashes and manifests provide reproducibility, not mathematical
  content; they are \file{import\_manifest.json} and \file{manifest.json} in
  the certificate directories, and \file{finite\_scalar\_meanings.json} and
  \file{finite20\_coalesced.jsonl} under \file{data/}.
\item A floating-point maximum, or a collection of positive sampled margins, is
  never used as a substitute for the interval inequalities.
\item The import script for the interval cells preserves the frozen integer
  records and their arithmetic conditions; it changes module wrappers, removes
  diagnostic printing, tightens a former $0.4689$ budget to $0.4688$, and uses
  the proved \lean{CertifiedCell} type for the final assembly.
\end{itemize}

Reproduction (from a checkout with Lean and the mathlib cache installed):
\begin{verbatim}
lake exe cache get
python3 scripts/build_certificates.py            # or --batch-size 2 for CI
\end{verbatim}
The batching limits the number of concurrent compiler processes; it does not
change any proof check. One build driver per checkout should be used, since
overlapping Lake invocations can observe partially written dependency outputs.

%=====================================================================
\section{Assembling the upper bound}
\label{sec:assembly}
%=====================================================================

\begin{theorem}[Pointwise universal bound]\label{thm:pointwise}
Let $\mu$ be admissible, $n\ge1$ and $x\in\R$. Then
$N(\mu,n,x)\le\tfrac{293}{625}$.
\end{theorem}

\begin{proof}
Write $\ell=\beta/\sqrt n$ and split into four exhaustive cases, which is
exactly the case analysis of \file{BerryEsseen/UniversalUpperBound.lean}.
\begin{enumerate}[leftmargin=2.2em]
\item If $\ell\le\tfrac1{20}$: Theorem~\ref{thm:small} gives
  $N\le\tfrac{183}{400}\le\tfrac{293}{625}$.
\item Otherwise $\ell\ge\tfrac1{20}$. If $\beta\ge2$:
  Theorem~\ref{thm:scalar}(a).
\item Otherwise $\beta<2$. If $n\ge20$: Theorem~\ref{thm:scalar}(b).
\item Otherwise $1\le n<20$ and $\beta<2$: Theorem~\ref{thm:finite}.
\end{enumerate}
Each branch of cases 2--4 internally falls back on Corollary~\ref{cor:large}
where $\ell$ exceeds the last certified endpoint $\tfrac65$.
\end{proof}

\begin{proof}[Proof of Theorem~\ref{thm:main-upper}]
By Lemma~\ref{lem:reduction} and Theorem~\ref{thm:pointwise},
$\Cii\le\tfrac{293}{625}=0.4688$.
\end{proof}

The Lean text of the case split is short enough to quote in full:
\begin{verbatim}
theorem normalizedError_universal_upper_bound {mu : ProbabilityMeasure R}
    (hmu : Admissible mu) {n : Nat} (hn : 0 < n) (x : R) :
    normalizedError mu n x <= 0.4688 := by
  by_cases hs : thirdAbsMoment mu / sqrt n <= 1 / 20
  . exact normalizedError_le_target_of_small_fraction hmu hn hs x
  have hell : 1 / 20 <= thirdAbsMoment mu / sqrt n := (lt_of_not_ge hs).le
  by_cases hb : 2 <= thirdAbsMoment mu
  . exact ScalarFullHAll.large_moment_normalizedError hmu hn hb hell x
  by_cases hn20 : 20 <= n
  . exact ScalarFullHAll.large_sample_normalizedError hmu hn20 hell x
  exact FiniteParameterBlocks.finite_normalizedError hmu hn (by omega)
    (lt_of_not_ge hb) x

theorem berryEsseenConstant_upper_bound : berryEsseenConstant <= 0.4688 := by
  apply (constant_le_iff _).mpr
  intro mu hmu n hn x
  exact normalizedError_universal_upper_bound hmu hn x
\end{verbatim}
(Unicode has been transliterated for legibility; the actual file uses
the Unicode symbols for $\R$, $\N$, $\le$ and $\mu$.)

\begin{corollary}
$\Cii<\infty$, and in particular for every centred variance-one law with finite
third absolute moment and every $n\ge1$,
\[
  \sup_{x\in\R}\bigl|F_n(x)-\Phi(x)\bigr| \ \le\ 0.4688\,\frac{\E|X|^3}{\sqrt n}.
\]
Combined with Theorem~\ref{thm:main-lower}, $0.40\le\Cii\le0.4688$.
\end{corollary}

\begin{leanref}
\lean{BerryEsseen.normalizedError\_universal\_upper\_bound} and
\lean{BerryEsseen.berryEsseenConstant\_upper\_bound} in
\file{BerryEsseen/UniversalUpperBound.lean}; \file{Solution.lean} re-exports the
two theorems under the comparator's names and also proves their conjunction
\lean{berry\_esseen\_constant\_bounds}.
\end{leanref}

%=====================================================================
\section{Verification status and trust boundary}
\label{sec:verification}
%=====================================================================

\subsection{Axioms and tactics}

The two theorems are proved with no \lean{sorry} and no custom axiom, and their
axiom closure is exactly
\[
  \{\lean{propext},\ \lean{Classical.choice},\ \lean{Quot.sound}\},
\]
the three standard mathlib axioms. In particular \lean{native\_decide} does not
occur anywhere in the $0.4688$ development: a search of
\file{BerryEsseen/} outside \file{BerryEsseen/Upper0423/}, together with
\file{Challenge.lean} and \file{Solution.lean}, returns no occurrence of
\lean{native\_decide} and no occurrence of \lean{sorry}. All $961{,}710$ finite
decisions are \lean{decide +kernel}, i.e.\ kernel evaluation of exact integer and
rational arithmetic (Table~\ref{tab:kernel}).

This matters for the trust boundary: a \lean{native\_decide} check delegates
evaluation to compiled code and adds an auxiliary axiom per check, whereas
\lean{decide +kernel} is checked by the same kernel that checks the rest of the
proof. The sharper $0.423$ development in the same repository makes the opposite
choice; see Section~\ref{sec:0423-note}.

\subsection{The challenge file and the comparator}

\file{Challenge.lean} contains all definitions of Section~\ref{sec:statement}
written out and states the two theorems with intentional proof holes. It
imports only mathlib, and only three files of it:
\file{Mathlib.MeasureTheory.Constructions.Pi},
\file{Mathlib.MeasureTheory.Measure.ProbabilityMeasure} and
\file{Mathlib.Probability.Distributions.Gaussian.Real}. \file{Solution.lean} proves the same two
statements from the supporting library.

\file{comparator/config.json} is
\begin{verbatim}
{ "challenge_module": "Challenge",
  "solution_module": "Solution",
  "theorem_names": [ "berry_esseen_constant_lower_bound",
                     "berry_esseen_constant_upper_bound" ],
  "permitted_axioms": [ "propext", "Quot.sound", "Classical.choice" ] }
\end{verbatim}
and the pinned tools are those of Table~\ref{tab:pins}. The runner script
\file{scripts/run\_comparator\_linux.sh} checks out these revisions, builds the
tools and invokes the Landrun sandbox under the upstream-recommended
\texttt{systemd-run} restriction; it does not use the macOS development
adapter. The workflow \file{.github/workflows/comparator.yml} can be dispatched
without inputs; it builds the certificates in bounded batches
(\texttt{--batch-size 2}), asserts \texttt{git diff --exit-code} so that no
source file was modified during the build, provisions swap for the replay, and
runs the comparator, uploading its log as an artifact. An optional
\texttt{proof-cache-tag} may supply precompiled supporting outputs; the workflow
accepts a cache only if its supporting source contents, project configuration
and dependency identities match the checkout, it verifies every archive part's
size and SHA-256 before extraction, and it excludes
\file{BerryEsseen/Defs}, \file{Challenge} and \file{Solution} so that those are
always rebuilt from the current source.

\begin{table}[ht]
\centering
\begin{tabular}{@{}ll@{}}
\toprule
Tool & Pinned git revision \\
\midrule
\file{leanprover/comparator} & \texttt{2312244ac716564a61cc0bf4e107d9abf1757a61} \\
\file{leanprover/lean4export} & \texttt{cacf989bd75f608700820f6afc595f32e7a99a4d} \\
\file{Zouuup/landrun} & \texttt{811cfff51ceaf3d9843708aa6d22e9b84ccac8b4} \\
\bottomrule
\end{tabular}
\caption{The pinned verification tools \cite{Comparator2026}.}
\label{tab:pins}
\end{table}

\subsection{Status, stated exactly}

\begin{itemize}[leftmargin=2.2em]
\item \textbf{Proved in Lean.} Theorems~\ref{thm:main-lower} and
  \ref{thm:main-upper}, with the axiom closure and tactic discipline above.
\item \textbf{Comparator: pending.} The final pinned, sandboxed Linux
  comparator run --- statement comparison against \file{Challenge.lean}, axiom
  whitelist check, and kernel replay of the dependency closure --- has
  \emph{not} been reported as passing. The repository states this in
  \file{README.md} (``the separate pinned, sandboxed Linux comparator check is
  pending''), in \file{comparator/README.md} (``The final comparator run is
  still pending'') and in every proof note. Nothing in this document should be
  read as a claim of comparator completion. A passing ordinary
  \lean{lake build} is reported separately from the comparator's statement
  comparison and kernel replay; any local macOS run using the upstream
  development adapter is to be identified separately and is not a sandboxed
  verification.
\item \textbf{Local build.} The full library was rebuilt from scratch on the
  machine used to prepare this document (a ten-core Apple Silicon laptop):
  \texttt{python3 scripts/build\_certificates.py --batch-size 4} followed by
  \texttt{lake build Solution Challenge} completed successfully, reporting
  $5{,}566$ and $6{,}117$ finished jobs, in $183$ minutes of wall-clock time
  ($583$ minutes of CPU time). On the resulting build,
  \texttt{\#print axioms} reports exactly
  \lean{[propext, Classical.choice, Quot.sound]} for both
  \lean{berry\_esseen\_constant\_lower\_bound} and
  \lean{berry\_esseen\_constant\_upper\_bound}, and their statements are the two
  displayed in Section~\ref{sec:statement}. This is an ordinary Lean build and
  axiom query; it is not the sandboxed comparator run, which remains pending.
\end{itemize}

\subsection{Sanity properties of the formalization}

Three features are worth checking against the informal statement of the
Berry--Esseen problem, because they are the usual places where a formalization
can be weaker than it appears.
\begin{enumerate}[leftmargin=2.2em]
\item \emph{The class of laws is not restricted.} The supremum in
  Definition~\ref{def:constant} ranges over all \lean{ProbabilityMeasure}s on
  $\R$ satisfying Definition~\ref{def:admissible}; there is no finiteness,
  boundedness, or discreteness hypothesis, and no placeholder definition.
\item \emph{Independence is genuine.} $\mu_n$ is the pushforward of a finite
  product measure under the normalized sum, and
  $\hat\mu_n(t)=f(t/\sqrt n)^n$ is proved, not assumed.
\item \emph{Finiteness is a theorem.} $\Cii$ is defined in $[0,\infty]$; the
  upper bound is what makes it finite.
\end{enumerate}

%=====================================================================
\appendix
\section{Index of Lean files by role}
%=====================================================================

{\footnotesize
\begin{longtable}{@{}p{5.6cm}p{9.2cm}@{}}
\toprule
File (under \file{BerryEsseen/} unless noted) & Role \\
\midrule
\endfirsthead
\toprule
File & Role \\
\midrule
\endhead
\file{../Challenge.lean} & self-contained statement with two proof holes; imports only mathlib \\
\file{../Solution.lean} & the two comparator targets and their conjunction \\
\file{Defs.lean} & Definitions~\ref{def:admissible}--\ref{def:constant} \\
\file{Constant.lean} & Lemma~\ref{lem:reduction} \\
\addlinespace
\multicolumn{2}{@{}l}{\emph{Lower bound}}\\
\file{BernoulliWitness.lean} & the witness law, its moments, the $2^6$-term exact probability \\
\file{GaussianBounds.lean} & normal density and CDF estimates, Lemma~\ref{lem:normal-third} \\
\file{LowerBound.lean} & Theorem~\ref{thm:main-lower} \\
\addlinespace
\multicolumn{2}{@{}l}{\emph{Characteristic-function estimates}}\\
\file{MomentExcess.lean} & Lemma~\ref{lem:beta-ge-one} \\
\file{IntegralQuadratic.lean}, \file{WeightedMoments.lean} & Lemma~\ref{lem:weighted-cs} \\
\file{TrigonometricAnchors.lean} & Lemma~\ref{lem:sine-identity} \\
\file{ImaginaryPart.lean} & Proposition~\ref{prop:imag} \\
\file{RealPart.lean} & Proposition~\ref{prop:real} \\
\file{CharacteristicBounds.lean} & coordinate combination, $\beta=1\Rightarrow f=\cos$ \\
\file{GaussianCosine.lean}, \file{GaussianPolynomial.lean} & Lemma~\ref{lem:gauss-cos} and its polynomial refinement \\
\file{ClassicalError.lean} & the classical third-order error bound \\
\file{SumBounds.lean} & Definition~\ref{def:bounds} \\
\addlinespace
\multicolumn{2}{@{}l}{\emph{Global modulus}}\\
\file{Symmetrization.lean} & Lemma~\ref{lem:symm} \\
\file{CosineMinorant.lean}, \file{CosineProfile.lean}, \file{ConvexSplice.lean}, \file{FullMinorant.lean} & Proposition~\ref{prop:minorant} \\
\file{FullModulus.lean} & Theorem~\ref{thm:modulus} \\
\addlinespace
\multicolumn{2}{@{}l}{\emph{Prawitz smoothing}}\\
\file{KernelRecurrence.lean} & the two-monotone-sequence argument \\
\file{PrawitzKernel.lean} & integrability and the shift recurrence \eqref{eq:shift} \\
\file{OscillatoryIntegral.lean} & interval form of Riemann--Lebesgue \\
\file{PrawitzMajorant.lean} & Theorem~\ref{thm:majorant} \\
\file{PrawitzProbability.lean} & Proposition~\ref{prop:prawitz-prob} \\
\file{PrawitzFubini.lean}, \file{CharacteristicProjection.lean}, \file{PrawitzFourier.lean} & Proposition~\ref{prop:prawitz-fourier} \\
\file{GaussianInversion.lean}, \file{FirstMomentFourier.lean} & Lemma~\ref{lem:normal-fourier} \\
\file{SmoothingIntegrability.lean}, \file{SmoothingSplit.lean}, \file{SmoothingAlgebra.lean}, \file{PrawitzSmoothing.lean} & Theorem~\ref{thm:smoothing} \\
\file{SmoothingMajorants.lean} & Corollary~\ref{cor:majorant-form} \\
\file{PrawitzIntegralForm.lean} & integral reformulation of the kernel \\
\addlinespace
\multicolumn{2}{@{}l}{\emph{Compression}}\\
\file{GeometricDamping.lean} & $A$, $G_q$, Lemma~\ref{lem:geom-monotone}, \eqref{eq:G-poly} \\
\file{DampedPower.lean}, \file{PowerComparison.lean} & difference of powers, replacement of $1/n$ by $v$ \\
\file{ScalarCompression.lean} & Proposition~\ref{prop:compression} for the actual i.i.d.\ law \\
\file{VectorRemainder.lean}, \file{VectorCompression.lean} & moment-band coefficients and the vector prefactor \\
\file{ParameterCaps.lean} & Lemma~\ref{lem:caps} \\
\file{CompressedSmoothing.lean}, \file{MixedSmoothing.lean} & the envelope functional \eqref{eq:cell-functional} and its integrability \\
\file{ParameterCellBounds.lean} & Lemma~\ref{lem:cdf-to-normalized} and the two cell interfaces \\
\addlinespace
\multicolumn{2}{@{}l}{\emph{Region I and II}}\\
\file{SmallFractionCoefficients.lean}, \file{SmallFractionPower.lean}, \file{SmallFractionKernels.lean}, \file{SmallFractionGaussian.lean}, \file{SmallFractionLow.lean}, \file{SmallFractionNormal.lean}, \file{SmallFractionDecay.lean}, \file{SmallFractionHigh.lean}, \file{SmallFractionRemainders.lean}, \file{SmallFractionBudget.lean}, \file{SmallFraction.lean} & Theorem~\ref{thm:small} \\
\file{Cantelli.lean}, \file{LargeFraction.lean}, \file{IndependentSums.lean} & Theorem~\ref{thm:cantelli}, Corollary~\ref{cor:large} \\
\addlinespace
\multicolumn{2}{@{}l}{\emph{Region III}}\\
\file{DecayPanels.lean}, \file{DecayWitnesses.lean} & Lemma~\ref{lem:decay-profile} \\
\file{SmoothingPanels.lean}, \file{LowSmoothingPanels.lean}, \file{VectorSmoothingPanels.lean}, \file{PanelSums.lean}, \file{GaussianTailBound.lean} & Propositions~\ref{prop:high-normal-panel}--\ref{prop:tail} \\
\file{Numerics/ScalarPanelCore.lean}, \file{Numerics/FixedCellCore.lean}, \file{Numerics/GeneralCellCore.lean}, \file{Numerics/UpperPanelCore.lean} & executable record formats and validators \\
\file{Numerics/Cell*Facts.lean}, \file{Numerics/CellApplication.lean}, \file{Numerics/CellPanelSums.lean}, \file{Numerics/CellCertificate.lean}, \file{Numerics/CertifiedCell*.lean}, \file{Numerics/HighGrid.lean}, \file{Numerics/IndexedPanelLists.lean} & record-to-integral soundness (Table~\ref{tab:soundness}) and Theorem~\ref{thm:cell-sound} \\
\file{Certificates/ScalarFullH/} & the $467$ cells, the tables, the coverings \\
\addlinespace
\multicolumn{2}{@{}l}{\emph{Region IV}}\\
\file{FiniteSampleBounds.lean}, \file{FiniteModulus.lean}, \file{FiniteKernels.lean}, \file{FiniteGeometricAverage.lean}, \file{FiniteSmoothing.lean} & Theorem~\ref{thm:finite-smoothing} \\
\file{FiniteRadicandRange.lean}, \file{FiniteDarboux.lean} & Lemma~\ref{lem:radicand-range}, Proposition~\ref{prop:darboux} \\
\file{FiniteLowPanels.lean}, \file{FiniteHighPanels.lean}, \file{FiniteNormalPanels.lean} & generic real-integral comparisons per panel kind \\
\file{FiniteLowCover.lean}, \file{FiniteHighCover.lean}, \file{FiniteNormalCover.lean} & adjacency and completeness of each panel list \\
\file{FiniteCellCertificate.lean}, \file{FiniteParameterCover.lean}, \file{FiniteMomentCap.lean} & Definition~\ref{def:finite-cell}, Theorem~\ref{thm:finite} \\
\file{Certificates/Finite*} & the checked collection of Table~\ref{tab:counts} \\
\addlinespace
\multicolumn{2}{@{}l}{\emph{Numerics infrastructure}}\\
\file{TaylorPanel.lean}, \file{Numerics/TaylorPanelSoundness.lean} & Theorem~\ref{thm:taylor-panel} \\
\file{Numerics/JetCore.lean}, \file{Numerics/IntervalSoundness.lean} & Definition~\ref{def:interval} and outward rounding \\
\file{Numerics/*Jet.lean}, \file{Numerics/EnclosesOn.lean}, \file{Numerics/JetSoundness.lean}, \file{Numerics/JetInputs.lean}, \file{Numerics/JetPanelRange.lean} & jet recurrences and their derivative semantics \\
\file{Numerics/ScalarEnclosure.lean} & exact-value scalar enclosures \\
\file{Numerics/ExpSeed*.lean}, \file{Numerics/TrigSeed*.lean}, \file{Numerics/Trig*Soundness.lean}, \file{TrigonometricRange.lean} & Sections~\ref{sec:exp-seeds}, \ref{sec:trig-seeds} \\
\file{Numerics/RadicandRefinement.lean} & radicand refinements \\
\file{Numerics/ScalarTable*.lean}, \file{Numerics/CotangentTable.lean}, \file{CotangentPolynomials.lean}, \file{CotangentCorrection.lean}, \file{ScalarRemainders.lean}, \file{ExponentialRemainder.lean} & table semantics and rational remainders \\
\file{Numerics/DarbouxCertificate.lean}, \file{Numerics/FiniteTailCertificate.lean}, \file{Numerics/GeometricWitness.lean}, \file{Numerics/DecayCertificate.lean} & the remaining per-record checkers \\
\bottomrule
\end{longtable}
}

%=====================================================================
\section{The four smoothing terms at a glance}
%=====================================================================

For orientation, Table~\ref{tab:terms} lists which estimate controls which of
the four terms of Theorem~\ref{thm:smoothing} in each region.

\begin{table}[ht]
\centering
\footnotesize
\begin{tabular}{@{}llll@{}}
\toprule
Term & Region I & Region III & Region IV \\
\midrule
low-frequency error
  & \eqref{eq:small-error} $+$ \eqref{eq:small-power}
  & $\mathcal E$ via \eqref{eq:compress}
  & $\mathcal L$ via \eqref{eq:finite-power} \\
high-frequency magnitude
  & Theorem~\ref{thm:modulus} $+$ Lemma~\ref{lem:small-decay}
  & $e^{-\Psi}$ via Theorem~\ref{thm:modulus}
  & $M_B^N$, line/cosine/anchor branches \\
normal correction
  & \eqref{eq:q-coarse}
  & Proposition~\ref{prop:high-normal-panel}(b)
  & $K_C$ via \eqref{eq:q-sandwich} \\
Gaussian tail
  & elementary, $T\ge60$
  & Proposition~\ref{prop:tail}
  & Section~\ref{sec:finite-tails} \\
\bottomrule
\end{tabular}
\caption{Which estimate controls which term. Region II (Cantelli) uses no
smoothing at all.}
\label{tab:terms}
\end{table}

\bigskip
\noindent\textbf{Acknowledgement of scope.} This document is a mathematical
exposition of a formal development. It reports the development's own recorded
status and provenance (Sections~\ref{sec:intro-status},
\ref{sec:provenance}, \ref{sec:verification}) and makes no claim of priority,
and no claim about the correctness of any other development, in particular not
of \cite{XiaoLi2026}.

\bibliographystyle{plain}
\bibliography{references}

\end{document}
